\documentclass[11pt,a4paper]{ltjsarticle}
% Common preamble; compile the four documents with LuaLaTeX.
\usepackage{fix-cm}
\usepackage[margin=25mm]{geometry}
\usepackage{luatexja-fontspec}
\setmainjfont{HaranoAjiMincho-Regular.otf}[BoldFont=HaranoAjiMincho-Bold.otf]
\setsansjfont{HaranoAjiGothic-Medium.otf}[BoldFont=HaranoAjiGothic-Bold.otf]
\usepackage{amsmath,amssymb,amsthm,mathtools,booktabs,enumitem}
\usepackage[unicode,hidelinks]{hyperref}
\setlength{\emergencystretch}{2em}
\setlist{nosep}
\newtheorem{theorem}{定理}[section]
\newtheorem{proposition}[theorem]{命題}
\newtheorem{lemma}[theorem]{補題}
\newtheorem{corollary}[theorem]{系}
\theoremstyle{definition}
\newtheorem{definition}[theorem]{定義}
\newtheorem{remark}[theorem]{注意}
\newtheorem{example}[theorem]{例}
\newtheorem{problem}[theorem]{問題}
\newcommand{\C}{\mathbb C}
\newcommand{\Z}{\mathbb Z}
\newcommand{\Q}{\mathbb Q}
\newcommand{\Cat}{\mathcal A}
\newcommand{\DD}{\mathcal D}
\newcommand{\OO}{\mathcal O}
\newcommand{\TT}{\mathcal T}
\newcommand{\PP}{\mathcal P}
\newcommand{\SSer}{\mathbb S}
\newcommand{\Db}{D^{\mathrm b}}
\newcommand{\Hom}{\operatorname{Hom}}
\newcommand{\End}{\operatorname{End}}
\newcommand{\Ext}{\operatorname{Ext}}
\newcommand{\RHom}{\operatorname{RHom}}
\newcommand{\qgr}{\operatorname{qgr}}
\newcommand{\Gr}{\operatorname{Gr}}
\newcommand{\QGr}{\operatorname{QGr}}
\newcommand{\Tor}{\operatorname{Tor}}
\newcommand{\coh}{\operatorname{coh}}
\newcommand{\thick}{\operatorname{thick}}
\newcommand{\gldim}{\operatorname{gldim}}
\newcommand{\GKdim}{\operatorname{GKdim}}
\newcommand{\rk}{\operatorname{rk}}
\newcommand{\id}{\operatorname{id}}
\newcommand{\Hilb}{\operatorname{Hilb}}
\newcommand{\versiondate}{2026年11月17日・v12}

\hypersetup{pdftitle={反Serre作用の三乗根による非可換射影平面の認識}}
\title{反Serre作用の三乗根による\\非可換射影平面の認識}
\author{}
\date{\versiondate\\研究原稿：一般認識問題は未解決}
\begin{document}
\maketitle
\begin{abstract}
反標準ampleな非可換曲面の内在的六公理に，heart上の反Serre作用のexactな三乗根を加えた場合を考える．
次数9の反標準Hilbert関数の下で，固定した三乗根の座標環$A$の三次元二次AS-regular性を，
一つのExt群と一つの生成障害の零性に還元する．
既知の標準対象の内在的特徴付けを用い，この判定が標準的な非可換射影平面としての
非標識のabel圏の認識とも同値であることを示す．標準クラスではexactな反Serre指数は3である．
可換な滑らかな射影多様体への特殊化は，三乗根の仮定なしに射影平面を認識する．
六公理だけからrank-one torsion-free対象の$\chi(I,I)=1$と例外性の同値を示し，
漸近Serre消滅と全数値rank $3$の下では内在的な線束列とexact指数$3$も復元する．
三次の非零normal元$g$があり，$B=A/(g)$がdomainなら，$A,B$は両側noetherianで，
飽和後のsection Hilbert関数は自動的に$\dim\widehat B_n=3n$となる．
従ってEuler列は$P(n)-\epsilon\mathbf1_{3\nmid n}$，Euler orbit rankは3または5である．
証明は$\chi_1$のVeroneseへの継承，高Veroneseの曲線分類，構造対象のcohomology比較を用いる．
高Veroneseの曲線が滑らかなら，剰余層の積から元のheartと次数3のshiftを復元する．
さらに全数値格子がrank $5$以下でunimodularなら，周期成分の判別式から$\epsilon=0$を得る．
従って滑らかな高Veronese境界と長さ$5$以下のfull exceptional collectionは認識の十分条件となる．
一般には有限section欠損の消滅を主張しない．

中心となる追加条件付き認識は，$B^{(3)}$が可換であることと，$A$が三乗が恒等となる線形自己同型による
三変数多項式環のZhang twistであることの同値である．
境界の一次生成・飽和性・幾何学的表示・$g$の中心性を仮定しない．
証明では，normal切断の高次Ext周期性，局所方程式指標，conductor次数式を組み合わせて
整数次数の忠実な偏極を持つ境界stackのstabilizerを除く．
反標準Veronese上でnormal元を中心化し，負次数双微分の消滅から通常の次数9 del Pezzo曲面へ帰着する．
さらに元のexactな三乗根を識別する．$B$自体が可換なら$A$は通常の多項式環となる．

一般の通常曲線表示を持つ境界についても，非可換THCRを保った中心化を与え，
認識を$\epsilon=0$という一つの障害へ還元する．
また，正次数normal元による全対象Serre消滅，次数1・2の整なnormal切断による認識，
一次生成三次境界の被覆のAS-regular型と分裂Ore型への二分，残余回転，mutation純粋性，
低次数増大，重み付き被覆，双有理的に幾何的な場合の認識を証明する．
元の公理から一般の認識障害が消えることは未証明であり，
既知結果と本稿で証明する組合せを区別する．文献上の新規性の確定は残されている．
\end{abstract}
\begingroup\fontsize{11}{14}\selectfont
\tableofcontents
\endgroup

\section{対象と認識問題}\label{sec:axioms}
本稿で認識する標準クラスは，三次元二次AS-regular代数$S$に対する
$\qgr S$と$\C$線形exact同値なabel圏である．通常の次数付き代数の表示は一意ではない．
三次元AS-regularという条件だけでは三次関係式を持つ別の型も含むため，二次型を明示する．
以下の六公理は，この既知のクラスを内在的に特徴付けるための候補条件であり，
それを満たすこと自体を標準クラスへの所属という定義に置き換えない．

可換射影幾何における対応する認識は，滑らかなdel Pezzo曲面の次数9の場合，
および一般次元の最大Fano指数の特徴付けである．後者は，ample Cartier divisor $H$と
$-K_X\equiv(n+1)H$から$X\simeq\mathbb P^n$を導く
\cite[Corollary 2(1)]{KachiKollar}．下の命題\ref{prop:v9-commutative-specialization}は，
可換な滑らかな場合には本稿の六公理だけで射影平面が決まることを確認する．
非可換の場合には通常の曲面の分類やKodaira消滅を直接使えないため，
その役割を担う負次数のExt消滅と生成を分離して扱う．

標準クラスを認識対象に選ぶ理由は，単に座標環のHilbert級数が一致するためではない．
Van den Bergh\cite[Definition 8.1, Theorem 8.1.1]{VdB}は，$\coh\mathbb P^2$の
形式変形を，基底上flatな三次元二次regular $\Z$-algebra族の$qgr$として記述する．
またAbdelgadir--Okawa--Ueda\cite[Theorem 7.3, Appendix A]{AOUCompact}は，
標準クラス内ではabel圏から直線対象の順序集合と$\Z$-algebraを復元する．
本稿ではこの復元を候補圏に無条件に適用せず，標準クラスに属すると仮定した方向で用いる．
これにより，固定した三乗根の認識が非標識の圏の認識と同値であることを得る
（命題\ref{prop:v9-marking-independence}）．

係数体を$\C$とする．$\Db(\Cat)$のSerre関手を$\SSer$と書き，
\[
 \theta=\SSer[-2],\qquad \tau=\theta^{-1}=\SSer^{-1}[2]
\]
と置く．以下では導来圏上の自己同値とそのheartへの制限を，区別が必要な場合を除いて同じ記号で表す．

\begin{definition}[基礎公理]\label{def:axioms}
次数9の候補$(\Cat,\OO)$に次の条件を課す．
\begin{enumerate}[label=(\arabic*)]
\item $\Cat$は本質的に小さい$\C$線形noetherian Ext-finite abel圏である．
\item $\Db(\Cat)$はSerre関手$\SSer$を持ち，$\theta$は$\Cat$を保つexact autoequivalenceである．
\item $\Cat_0$を有限長対象のSerre部分圏，$\Cat_1$を$\Cat/\Cat_0$で有限長になる対象の逆像とする．
商圏$\Cat/\Cat_1$における$\OO$の像は非零単純対象であり，
$T\in\Cat_1$に対して$\Hom(T,\OO)=0$である．
\item $\OO$はexceptionalである：$\End(\OO)=\C$かつ$\Ext^i(\OO,\OO)=0$ for $i>0$．
\item $(\OO,\tau)$はArtin--Zhangのample pairである．
\item 全ての$n\geq0$に対して
\begin{equation}\label{eq:anticanonical}
 \dim\Hom(\OO,\tau^n\OO)=1+\frac92n(n+1)
 =\binom{3n+2}{2}
\end{equation}
である．
\end{enumerate}
\end{definition}

ここで$(\OO,\tau)$のample性とは，任意の$F\in\Cat$に対して有限直和からの全射
$\bigoplus_i\tau^{-m_i}\OO\twoheadrightarrow F$，$m_i\geq0$が存在し，
さらに任意の全射$F\twoheadrightarrow G$に対して
$\Hom(\tau^{-n}\OO,F)\to\Hom(\tau^{-n}\OO,G)$が$n\gg0$で全射になることをいう．

\begin{proposition}[可換な特殊化の認識]\label{prop:v9-commutative-specialization}
$X$を滑らかな連結射影複素多様体，$F\in\coh X$とし，
$(\coh X,F)$が定義\ref{def:axioms}の六公理を満たすとする．
このとき$X\simeq\mathbb P^2$であり，$F$は線束である．
従って線束によるtensor同値まで
$(\coh X,F)\simeq(\coh\mathbb P^2,\OO_{\mathbb P^2})$である．
この主張には反Serre関手の三乗根の存在を仮定しない．
\end{proposition}
\begin{proof}
$d=\dim X$とするとSerre関手は$(-)\otimes\omega_X[d]$である．
その$[-2]$シフトが標準heartを保つので$d=2$である．
滑らかな曲面上の$\coh X$のGabriel filtrationでは，$\Cat_0$は零次元supportの層，
$\Cat_1$は一次元以下のsupportの層である．
従って公理(3)より$F$はrank one torsion-freeで，
$F\simeq\mathcal L\otimes I_Z$と書ける．ここで$\mathcal L=F^{**}$は線束，$Z$は零次元部分schemeである．
滑らかな曲面では$\mathcal Hom(I_Z,\OO_X)\simeq\OO_X$である．
これは$\mathcal Ext^1(\OO_Z,\OO_X)=0$を
$0\to I_Z\to\OO_X\to\OO_Z\to0$に適用すれば分かる．
従って任意の$m$に対し，任意の局所射
$F\otimes\omega_X^m\to\OO_X$の像は$I_Z$に含まれる．
AZ生成条件を対象$\OO_X$に適用すると，これらの有限直和から$\OO_X$への全射が存在するので，
$I_Z=\OO_X$，すなわち$Z=\varnothing$である．よって$F=\mathcal L$は線束である．

線束$\mathcal L^{-1}$によるtensorで構造対象を$\OO_X$に取り替え，
$H=\omega_X^{-1}$と置く．$H$のample性をAZ生成条件から直接確認する．
任意のvery ample線束$A$を一つ固定する．有限個の射
$H^{-m_i}\to A^{-1}$，$m_i\geq0$が全射を与えるので，
section $s_i\in H^0(X,H^{m_i}\otimes A^{-1})$は，各点で少なくとも一つ非零である．
非零の$s_i$について$m_i>0$であり，その零divisorを$D_i$とすると
$m_iH\sim A+D_i$である．任意の既約曲線$C$上の点を一つ選び，そこで非零となる$s_i$を取れば，
$C$は$D_i$の成分でないので
\[
 m_i(H\cdot C)=A\cdot C+D_i\cdot C>0.
\]
従って$H$は全ての曲線に正の次数を持つ．さらに一つの非零$s_i$について
\[
 m_i H^2=H\cdot A+H\cdot D_i>0
\]
であるから，Nakai--Moishezon判定により$H$はampleである．

公理(6)と$F$の線束性から
\[
 h^0(X,H^m)=1+\frac92m(m+1)\qquad(m\geq0).
\]
$m\gg0$でSerre消滅とRiemann--Rochを用いれば
\[
 h^0(X,H^m)=\chi(\OO_X)+\frac{K_X^2}{2}m(m+1).
\]
係数比較により$K_X^2=9$である．従って$X$は次数9の滑らかなdel Pezzo曲面であり，
その古典的分類\cite[\S2.1, \S3.1.1]{DolgachevMartin}から$X\simeq\mathbb P^2$を得る．
\end{proof}

本稿ではさらに，heart上の$\C$線形exact autoequivalence$s$と自然同型
\begin{equation}\label{eq:root}
 s^3\simeq\tau
\end{equation}
を固定する．導来圏だけに存在する，heartを保たない三乗根は扱わない．
$\OO(i)=s^i\OO$，$\DD=\Db(\Cat)$と書く．条件(4)は$\OO$自身の例外性であり，
一般の$\Ext^{>0}(\OO,s^nF)$の消滅を含まない．

\paragraph{中心となる認識定理．}
$A=\bigoplus_{n\geq0}\Hom(\OO,s^n\OO)$と置き，$0\ne g\in A_3$をnormal元，$B=A/(g)$をdomainとする．
定理\ref{thm:v8-cubic-standard-sections}は，可換性やEuler rankの有限性を追加せず
\[
 \dim\widehat B_n=3n\quad(n>0),\qquad
 q_n=P(n)-\epsilon\mathbf1_{3\nmid n}\quad(n\in\Z)
\]
を与える．ここで$\widehat B_n=\Hom_{\qgr B}(\pi B,\pi B(n))$，
$P(n)=\binom{n+2}{2}$，$\epsilon=\dim\Ext^1(\OO,s^{-1}\OO)$である．
Euler orbit rankは$\epsilon=0$なら3，正なら5となる．
この数値的結論と，元の境界の一次生成・飽和性とは区別する．
高Veronese表示の曲線が滑らかなら，定理\ref{thm:v10-smooth-veronese-heart}は
$(\qgr B,D,(1))$を通常の次数3の曲線tripleとして復元する．
また周期成分のEuler行列の判別式は$3^5\epsilon^2$であり，
全数値格子がrank $5$以下でunimodularなら$\epsilon=0$となる．
これらを合わせた系\ref{cor:v10-five-exceptionals-recognition}は，滑らかな高Veronese境界と
長さ$5$以下のfull exceptional collectionから固定した$A$を認識する．
この例外列にstrong性やrank-one性は仮定しない．

定理\ref{thm:v8-commutative-veronese-boundary}は，さらに
\[
 B^{(3)}\text{ が可換}
 \ \Longleftrightarrow\ A^{(3)}\text{ が可換}
 \ \Longleftrightarrow\ A\simeq\C[x_0,x_1,x_2]^\sigma,
 \quad\sigma^3=1
\]
を証明する．$\sigma$は次数1上の線形自己同型から誘導される．
構造対象と固定した$s$を含む表示は
$(\coh\mathbb P^2,\OO,\alpha^*(-)\otimes\OO(1))$，$\alpha^3=1$である．
この場合$\epsilon=0$で，境界は飽和したintegral plane cubicのTHCRとなる．
境界の一次生成・飽和性・幾何学的表示・$g$の中心性は追加仮定ではない．
$B$自体が可換なら，定理\ref{thm:v7-commutative-cubic-recognition}により$A$は通常の多項式環となり，
固定した$s$も$\OO(1)$によるtensorに対応する．

\paragraph{既知結果との関係．}
Artin--Zhangの再構成定理\cite[Theorem 4.5]{AZ}は，ample pairからnoetherianな座標環を復元する．
Artin--Stafford\cite[Theorems 0.1--0.2, 0.4--0.5]{ASquadratic}のGK次元2のdomainの分類
（共著者による再掲\cite[Theorems 5.1--5.2, 5.5--5.6]{SVdBsurvey}）は，高Veroneseの曲線表示を与える外部入力である．
本稿では両側有限性と$\chi_1$の継承を証明し，その分類から全次数でのgenus-one THCRを得る．
Bondal--Van den Bergh\cite[\S4.1, (4.1), Lemmas 4.1.4--4.1.6]{BvdBGenerators}の
truncation cohomology公式を用いて，圏同値を仮定しない有限Veroneseのcohomology比較を証明する．
その後，正負両側のRiemann--Rochと境界双対性により元の全次数のsection列を決定する．
Zhang twistの一般理論\cite{ZhangTwist}と，既にskew polynomialであるambientに対する
skew cubicの曲線表示\cite[\S6.1]{BBvGCentral}は既知である．
本稿の第三Veronese可換認識は，そのambientの生成・twist表示を内在的な仮定から導く点にある．
一方，Mori--Ueyama\cite[Corollary 4.2]{MU}のAS-regular認識にはfull geometric helixが必要である．
定理\ref{thm:finite}は，このhelix仮定を一つのExt群と一つの生成条件に還元する．
純粋なhelixの幾何学性はBridgeland--Stern\cite[Theorem 5.3(ii)]{BS}による．
三次境界を扱う節\ref{sec:elliptic}の曲線上の関係式とoverlapの幾何は，
Artin--Tate--Van den Bergh\cite[Theorem 6.6, \S7]{ATV1990}に対応する既知の内容である．
本稿ではintegralな場合に評価核のstabilityから必要な性質を証明し，
同論文のstandard algebraの判定\cite[Theorem 1]{ATV1990}を用いる
（\cite[Theorem 2.2]{IM}にも再掲されている）．
さらに，任意のregular normalな三次被覆をAS-regular型と分裂Ore型に二分し，
後者を有限ホモロジー次元で除く．中心的な滑らかな楕円境界に対する
Rogalski--Sierra--Stafford\cite[Proposition 4.5, Remark 4.6]{RSS}の機構との関係も明記する．
残余回転のSerre関係はKuznetsov--Smirnov\cite[Theorem 2.8]{KS}の構成に対応するが，
本稿では飽和性を追加せず直接証明し，AZ生成条件から一対象の回転軌道による生成を加える．
障害加群の有限性については，Yekutieli--Zhang\cite[Theorem 2.3]{YZ}の
Serre双対性の逆向きだけでは下位コホモロジーの有限性は得られない．
本稿はこの有限性を既知定理の仮定に置き換えず，左右作用とBass空間で不足箇所を特定する．
normal hypersurfaceの高次周期性はtwisted matrix factorizationの既知理論
\cite[Theorem 4.2]{CCKM}と関係するが，同定理の有限global dimensionを前提にせず，
本稿では$\qgr A$の有限homological dimensionからExt周期性を直接導く．
その後，構造対象を固定するnormal twistとAZ生成条件を使って局所的な単純対象を制約する．
nodal stackの局所形式は\cite[Introduction]{LPstack}に従い，
同文献のAuslander orderによる圏の解消と，本稿のnormal被覆の非存在は異なる主張として扱う．
局所平面性にはGulliksen\cite[Proposition 6]{Gulliksen1967}の有界Betti数の判定を用い，
構造対象の双対性から双対化層を同定する箇所ではLevchenko\cite[Proposition 2]{Levchenko}を用いる．
equivariantな双対化指標とconductorはNironi\cite{NironiDuality}，
局所種数1のGorenstein特異点の分類はSmyth\cite[Appendix A, Proposition A.3]{SmythGenusOne}に依存する．
これらは既知の外部入力であり，本稿ではnormal切断からの高次Ext周期性とAZ生成条件を組み合わせて，
方程式指標の忠実性とstabilizerの消滅を証明する．
滑らかな正種数曲線のsection環について，負次数multiderivationの消滅とfiltered deformationの可換性は
Rains\cite[Corollaries 2.5--2.6]{RainsFiltered}の既知結果である．本稿は必要な双微分の議論を
節点・尖点にも含めて直接証明し，normal automorphismのscalar性と反標準Veroneseの再構成を介して認識へ適用する．
最後の通常の次数9 del Pezzo曲面の識別は古典的分類
\cite[\S2.1, \S3.1.1]{DolgachevMartin}による．この組合せ全体の文献上の新規性は未確認である．

内在的な線束対象については，補題\ref{lem:v10-negative-line-hom}と
命題\ref{prop:v10-lines-exceptional-order}を六公理から直接証明し，AOUの認識を候補圏へ先取りしない．
条件\eqref{eq:v10-SV}の下で数値類による一意性を証明し，全数値rank $3$なら
定理\ref{thm:v10-numerical-line-lattice}により線束列を復元する．この段階でfull性は主張しない．
境界heartの復元では，高次数の分数可逆層の積が次数一致によって等号になることを用い，
余有限cocycleを一次cocycleへ延長する．Artin--Staffordのorder分類との新規性比較は未完了である．
判別式の計算は明示的Euler列からの初等的帰結であり，独立した新規性を主張しない．

本稿で閉じる追加仮定を区別する．
第一に，正次数normal元の存在だけでExt障害の有限次元性を得る
（命題\ref{prop:normal-vanishing}）．これは商の正則性やfull $\chi$を仮定しない．
第二に，次数1または2のnormal元の商がdomainなら認識まで得る
（定理\ref{thm:normal-domain-recognition}，\ref{thm:quadratic-normal-recognition}）．その証明では，切断商の次数supportに対する
Serre双対性の障害と増大度評価から，まず全次数のHilbert関数を決定する．
さらに切断商の識別と生成の議論から二次表示を導く．
この後でIyudu--Shkarin\cite[\S12.1]{IS}の，標準Hilbert級数を持つ三生成元二次代数の分類を用いる．
本稿の役割は，この既知分類の前提となる二次性と標準Hilbert級数を内在的条件から導くことにある．
次数2では，Artin--Stafford\cite[Theorem 0.2]{ASquadratic}の曲線分類により商を
$\mathbb P^1$の次数2のTHCRと識別する．さらに有限ホモロジー次元で分裂拡大を除き，
Cassidy--Shelton\cite[\S6, (6.1)--(6.2)]{CassidyShelton}の二つの環の変更の完全列と，
conicのKoszul双対の具体的なgraded primenessから二次性を導く．
Shelton--Tingey\cite{SheltonTingey}が示すKoszul性の下降と本稿の持上げは，仮定の方向を区別する．
Hu--Mori--Wu\cite[Definitions 1.1--1.2]{HuMoriWu}のcentral conic分類は，
ambient代数が既にquantum polynomial algebraであるクラスを対象とする．
本稿では内在的な圏の条件からambient代数の認識を先に行い，その後でこのような境界分類と比較できる形にする．
切断商が全次数で飽和する場合には，次数に関係なく$E$全体も消滅する
（命題\ref{prop:saturated-boundary}）．
第三に，次数3のnormal商がdomainかつ一次生成なら，
Artin--Staffordの全次数の埋込みと種数1曲線上の積の全射性から，
商を次数3のTHCRと低次数を含めて識別する．この段階では曲線の滑らかさを仮定しない．
上記の被覆の二分により，定理\ref{thm:cubic-generated-boundary-recognition}は
商の一次生成性とambient代数のAS-regular性が同値であることを示す．
曲線上の幾何学的補題は既知結果の再証明であり，この認識への組合せの新規性は未確認である．
一次生成を外した範囲では，$C=\widehat B/B$，$\epsilon=\dim E_{-1}$に対して
\[
 \dim C=2\epsilon,
 \qquad E_{-1}=0\ \Longleftrightarrow\ E=0\ \Longleftrightarrow\ B=\widehat B
\]
を証明する（命題\ref{prop:cubic-defect-mass}）．
この範囲では定理\ref{thm:v8-cubic-standard-sections}により$\widehat b_n=3n$は自動的に成立し，
Euler列とrankは上記の形に決まる．有限rankを仮定した独立の数値的証明も
命題\ref{prop:v7-finite-rank-cubic-euler}に与える．$\epsilon>0$の排除は未解決である．
可換な三次domain境界については，定理\ref{thm:v7-commutative-boundary-geometry}が
構造対象とshiftを保つ通常の種数1曲線表示を導き，
定理\ref{thm:v7-commutative-cubic-recognition}が有限欠損と生成障害の両方を除く．
第三Veroneseだけが可換な場合にも，定理\ref{thm:v8-commutative-veronese-boundary}が認識を与える．
通常の曲線表示が与えられた一般の非可換境界では，命題\ref{prop:v8-geometric-central-reduction}が
非可換THCRを保って反標準環を中心化し，系\ref{cor:v8-geometric-one-obstruction}が
飽和性・一次生成性・AS-regular性・rank条件・$h_1=3$・$h_2=6$を同値にする．
RSS\cite[Notation 3.2, Theorem 4.3]{RSS}のdegree-nine認識には，さらに特定のidealとdivisorに関する条件が必要である．
それらが成立すれば，認識された圏の$K_0$のrankから元の固定した三乗根も認識できる．
一方，系\ref{cor:v6-commutative-degree-two-recognition}は，標準Hilbert列と次数$2$までの生成を仮定する代わりに，
六公理やambientのSerre関手を必要としない被覆定理として独立に保持する．
系\ref{cor:v6-balanced-node-obstruction}および定理\ref{thm:v7-faithful-equation-character}は，
構造対象を保つ曲線stack表示を持つnormal境界一般にも適用される．
一般の非可換境界からこの表示を得たという主張ではない．
第四に，一次生成・両側noetherian・双有理的に幾何的という別の追加条件の下では，
Rogalski--Stafford\cite[Theorem 1.1, Corollary 1.2]{RSsurface}の曲面分類を適用できる．
同分類のnaive blowupの場合はExt-finitenessで排除し，残る曲面上で三乗根と次数9を使って
$(\mathbb P^2,\OO,\alpha^*(-)\otimes\OO(1))$を得る（定理\ref{thm:birational-recognition}）．
これらの追加条件が全公理から従うという主張ではない．

最小増大の候補$\dim A_1=4$については，normal affine semigroupを台とする型を
Ehrhart reciprocity\cite{StanleyReciprocity,BeckDevelin}で排除する．
他方，研究ノートには全反標準Hilbert列とcanonical moduleのシフトを満たす
noetherian Gorenstein domainを構成する．その例ではAZ生成条件と有限ホモロジー次元が失敗する．
従って，最小増大を一般に排除するには，これら圏論的条件の実際の使用がなお必要である．

三乗根の存在は，既知の三次元二次AS-regular型を排除しない．
実際，Itaba--Matsuno\cite[Theorems 1.1, 4.4]{IM}によれば，任意の三次元二次AS-regular代数は
Calabi--Yau AS-regular代数とgraded Morita同値である．
その$qgr$上の次数シフト$(1)$を移せば，三乗が$\SSer^{-1}[2]$となるexact自己同値を得る．
一般の表示における元の次数シフトをそのまま三乗根と見なす主張ではない．

\paragraph{結論の対象．}
軌道環の同型類は固定した$(\Cat,\OO,s)$に依存するが，命題\ref{prop:v9-marking-independence}により，
本稿のrank-one条件と例外性の下では，そのAS-regular性は非標識のabel圏が標準クラスに属することと同値である．
従って一つの許容される三乗根で認識障害が非零なら，同じabel圏の別のAS型表示も存在しない．
この主張は，heartを変更する導来同値や，rank-one条件を失う構造対象の変更を扱わない．

\section{公理から直接従う構造}\label{sec:elementary}
\begin{lemma}\label{lem:dimension}
$\operatorname{hdim}\Cat=2$である．また，全ての$\OO(i)$は条件(3)の意味でrank oneかつtorsion-freeである．
\end{lemma}
\begin{proof}
$X,Y\in\Cat$についてSerre双対性は
\[
 \Ext^i(X,Y)^*\simeq\Hom_{\DD}(Y,\theta X[2-i])
\]
を与える．$\theta X\in\Cat$なので$i>2$では零となる．
一方$\Ext^2(\OO,\theta\OO)^*\simeq\End(\theta\OO)=\C$である．
有限長性，およびSerre商での有限長性はexact自己同値で保たれるので，
$s$は$\Cat_0,\Cat_1$を保ち，後半が従う．
\end{proof}

\begin{lemma}\label{lem:mono}
軌道対象間の非零射$\OO(i)\to\OO(j)$は単射である．さらに
\begin{equation}\label{eq:negative}
 \Hom(\OO,\OO(-m))=0\quad(m>0)
\end{equation}
である．
\end{lemma}
\begin{proof}
非零射$f$が$\Cat/\Cat_1$で零なら$\operatorname{im}f\in\Cat_1$となり，標的のtorsion-free性に反する．
商圏で$f$は単純対象間の同型なので$\ker f\in\Cat_1$であり，始域のtorsion-free性から$\ker f=0$である．
もし$\OO\to\OO(-m)$が非零なら，その三つのtwistの合成で$\OO(3m)\hookrightarrow\OO$を得る．
従って$\Hom(\OO,\OO(3m))\hookrightarrow\End(\OO)=\C$となるが，
式\eqref{eq:anticanonical}により左辺の次元は1より大きい．
\end{proof}

\begin{proposition}\label{prop:reconstruction}
$(\OO,s)$もampleである．積を$a\in A_m,b\in A_n$に対して
$a*b=s^n(a)\circ b$と定めた環
\begin{equation}\label{eq:orbitring}
 A=A_s=\bigoplus_{n\geq0}\Hom(\OO,\OO(n))
\end{equation}
はconnected right-noetherian graded domainで，right $\chi_1$を満たし，
\[
 (\Cat,\OO,s)\simeq(\qgr A,\pi A,(1)).
\]
第三Veronese$A^{(3)}$は反標準軌道環と同型で，
\begin{equation}\label{eq:veronese}
 H_{A^{(3)}}(t)=\frac{1+7t+t^2}{(1-t)^3}.
\end{equation}
\end{proposition}
\begin{proof}
AZの生成条件は$\tau^{-m}\OO=s^{-3m}\OO$の直和だけで成立する．
全射$F\twoheadrightarrow G$について$n=3m+r$，$0\leq r<3$と書けば
\[
 \Hom(s^{-n}\OO,F)=\Hom(\tau^{-m}\OO,s^rF).
\]
三つの剰余類それぞれに元のAZの全射性条件を適用して，$(\OO,s)$のample性を得る．
Artin--Zhangの定理により再構成，right noetherian性，right $\chi_1$が従う．
元の$(\OO,\tau)$にも同じ定理を適用できるので，反標準環もright $\chi_1$を満たす．
非零斉次元は単射で表され，積も非零である．最高次数成分を取れば一般の非零元の積も非零なので$A$は整域である．
自然同型$s^3\simeq\tau$を反復して軌道環を同定し，\eqref{eq:anticanonical}を総和すると\eqref{eq:veronese}を得る．
\end{proof}

\begin{proposition}[三乗根次数での増大度]\label{prop:growth}
$h_n=\dim A_n$とすると，追加の消滅仮定なしに
\begin{equation}\label{eq:growth}
 h_n=\frac12n^2+O(n),\qquad \GKdim A=3
\end{equation}
が成立する．これは全次数でのHilbert多項式の存在や一次生成性を主張しない．
\end{proposition}
\begin{proof}
元のAZ生成条件を$\OO(r)$，$r=1,2$に適用する．
全射を与える有限直和の成分に非零射$\OO(-3m)\to\OO(r)$があるので，
$t_r\equiv r\pmod3$かつ$A_{t_r}\ne0$となる$t_r\geq0$が存在する．$t_0=0$とする．
$n\equiv r\pmod3$のとき，固定した非零斉次元による乗法と整域性から
\[
 h_{n-t_r}\leq h_n\leq h_{n+t_{-r}},
\]
ただし添字$-r$は$\Z/3\Z$で解釈し，$n$は十分大きいとする．
両端は3の倍数次数なので，それぞれ\eqref{eq:anticanonical}で与えられる．
引数の差が固定された二次式による挟み撃ちで\eqref{eq:growth}の第一式を得る．
$A$は有限生成であり，累積次元の増大は三次だから$\GKdim A=3$である．
\end{proof}


\subsection{内在的な線束対象：例外性，半順序，数値類による一意性}
\label{subsec:v10-intrinsic-lines}
以下では六公理を仮定し，$\tau=\SSer^{-1}[2]$とする．
負方向のHom消滅と自己Extの命題には三乗根の存在を仮定しない．
$\Cat_1$を六公理のSerre部分圏とし，
$I$がrank one torsion-freeであるとは，
$\Cat/\Cat_1$における像が非零単純で，$\Hom(\Cat_1,I)=0$であることをいう．
線束対象の集合を
\[
 \mathcal L(\Cat)
 =\{I\mid I\text{ はrank one torsion-free},\ \chi(I,I)=1\}/\simeq
\]
と置く．既知の非可換射影平面におけるAOUの分類を，以下の候補圏に適用することはしない．

\begin{lemma}[負の反標準方向のHomの有界性]
\label{lem:v10-negative-line-hom}
$I,J$をrank one torsion-free対象とする．任意の非零射$I\to J$は単射である．
また，ある$a\geq0$について非零射$\tau^{-a}\OO\to I$が存在し，
そのような$a$を一つ固定すると，任意の$m\geq0$に対して
\[
 \Hom(I,\tau^{-m}J)\ne0
 \quad\Longrightarrow\quad
 1+\frac92m(m+1)\leq\dim\Hom(\tau^{-a}\OO,J)
\]
が成立する．特に$\Hom(I,\tau^{-m}J)=0$ for $m\gg0$である．
\end{lemma}
\begin{proof}
非零射の像が$\Cat_1$に属せば，標的のtorsion-free性に反する．
従ってSerre商における射は単純対象間の非零射であり，同型である．
そのkernelは$\Cat_1$に属するので，始域のtorsion-free性により零となる．

AZ生成条件で$\bigoplus_i\tau^{-a_i}\OO\twoheadrightarrow I$を取り，
一つの非零成分を選べば$a$を得る．この射も前半により単射である．
$0\ne f:I\to\tau^{-m}J$が存在すれば，$\tau^m f$と$\tau^m(\tau^{-a}\OO\to I)$の合成は
単射$\tau^{m-a}\OO\hookrightarrow J$を与える．従って
\[
 \Hom(\tau^{-a}\OO,\tau^{m-a}\OO)
 \hookrightarrow\Hom(\tau^{-a}\OO,J).
\]
左辺の次元は六公理の反標準Hilbert関数によって
$1+\frac92m(m+1)$である．右辺はExt-finitenessにより有限次元であり，主張を得る．
\end{proof}

\begin{proposition}[rank one対象の自己Extと内在的半順序]
\label{prop:v10-lines-exceptional-order}
任意のrank one torsion-free対象$I$について
\[
 \End(I)=\C,\qquad
 \Hom(I,\tau^{-m}I)=0\quad(m>0),\qquad
 \Ext^2(I,I)=0
\]
である．従って
\[
 I\in\mathcal L(\Cat)
 \quad\Longleftrightarrow\quad
 I\text{ はrank one torsion-freeかつexceptional}
\]
である．さらに
\[
 [I]\preceq[J]\quad\Longleftrightarrow\quad\Hom(I,J)\ne0
\]
は$\mathcal L(\Cat)$上の半順序であり，exact自己同値はこの半順序を保つ．
\end{proposition}
\begin{proof}
前補題により，非零endomorphismは単射なので，二つの非零endomorphismの積は非零である．
$\End(I)$は有限次元$\C$-代数で零因子を持たないためdivision algebraであり，
$\C$が代数閉であることから$\End(I)=\C$である．

$I\to\tau^{-m}I$，$m>0$が非零なら，そのtwistを反復合成して
$I\hookrightarrow\tau^{-km}I$を全ての$k\geq1$について得る．
前補題の有界性に反するので，このような射は存在しない．
Serre双対性と$\SSer=\tau^{-1}[2]$から
\[
 \Ext^2(I,I)\simeq\Hom(I,\tau^{-1}I)^*=0.
\]
六公理から$\operatorname{hdim}\Cat=2$であるため，
$\chi(I,I)=1-\dim\Ext^1(I,I)$であり，例外性との同値が従う．

半順序の反射性は恒等射，推移性は単射の合成による．
$0\ne f:I\to J$と$0\ne g:J\to I$が存在すれば，$gf$は非零scalar倍の恒等射である．
よって$f$はsplit monoであり，$g$自身もmonoなので$f,g$は同型である．
これが反対称性を与える．exact自己同値はSerre filtration，rank，torsion-free性，Euler形式，
およびHomの非消滅を保つ．
\end{proof}

\begin{remark}[六公理の例外性の書き方]
公理(4)以外の五条件を仮定し，代わりに$\chi(\OO,\OO)=1$を課しても，
上の証明から公理(4)が従う．Serre関手のheart-exact性によるhomological dimension $2$，
rank-one torsion-free性，AZ1，反標準Hilbert関数，Ext-finitenessだけを使っているためである．
これは例外性の条件をEuler条件へ書き換えられるという主張であり，
そのEuler条件自体が残りの公理から従うという主張ではない．
負方向消滅と自己例外性の証明は，反標準切断次元が非有界であることを使い，
次数9の具体式を必要とするのは明示的な上界だけである．
また半順序の全順序性や，線束対象の集合が整数添字の一軌道となることは，この命題の結論ではない．
\end{remark}

次の条件を明示的に追加する：
\begin{equation}\label{eq:v10-SV}
 \text{各 }F\in\Cat\text{ に対し }
 \Ext^{>0}(\OO,\tau^nF)=0\quad(n\gg0).
 \tag{SV}
\end{equation}
exactな三乗根を固定した現行の設定では，命題\ref{prop:serrevanishing}により
$E=\bigoplus_n\Ext^1(\OO,s^n\OO)$の有限次元性と同値である．
特に命題\ref{prop:normal-vanishing}のnormal元の仮定の下で成立する．
以下で$K_{\mathrm{num}}(\Cat)$は，$K_0(\Cat)$をEuler形式のradicalで割った群とする．
Serre双対性とSerre作用の可逆性から左右のradicalは一致し，整数値Euler形式のradicalは飽和部分群なので，この商はtorsion-freeである．
数値同値とはこの商で同じ類を持つことをいう．

\begin{lemma}[heartの正の対象は数値的に消えない]
\label{lem:v10-heart-numerical-detection}
条件\eqref{eq:v10-SV}の下で，$F\in\Cat$が
\[
 \chi(\tau^{-n}\OO,F)=0\quad(n\gg0)
\]
を満たせば$F=0$である．特に$\Cat$の非零対象は数値的に零ではない．
\end{lemma}
\begin{proof}
条件\eqref{eq:v10-SV}により，十分大きい$n$について
\[
 0=\chi(\tau^{-n}\OO,F)
   =\dim\Hom(\tau^{-n}\OO,F).
\]
そのような$n$の下界を$N$とする．AZ1を$\tau^N F$に適用し，$\tau^{-N}$で戻すと，
$m_i\geq0$に対する全射
$\bigoplus_i\tau^{-(N+m_i)}\OO\twoheadrightarrow F$を得る．
全成分が零なので$F=0$である．
\end{proof}

\begin{proposition}[線束対象の数値類による一意性]
\label{prop:v10-numerical-line-rigidity}
六公理と条件\eqref{eq:v10-SV}を仮定する．
$I,J\in\mathcal L(\Cat)$が同じ数値類を持てば$I\simeq J$である．
従って
\[
 \mathcal L(\Cat)\longrightarrow K_{\mathrm{num}}(\Cat),\qquad I\longmapsto[I]
\]
は単射である．
\end{proposition}
\begin{proof}
$h_I(n)=\dim\Hom(\OO,\tau^nI)$と置き，$h_J$も同様に定める．
数値同値と条件\eqref{eq:v10-SV}から，$n\gg0$では
\[
 h_I(n)=\chi(\tau^{-n}\OO,I)
       =\chi(\tau^{-n}\OO,J)=h_J(n).
\]
一方，補題\ref{lem:v10-negative-line-hom}の単射$\tau^{-a}\OO\hookrightarrow I$により
\[
 h_I(n)\geq 1+\frac92(n-a)(n-a+1)\qquad(n\geq a)
\]
なので，$h_I(n)$は非有界である．

$\Ext^2(I,J)\ne0$なら，Serre双対性から$J\to\tau^{-1}I$という非零射が存在し，
これは単射である．従って全$n$で$h_J(n)\leq h_I(n-1)$となる．
十分大きい$n$での$h_I(n)=h_J(n)$と合わせると$h_I$はeventually nonincreasingとなり，
上の非有界性に反する．よって$\Ext^2(I,J)=0$である．

数値同値より$\chi(I,J)=\chi(I,I)=1$だから
\[
 \dim\Hom(I,J)=1+\dim\Ext^1(I,J)>0.
\]
非零射$f:I\to J$を取れば単射であり，そのcokernel $Q$は数値的に零である．
補題\ref{lem:v10-heart-numerical-detection}により$Q=0$なので$f$は同型である．
\end{proof}

\begin{remark}[数値的一意性が解決する範囲]
この命題は，内在的な線束対象の復元を，許される数値類の分類へ還元する．
特に，ある$I\in\mathcal L(\Cat)$と整数$m$について$[I]=[s^m\OO]$を示せれば，
条件\eqref{eq:v10-SV}の下では直ちに$I\simeq s^m\OO$となる．
しかし，全線束対象の数値類がこの軌道で尽くされることは別の問題である．
またheart内の数値的検出は，導来圏の残余圏やmutation障害$\PP$の消滅を意味しない．
複体の数値類では異なるcohomology次数の寄与が相殺し得るためである．
ここでの証明は六公理の候補圏に直接適用するものであり，標準ASクラスへの所属を前提としていない．
これらの一般性での文献上の新規性は未確認である．
\end{remark}


\begin{proposition}[反標準Hilbert関数の最高次係数とgeneric rank]
\label{prop:v10-rank-hilbert-leading}
六公理の下で，任意の$F\in\Cat$について
\[
 \dim\Hom(\OO,\tau^nF)
 =\frac92\operatorname{rk}(F)n^2+O(n)\qquad(n\to+\infty)
\]
が成立する．ここで$\operatorname{rk}(F)$は$\Cat/\Cat_1$における長さである．
この命題は三乗根や漸近Serre消滅を仮定しない．
さらに条件\eqref{eq:v10-SV}の下では，rankは$K_{\mathrm{num}}(\Cat)$上の整数値準同型に降下する．
\end{proposition}
\begin{proof}
$R=\bigoplus_{n\geq0}\Hom(\OO,\tau^n\OO)$を反標準軌道環とする．
AZ再構成により$R$はconnected right noetherianなgraded domainで，
$(\Cat,\OO,\tau)\simeq(\qgr R,\pi R,(1))$である．
また
\[
 h_R(n)=1+\frac92n(n+1)\qquad(n\geq0).
\]

まず非零斉次元全体に関する左右のgraded Ore性を，この増大度から直接確認する．
非零斉次元$a,b$の次数を$d,e$とすると，$R_n$内の部分空間
$aR_{n-d},bR_{n-e}$の交わりの次元は
\[
 h_R(n-d)+h_R(n-e)-h_R(n)=\frac92n^2+O(n)>0
\]
以上である．従って$a,b$は非零の斉次右公倍元を持つ．
$R_{n-d}a,R_{n-e}b$にも同じ議論が使え，左公倍元も存在する．
この斉次Ore局所化を$D$と書く．非零斉次元が可逆なgraded division ringである．

有限生成graded右$R$-加群$M$に対して，$M\otimes_R D$の斉次$D$-基底を
$M$の斉次元から選ぶことができる．その個数を$r$とする．
これにより単射
\[
 \bigoplus_{i=1}^rR(-a_i)\hookrightarrow M
\]
を得て，そのcokernel $T$は有限生成のgraded torsion加群となる．
$T$の斉次生成元$t_j$はそれぞれ正次数の非零斉次元$c_j$で右から消えるので，
適切なshiftを付けた$R/c_jR$の有限直和から$T$への全射がある．
従って
\[
 \dim T_n=O(n),\qquad
 \dim M_n=\frac92r n^2+O(n).
\]
ここで$\dim(R/cR)_n=h_R(n)-h_R(n-\deg c)=O(n)$を用いた．

非零斉次元$c$に対する$\pi R(-\deg c)\to\pi R$はrank one torsion-free対象間の単射であり，
generic quotientでは単純対象間の同型となる．従って$\pi(R/cR)\in\Cat_1$である．
上の有限直和表示から$\pi T\in\Cat_1$も従う．
一方，$\pi R(-a_i)$はgeneric quotientで$\OO$と同じ単純対象に移るので，
上の単射の完全列から$\operatorname{rk}(\pi M)=r$である．
この議論は，$\pi M\in\Cat_1$と$M$が斉次Ore torsionであることの同値も与える．

各$F\in\Cat$に対して，
$M=\bigoplus_{n\geq0}\Hom(\OO,\tau^nF)$は有限生成graded右$R$-加群であり，$\pi M\simeq F$である．
有限生成性を明示すれば，AZ1による有限和$\bigoplus_i\tau^{-b_i}\OO\twoheadrightarrow F$に
AZ2を適用すると，十分高い次数の$M$は有限個の$R(-b_i)$のtruncationの商となる．
$R$のright noetherian性と残りの有限個の次数の有限次元性から$M$全体も有限生成である．
AZ再構成による$\pi M\simeq F$と上の計算から，求める最高次係数を得る．

条件\eqref{eq:v10-SV}の下では，十分大きい$n$で上のHom次元は
$\chi(\tau^{-n}\OO,F)$に一致する．従って数値的に同じ$F,G$は同じrankを持つ．
$K_0$の任意の元は二つの対象の類の差として表せるので，rankはEuler radicalを消し，
$K_{\mathrm{num}}$に降下する．
\end{proof}

\begin{theorem}[数値rank三における内在的な線束列の復元]
\label{thm:v10-numerical-line-lattice}
六公理，exactな三乗根$s^3\simeq\tau$，条件\eqref{eq:v10-SV}，および
\[
 \dim_{\Q}\bigl(K_{\mathrm{num}}(\Cat)\otimes_{\Z}\Q\bigr)=3
\]
を仮定する．このとき
\[
 \mathcal L(\Cat)=\{s^n\OO:n\in\Z\}/\simeq,
 \qquad
 \Hom(s^i\OO,s^j\OO)\ne0\ \Longleftrightarrow\ i\leq j.
\]
従って線束対象のHom半順序は整数の全順序となり，任意のexact自己同値はその添字の平行移動を誘導する．
さらにexactな反Serre根の可能な正整数次数は$\{1,3\}$である．
\end{theorem}
\begin{proof}
$L_i=s^i\OO$と置く．全数値群のrankが3なので，軌道上のEuler rankは3以下である．
条件\eqref{eq:v10-SV}と定理\ref{thm:numericalrank}から，
全整数$n$について
\[
 q_n=\chi(L_0,L_n)=\frac{(n+1)(n+2)}2
\]
であり，$L_0,L_1,L_2$のEuler行列は
\[
 G=\begin{pmatrix}1&3&6\\0&1&3\\0&0&1\end{pmatrix}
\]
となる．これら三つの類は$K_{\mathrm{num}}\otimes\Q$の基底であり，
$\det G=1$なので整数基底でもある．実際，任意の数値類$v$に対し，
$\chi(L_i,v)$，$i=0,1,2$を列に並べた整数ベクトルに$G^{-1}$を掛ければ整数係数を得る．
その線形結合との差は全ての$L_i$と左直交し，従って全数値群と左直交して零となる．

$I\in\mathcal L(\Cat)$に対して
$[I]=a[L_0]+b[L_1]+c[L_2]$と書く．
命題\ref{prop:v10-rank-hilbert-leading}によりrankは数値群に降下するので
\[
 a+b+c=1.
\]
$d=b+2c\in\Z$と置き，Euler行列を用いると
\[
 1=\chi(I,I)=1+b+4c-d^2.
\]
従って
\[
 a=\frac{(d-1)(d-2)}2,\qquad
 b=d(2-d),\qquad c=\frac{d(d-1)}2.
\]
一方，$\chi(L_i,L_d)=q_{d-i}$と同じ行列$G$を用いれば，これは$[L_d]$の係数そのものである．
命題\ref{prop:v10-numerical-line-rigidity}から$I\simeq L_d$を得る．
逆に全ての$L_d$はrank one torsion-freeかつexceptionalなので，集合の等式が従う．

$i>j$なら既存の負次数Hom消滅により$\Hom(L_i,L_j)=0$である．
$i\leq j$なら$\Ext^2(L_i,L_j)=0$であるため
\[
 \dim\Hom(L_i,L_j)
 =q_{j-i}+\dim\Ext^1(L_i,L_j)>0.
\]
従ってHom半順序は$\Z$の順序と一致する．
exact自己同値はこの順序集合を保つので，整数添字を平行移動する．
$\tau=s^3$は添字を3進める．$t^r\simeq\tau$となるexact自己同値$t$の移動量を$d(t)$とすれば
$r\,d(t)=3$であり，$r=1,3$に限られる．$t=\tau,s$により両方が実現する．
\end{proof}

\begin{remark}[この復元と圏の認識との間に残る生成条件]
上の定理は，標準ASクラスへの所属を仮定せずに，そのクラスと同じ内在的な線束対象の列を復元する．
全数値群のrank3は軌道上のrank3より強い仮定であり，両者を取り違えない．
またこの定理から$\dim\Hom(L_0,L_1)=3$や，$(L_0,L_1,L_2)$のfull性が従うとはしていない．
既存の有限認識定理と合わせると，この設定で残る認識条件は
$L_3\in\thick(L_0,L_1,L_2)$，すなわち生成障害$\PP=0$である．
数値的に零のheart対象の非存在は，導来残余圏の非存在と同義ではない．
本定理の文献上の新規性は未確認である．
\end{remark}

\section{二つの障害による認識定理}\label{sec:finite}
以下，$L_i=\OO(i)$，$E_n=\Ext^1(\OO,\OO(n))$，$e_n=\dim_\C E_n$と置く．
Serre双対性と\eqref{eq:negative}から，全整数$n$について
\begin{equation}\label{eq:duality}
 \Ext^2(\OO,\OO(n))\simeq A_{-n-3}^{*},
 \qquad E_n^*\simeq E_{-n-3},
\end{equation}
ただし負次数の$A$は零とする．従って
\begin{equation}\label{eq:exceptionalcriterion}
 (L_0,L_1,L_2)\text{が例外列}\quad\Longleftrightarrow\quad E_{-1}=0.
\end{equation}
実際，後向きの差$-1,-2$では$\Hom$と$\Ext^2$は零であり，残る二つの$\Ext^1$は互いに双対である．

\begin{theorem}[固定した三乗根に対する有限判定]\label{thm:finite}
定義\ref{def:axioms}と\eqref{eq:root}の下で，$\TT=\thick(L_0,L_1,L_2)$と置く．
次の条件は同値である．
\begin{enumerate}[label=(\roman*)]
\item $E_{-1}=0$かつ$L_3\in\TT$．
\item $(L_0,L_1,L_2)$は$\DD$のfull exceptional collectionである．
\item $(L_i)_{i\in\Z}$は周期3のfull geometric helixである．
\item \eqref{eq:orbitring}の環$A$は，三次元二次Koszul AS-regular代数で，Gorenstein parameterは3，
Hilbert級数は$(1-t)^{-3}$である．
\end{enumerate}
\end{theorem}

\begin{proof}
まず(i)$\Rightarrow$(ii)を示す．\eqref{eq:exceptionalcriterion}より$\TT$は有限例外列で生成されるのでadmissibleである．
仮定より$\SSer^{-1}L_0=L_3[-2]\in\TT$である．
任意の$Y\in\TT$に対して
\[
 \Hom_{\TT}(\SSer^{-1}L_0,Y)\simeq\Hom_{\TT}(Y,L_0)^*
\]
が自然に成立するため，この対象は$\TT$の内部のinverse Serre関手の$L_0$における値でもある．
内部で例外列を回転すると
\begin{equation}\label{eq:rotation}
 \TT=\langle L_1,L_2,\SSer_{\TT}^{-1}L_0\rangle
 =\langle L_1,L_2,L_3\rangle=s\TT.
\end{equation}
よって全ての$L_i$が$\TT$に属する．

AZ生成条件だけで$(L_i)$は$\DD$を検出する．
実際，非零$X\in\DD$の非零cohomologyの最小次数を$q$とする．
ある$n\geq0$と非零射$L_{-n}\to H^q(X)$が存在する．
その$[-q]$シフトと標準切断$H^q(X)[-q]=\tau_{\leq q}X\to X$を合成した射は，
$q$次cohomology上で非零である．従って$\Hom(L_{-n},X[q])\ne0$となる．
これにより$\TT^\perp=0$であり，admissibilityから$\TT=\DD$を得る．
この段階では高次Extの漸近消滅は使っていない．

次に(ii)$\Rightarrow$(iii)．$s$で移せば全ての連続する三対象がfull exceptionalであり，
$L_{i-3}=\SSer(L_i)[-2]$である．従って$(L_i)$はtype$(3,3)$のhelixである．
全ての対象が同じheart$\Cat$に入るのでpureであり，$\phi(L_i)=i$によりsingletonからなるblock structureを持つ．
Bridgeland--Stern\cite[Definition 4.10, Theorem 5.3(ii)]{BS}により，このhelixはgeometricである．
適用に必要なfinite type性は補題\ref{lem:dimension}とExt-finitenessから従い，
algebraicityはabel圏の有界導来圏の標準dg enhancementから従う．
従って全ての$i\leq j$，$k\ne0$に対して$\Hom(L_i,L_j[k])=0$である．

(iii)$\Rightarrow$(iv)では，まず生成元数を決定する．
$a=\dim\Hom(L_0,L_1)$，$b=\dim\Hom(L_0,L_2)$とすると，$e_i=[L_i]$を基底としたEuler行列は
\[
 G=\begin{pmatrix}1&a&b\\0&1&a\\0&0&1\end{pmatrix}.
\]
$v=[L_3]=[\SSer^{-1}L_0]$に対し，Serre双対性から
$\chi(v,e_j)=\chi(e_j,e_0)$なので
\[
 v=e_0-ae_1+(a^2-b)e_2.
\]
$s$はEuler形式を保つから
$a^2-b=\chi(e_2,v)=\chi(e_1,e_2)=a$，従って$b=a^2-a$である．
geometricityと$\dim\Hom(L_0,L_3)=10$を使うと
\[
 10=\chi(e_0,v)=1-a^2+ab=1+a^3-2a^2.
\]
従って$(a-3)(a^2+a+3)=0$であり，$a=3,b=6$を得る．

Mori--Ueyama\cite[Corollary 4.2]{MU}を適用する．canonical bimoduleに対応するexact自己同値は$\theta$，
ample pairは$(\OO,s)$であり，full geometric helixは既に得られた．
よって$A$はglobal dimension 3，Gorenstein parameter 3のAS-regular代数である．
周期が$\operatorname{hdim}\Cat+1$に等しいので同系によりKoszulである．
三つの一次生成元とAS双対性によって最小分解は
\begin{equation}\label{eq:koszul}
 0\to A(-3)\to A(-2)^{\oplus3}\to A(-1)^{\oplus3}
 \to A\to\C\to0
\end{equation}
となり，$H_A(t)=(1-t)^{-3}$を得る．right noetherian性は命題\ref{prop:reconstruction}による．

(iv)$\Rightarrow$(iii)はAS-regular代数の標準helix，すなわちMori--Ueyamaの定理の逆向きとAZのtriple equivalenceによる．
(iii)$\Rightarrow$(ii)および(ii)$\Rightarrow$(i)は定義と\eqref{eq:exceptionalcriterion}から直ちに従う．
\end{proof}

\begin{corollary}[heart内の生成証明書]\label{cor:certificate}
定理\ref{thm:finite}の条件が成立すると
\begin{equation}\label{eq:certificate}
 0\to\OO\to\OO(1)^{\oplus3}\to\OO(2)^{\oplus3}\to\OO(3)\to0
\end{equation}
が$\Cat$で完全である．逆に，$E_{-1}=0$とこの形の完全列の存在を示せば定理の条件が成立する．
\end{corollary}
\begin{proof}
\eqref{eq:koszul}を$qgr$に送り，次数を3ずらす．逆向きでは完全列から$L_3\in\thick(L_0,L_1,L_2)$が従う．
\end{proof}

\begin{definition}[生成障害]\label{def:defect}
$E_{-1}=0$が成立するとき，標準の左mutation
$L_E(X)=\operatorname{Cone}(\RHom(E,X)\otimes E\to X)$を用いて
\begin{equation}\label{eq:defect}
 \PP=L_{L_0}L_{L_1}L_{L_2}(L_3)
\end{equation}
と置く．右端から順に作用させる．
\end{definition}
例外列を右端から順に取り除くので$\PP\in\TT^\perp$であり，
\[
 \PP=0\quad\Longleftrightarrow\quad L_3\in\TT.
\]
したがって定理\ref{thm:finite}は，有限次元ベクトル空間$E_{-1}$と一つの導来対象$\PP$の零性による判定である．
これは有限個のデータへの還元であり，任意の抽象圏に対する計算アルゴリズムの存在を主張しない．
特に$[\PP]=0$という$K_0$上の条件で，$\PP=0$を代用してはならない．

\subsection{非標識の圏と固定した三乗根の認識}\label{subsec:v9-marking-independence}
本稿の定義\ref{def:axioms}のrank-one条件は，既知の非可換射影平面において
構造対象の次数シフトを内在的に識別する条件と一致する．
この事実により，固定した座標環の認識と非標識のabel圏の認識との関係を明確にできる．
以下で「標準的な非可換射影平面」とは，三次元二次AS-regular代数$S$に対する
$\qgr S$と$\C$線形exact同値なabel圏をいう．これは本稿の候補公理を定義として採用することとは区別する．

\begin{lemma}[既知の非可換射影平面における直線対象と平行移動]
\label{lem:v9-intrinsic-line-translation}
$S$を三次元二次AS-regular代数，$\mathcal C=\qgr S$，
$\mathcal L_n=\pi S(n)$とする．$\mathcal C_0$を有限長対象の部分圏，
$\mathcal C_1$を$\mathcal C/\mathcal C_0$で有限長になる対象の逆像とする．
次の条件を満たす対象$I\in\mathcal C$は，ある唯一の整数$n$に対する$\mathcal L_n$である：
\[
 \Hom(\mathcal C_1,I)=0,\qquad
 I\text{の}\mathcal C/\mathcal C_1\text{における長さは}1,\qquad
 \chi(I,I)=1.
\]
さらに任意の$\C$線形exact自己同値$F$に対して，ある整数$d(F)$が存在し，
全ての$n\in\Z$で$F(\mathcal L_n)\simeq\mathcal L_{n+d(F)}$となる．
$d(FG)=d(F)+d(G)$である．
\end{lemma}
\begin{proof}
第一の主張はAbdelgadir--Okawa--Ueda\cite[Corollary A.7]{AOUCompact}である．
同文献の$\mathcal C_i$は有限長のSerre filtrationで定義されており，
本稿の$\Cat_0,\Cat_1$と同じである．また
\[
 \Hom_{\mathcal C}(\mathcal L_i,\mathcal L_j)\ne0
 \quad\Longleftrightarrow\quad i\leq j,
 \qquad
 \dim\Hom_{\mathcal C}(\mathcal L_i,\mathcal L_{i+r})
 =\binom{r+2}{2}\quad(r\geq0).
\]
特に$\mathcal L_i\simeq\mathcal L_j$なら$i=j$である．
exact自己同値は有限長対象，そのSerre商，対象の長さ，およびEuler形式を保つ．
従って$F$は上記の内在的な対象集合を置換し，Homの非消滅関係によってその順序を保つ．
$\Z$の順序を保つ全単射は平行移動なので，$F$は$n\mapsto n+d(F)$を誘導する．
合成についての式も直ちに従う．
\end{proof}

\begin{proposition}[認識の標識への非依存性]\label{prop:v9-marking-independence}
定義\ref{def:axioms}を満たす$(\Cat,\OO)$を固定する．
\begin{enumerate}[label=(\roman*)]
\item $\Cat$が標準的な非可換射影平面なら，heart上でexactな三乗根
$s^3\simeq\SSer^{-1}[2]$が存在する．
\item $\Cat$が標準的な非可換射影平面なら，そのような三乗根$s$をどれに取っても，
$(\OO,s\OO,s^2\OO)$はfull strong exceptional collectionであり，
\[
 A_s=\bigoplus_{n\geq0}\Hom(\OO,s^n\OO)
\]
は三次元二次Koszul AS-regular代数となる．
\item 従って，exactな三乗根$s$が固定されたとき，次は同値である：
\[
 \begin{aligned}
 &\Cat\text{が標準的な非可換射影平面}\\
 &\quad\Longleftrightarrow\ A_s\text{が三次元二次AS-regular}\\
 &\quad\Longleftrightarrow\ E_{-1}=0\text{かつ }s^3\OO\in\thick(\OO,s\OO,s^2\OO).
 \end{aligned}
\]
\end{enumerate}
この命題にnormal元の存在や境界の仮定は不要である．
\end{proposition}
\begin{proof}
まず$\Cat\simeq\qgr S$を仮定する．この同値は有限長Serre filtrationを保つ．
公理(3)と公理(4)から，$\OO$の像は補題\ref{lem:v9-intrinsic-line-translation}の三条件を満たす．
従って$\OO$の像はある$\mathcal L_m=\pi S(m)$である．

(i)を示す．Itaba--Matsuno\cite[Theorems 1.1, 4.4]{IM}により，$S$は
三次元二次Calabi--Yau AS-regular代数$T$とgraded Morita同値である．
同論文Lemma 2.5により，この同値はtwisting systemによる同値として選べるため，
有限次元加群を保ち，$\qgr S\simeq\qgr T$というexact同値を誘導する．
その下で$\OO$の像は，補題\ref{lem:v9-intrinsic-line-translation}を$T$に適用することにより
$\pi T(a)$となる．同値に次数シフト$(-a)$を合成し，$\OO$が$\pi T$へ移るように正規化する．
この正規化には，graded Morita同値がもとの全shift関手を自然に絡めることは要求しない．

$T$はCalabi--YauでGorenstein parameterが3なので，$\Db(\qgr T)$のSerre関手は
自然同型として$(-3)[2]$である．従って$\qgr T$の標準shift$(1)$を上のexact同値で移せば，
$\Cat$のexact自己同値$s_0$で$s_0^3\simeq\SSer^{-1}[2]$を満たすものが得られる．
Serre関手を移す際の同型は，Serre関手の一意性による自然同型である．

(ii)では，$s$を任意のexactな三乗根とする．$\tau=\SSer^{-1}[2]$と置く．
$\qgr S$へ移された$s,\tau$はいずれもexact自己同値なので，
補題\ref{lem:v9-intrinsic-line-translation}により，直線対象の添字をそれぞれ整数$d,e$だけ平行移動する．
公理(6)の$n=1$は
\[
 10=\dim\Hom(\OO,\tau\OO)
   =\dim\Hom(\mathcal L_m,\mathcal L_{m+e})
\]
を与える．右辺は$e<0$なら零，$e\geq0$なら$\binom{e+2}{2}$なので，$e=3$である．
$s^3\simeq\tau$より$3d=e=3$，従って$d=1$となる．よって
\[
 (\OO,s\OO,s^2\OO)\longmapsto
 (\mathcal L_m,\mathcal L_{m+1},\mathcal L_{m+2}).
\]
右辺は標準のfull strong exceptional collectionである
\cite[Theorem 7.1]{AOUCompact}．従って左辺もそうであり，
定理\ref{thm:finite}により$A_s$は三次元二次Koszul AS-regularとなる．

(iii)の後半の同値は定理\ref{thm:finite}である．$A_s$がAS-regularなら，
AZ再構成による$\Cat\simeq\qgr A_s$から$\Cat$は標準的な非可換射影平面である．
逆向きは(ii)で証明した．
\end{proof}

\begin{remark}[対象の認識と標識の分類]\label{rem:v9-marking-status}
この命題は，公理(3)のrank-one条件と公理(4)の例外性を保った範囲で，
一つの許容される三乗根の認識障害が非零なら，$\Cat$自体が標準的な非可換射影平面ではないことを意味する．
一方，同じ圏を与えるgraded algebraやexact三乗根が一意になるとは主張しない．
Zhang twistによる座標環の違いと，非標識のabel圏の認識は区別される．
また，$\OO$を任意の例外対象へ変更してrank-one条件を失う場合や，
heartを保たない導来自己同値だけを扱う場合には，本命題は適用できない．
\end{remark}

\begin{corollary}[標準的な非可換射影平面のexactな圏論的指数]
\label{cor:v9-exact-index}
$\mathcal C$を標準的な非可換射影平面とし，
$\tau=\SSer^{-1}[2]$とする．heart上exactな自己同値による根の次数の集合を
\[
 R_{\mathrm{ex}}(\mathcal C)
 :=\{r\in\Z_{\geq1}\mid
       \text{あるexact自己同値 }t:\mathcal C\to\mathcal C
       \text{ に対し }t^r\simeq\tau\}
\]
と定める．ここで同型は$\Db(\mathcal C)$上の関手の自然同型を意味する．
このとき
\[
 R_{\mathrm{ex}}(\mathcal C)=\{1,3\},\qquad
 i_{\mathrm{ex}}(\mathcal C):=\max R_{\mathrm{ex}}(\mathcal C)=3.
\]
この主張は六公理を追加せず，標準的な非可換射影平面そのものに対して成立する．
\end{corollary}
\begin{proof}
Itaba--Matsuno\cite[Theorem 4.4]{IM}により，
$\mathcal C\simeq\qgr T$となる三次元二次Calabi--Yau AS-regular代数$T$を選べる．
この表示では$\tau\simeq(3)$である．
補題\ref{lem:v9-intrinsic-line-translation}を$T$に適用すると，
全てのexact自己同値$t$は直線対象$\pi T(n)$の添字を整数$d(t)$だけ平行移動する．
従って$d(\tau)=3$であり，$t^r\simeq\tau$なら
\[
 r\,d(t)=d(t^r)=d(\tau)=3.
\]
$r$は正整数，$d(t)$は整数だから，$r=1$または$r=3$でなければならない．
逆に$r=1$は$t=\tau$によって，$r=3$はこのCalabi--Yau表示の次数シフト$(1)$を
$\mathcal C$へ移送することによって実現する．
\end{proof}

\begin{remark}[指数の帰結の範囲]\label{rem:v9-exact-index-scope}
系\ref{cor:v9-exact-index}は，AOUの直線対象の内在的特徴付けと
Itaba--MatsunoのCalabi--Yau表示を組み合わせた短い帰結であり，
ここでは独立した新規性を主張しない．
これは標準的な非可換射影平面がexactな圏論的指数3を持つことを示すが，
一般の六公理の候補圏に対する指数の上界や，指数3からの認識を証明するものではない．
また，変形族全体で根を互換的に選べるという主張も含まない．
\end{remark}

\subsection{残余回転と生成障害の軌道}\label{subsec:residualrotation}
生成障害$\PP$は，残余圏全体を一つの自己同値の軌道によって生成する．
その自己同値とSerre関手の関係は，Kuznetsov--Smirnovの残余偏極の標準構成
\cite[Theorem 2.8 and its proof]{KS}の特殊化である．
同文献は飽和圏を仮定するが，以下ではその仮定を追加せず，
本稿に必要な自己同値とSerre双対性を射影三角形から直接示す．
さらにAZ生成条件を用いて，$\PP$の軌道が残余圏を厚く生成することを示す．
この生成の記述について，文献に対する新規性は未確認である．

\begin{lemma}[軌道対象による厚い生成]\label{lem:orbitthickgeneration}
基礎公理の下で
\[
 \DD=\thick\{L_{-3m}:m\geq0\}
      =\thick\{L_n:n\in\Z\}.
\]
より具体的に，任意の$M\in\Cat$は，各項が$L_{-3m}$の有限直和である
三項複体の直和因子である．この主張には高次Extの漸近消滅を用いない．
\end{lemma}
\begin{proof}
AZ生成条件を三回適用し，$P_i$を$L_{-3m}$の有限直和として完全列
\[
 0\longrightarrow K_3\longrightarrow P_2\longrightarrow P_1
 \longrightarrow P_0\longrightarrow M\longrightarrow0
\]
を作る．$P_2,P_1,P_0$を次数$-2,-1,0$に置いた複体を$Q$とする．
標準切断による三角形は
\[
 K_3[2]\longrightarrow Q\longrightarrow M
 \longrightarrow K_3[3]
\]
である．最後の射は$\Ext^3(M,K_3)=0$に属するので，この三角形は分裂する．
従って$Q\simeq K_3[2]\oplus M$であり，$M$は上記の厚い部分圏に属する．
有界複体は有限個のcohomology対象から切断三角形によって作られるので，結論を得る．
\end{proof}

以後この小節では$E_{-1}=0$を仮定する．
\[
 \mathcal R=\TT^\perp,
 \qquad \mathsf F=L_{L_0}L_{L_1}L_{L_2}:\DD\longrightarrow\mathcal R
\]
と置き，包含を$i:\mathcal R\hookrightarrow\DD$と書く．
半直交分解の向きは
\begin{equation}\label{eq:residualsod}
 \DD=\langle\mathcal R,\TT\rangle
     =\langle\mathcal R,L_0,L_1,L_2\rangle,
 \qquad \RHom(\TT,\mathcal R)=0
\end{equation}
である．各$X\in\DD$には自然な射影三角形
\begin{equation}\label{eq:residualprojection}
 T_X\longrightarrow X\longrightarrow i\mathsf F X\longrightarrow T_X[1],
 \qquad T_X\in\TT
\end{equation}
がある．この三角形に$\RHom(-,iY)$，$Y\in\mathcal R$を適用すれば，
$\mathsf F$が$i$の左随伴であることが分かる．以下では包含$i$を必要な場合以外は省略する．
特に$\mathsf F L_0=\mathsf F L_1=\mathsf F L_2=0$，$\mathsf F L_3=\PP$である．

\begin{theorem}[残余回転と一対象の軌道生成]\label{thm:residualrotation}
関手
\begin{equation}\label{eq:residualrotation}
 \rho=L_{L_0}\circ s\big|_{\mathcal R}
\end{equation}
は$\mathcal R$の自己同値である．右mutationを
\[
 R_{L_0}(X)=
 \operatorname{Cone}\bigl(X\longrightarrow
 \RHom(X,L_0)^*\otimes L_0\bigr)[-1]
\]
と正規化すると，逆関手は
\begin{equation}\label{eq:residualrotationinverse}
 \rho^{-1}\simeq s^{-1}\circ R_{L_0}\big|_{\mathcal R}
\end{equation}
で与えられる．また$\mathcal R$はSerre関手$\SSer_{\mathcal R}$を持ち，
\begin{equation}\label{eq:residualserre}
 \rho^3\simeq\SSer_{\mathcal R}^{-1}[2]
\end{equation}
である．生成障害は次の意味で残余圏全体を生成する：
\begin{equation}\label{eq:residualorbitgeneration}
 \mathcal R=\thick\{\rho^k\PP:k\in\Z\}.
\end{equation}
\end{theorem}
\begin{proof}
\emph{回転が残余圏を保つこと．}
$X\in\mathcal R$なら$sX$は$L_1,L_2,L_3$の右直交に入る．
左mutationの三角形
\[
 \RHom(L_0,sX)\otimes L_0\longrightarrow sX
 \longrightarrow L_{L_0}(sX)
\]
において，最初の二対象はいずれも$L_1,L_2$の右直交に入る．
最初の対象については例外列の半直交性$\RHom(L_j,L_0)=0$，$j=1,2$を使う．
従って最後の対象も$L_1,L_2$の右直交に入り，mutationの定義により$L_0$にも右直交である．
これで$\rho X\in\mathcal R$が従う．

\emph{逆関手．}
$X\in\mathcal R$について$Y=R_{L_0}X$と置く．
右mutationの三角形
\[
 Y\longrightarrow X\longrightarrow\RHom(X,L_0)^*\otimes L_0
\]
から，$Y\in{}^\perp L_0$かつ$\RHom(L_1,Y)=\RHom(L_2,Y)=0$である．
さらに$\SSer L_3\simeq L_0[2]$なので，Serre双対性から
$\RHom(L_3,Y)=0$も従う．従って$s^{-1}Y\in\mathcal R$であり，
式\eqref{eq:residualrotationinverse}の右辺は$\mathcal R$の自己関手になる．

一方，$X\in\mathcal R$なら$\RHom(L_3,sX)=0$であり，
$\SSer L_3=L_0[2]$から$\RHom(sX,L_0)=0$となる．
したがって$sX\in{}^\perp L_0$である．
例外対象を通る左右mutationは
\[
 L_{L_0}:{}^\perp L_0\longrightarrow L_0^\perp,
 \qquad R_{L_0}:L_0^\perp\longrightarrow{}^\perp L_0
\]
の間の互いに逆な同値である．これは二つのevaluation三角形に
$\RHom(L_0,-)$および$\RHom(-,L_0)$を適用して得られる標準のmutation恒等式である．
従って$X\in\mathcal R$について自然に
\[
 s^{-1}R_{L_0}L_{L_0}sX\simeq X,
 \qquad L_{L_0}s\bigl(s^{-1}R_{L_0}X\bigr)\simeq X
\]
となり，逆関手の主張が従う．

\emph{内部Serre関手との関係．}
自己同値によるmutationの移送
$sL_{L_j}s^{-1}\simeq L_{L_{j+1}}$から
\begin{equation}\label{eq:residualcubecalculation}
 \rho^3\simeq L_{L_0}L_{L_1}L_{L_2}s^3\big|_{\mathcal R}
       =\mathsf F s^3i
       \simeq\mathsf F\SSer^{-1}i[2].
\end{equation}
左随伴の性質とambient Serre双対性により，任意の$X,Y\in\mathcal R$について
\[
 \Hom_{\mathcal R}(\mathsf F\SSer^{-1}iX,Y)
 \simeq\Hom_{\DD}(\SSer^{-1}iX,iY)
 \simeq\Hom_{\mathcal R}(Y,X)^*.
\]
しかも式\eqref{eq:residualcubecalculation}により
$\mathsf F\SSer^{-1}i\simeq\rho^3[-2]$は自己同値である．
よってこれは$\mathcal R$のinverse Serre関手であり，式\eqref{eq:residualserre}を得る．
この議論は$\DD$の飽和性やsmoothnessを用いない．

\emph{生成障害の軌道による生成．}
$\mathcal U=\thick\{\rho^k\PP:k\in\Z\}\subset\mathcal R$と置く．
式\eqref{eq:residualprojection}に$s$，次いで$\mathsf F$を適用すると三角形
\begin{equation}\label{eq:residualrecursion}
 \mathsf F(sT_X)\longrightarrow\mathsf F(sX)
 \longrightarrow\rho(\mathsf F X)\longrightarrow\mathsf F(sT_X)[1]
\end{equation}
を得る．ここで$s\mathsf F X$は$L_1,L_2$の右直交に入るので，
$\mathsf F(s\mathsf F X)\simeq L_{L_0}s\mathsf F X=\rho\mathsf F X$を使った．
また$sT_X\in\thick(L_1,L_2,L_3)$だから
\[
 \mathsf F(sT_X)\in\thick(\mathsf FL_1,\mathsf FL_2,\mathsf FL_3)
 =\thick(\PP)\subset\mathcal U.
\]
従って式\eqref{eq:residualrecursion}を$X=L_n$に適用すれば，
$\mathsf F L_n\in\mathcal U$と$\mathsf F L_{n+1}\in\mathcal U$は同値である．
$\mathcal U$が$\rho$と$\rho^{-1}$の両方で保たれることを使っている．
$\mathsf F L_0=0$を出発点として，全整数$n$について$\mathsf F L_n\in\mathcal U$を得る．
補題\ref{lem:orbitthickgeneration}と$\mathsf F$のexact性から
$\mathsf F(\DD)\subset\mathcal U$である．
$\mathsf F i\simeq\id_{\mathcal R}$なので$\mathcal R\subset\mathcal U$となり，結論が従う．
\end{proof}

\begin{remark}\label{rem:residualstatus}
式\eqref{eq:residualorbitgeneration}は$\PP=0$なら残余圏全体が零になる理由を明確にするが，
$\PP$が零であること自体の証明ではない．また，$\rho^3=\SSer_{\mathcal R}^{-1}[2]$は
$\mathcal R$が2-Calabi--Yauであることを意味しない．
$\rho$があるbounded heartを保つことや，その軌道がAZ ample pairをなすことも，ここでは主張していない．
\end{remark}

\subsection{二重mutationの純粋性による生成判定}\label{sec:mutationpurity}
$E_{-1}=0$を仮定し，
\[
 Q_2=L_{L_2}(L_3),\qquad Q=L_{L_1}(Q_2),\qquad
 \PP=L_{L_0}(Q)
\]
と置く．左mutationの正規化は主原稿と同じ
$L_E(X)=\operatorname{Cone}(\RHom(E,X)\otimes E\to X)$である．

\begin{lemma}\label{lem:mutationamplitude}
上の設定では
\[
 Q_2\in\DD^{[-1,0]},\qquad Q\in\DD^{[-2,0]},\qquad
 \PP\in\DD^{[-3,0]}.
\]
また$Q$はexceptionalであり，
\begin{equation}\label{eq:mutationcomparison}
 \RHom(Q,L_0)\simeq\C[-2].
\end{equation}
$H^{-2}(Q)$は$L_1$の有限直和の部分対象で，従ってtorsion-freeである．
\end{lemma}
\begin{proof}
$i<j$なら$\RHom(L_i,L_j)$の非零cohomologyは次数0,1に限られる．
負次数はheartの性質，次数2はSerre双対性と負次数Hom消滅，
それより高い次数は$\operatorname{hdim}\Cat=2$によって消える．
mutation三角形と，それに$\RHom(L_i,-)$を適用した三角形から，順に
\[
\begin{array}{c|c}
\text{対象または複体}&\text{cohomologyの範囲}\\ \hline
\RHom(L_2,L_3)&[0,1]\\
Q_2&[-1,0]\\
\RHom(L_i,Q_2)\ (i=0,1)&[-1,1]\\
Q&[-2,0]\\
\RHom(L_0,Q)&[-2,1]\\
\PP&[-3,0]
\end{array}
\]
を得る．例えば$\RHom(L_1,Q_2)$は，次数$[0,2]$にある
$\RHom(L_1,L_2)\otimes\RHom(L_2,L_3)$から
次数$[0,1]$にある$\RHom(L_1,L_3)$への射のconeなので，次数$[-1,1]$にある．
$i=0$の場合も同様である．
次いで$\RHom(L_0,Q)$は，次数$[-1,2]$にある
$\RHom(L_0,L_1)\otimes\RHom(L_1,Q_2)$から
次数$[-1,1]$にある$\RHom(L_0,Q_2)$への射のconeなので，次数$[-2,1]$にある．
従ってこの計算はforward $\Ext^1$の消滅を仮定していない．
また第二のmutation三角形の長完全列から
\[
 H^{-2}(Q)\hookrightarrow
 H^{-1}\RHom(L_1,Q_2)\otimes L_1
\]
が得られる．

$(L_1,L_2,L_3)$は$(L_0,L_1,L_2)$を$s$で移した例外列である．
例外対の左mutationは例外性を保つので，二回移した$Q$もexceptionalである．
この基本事実は，例外対$(E,X)$に対して$\RHom(X,E)=0$と
$\RHom(E,L_EX)=0$を二つのmutation三角形に用い，
$\RHom(L_EX,L_EX)\simeq\RHom(X,X)$を得れば確認できる．
最後に$\RHom(L_1,L_0)=\RHom(L_2,L_0)=0$なので，
二つの三角形に$\RHom(-,L_0)$を適用して
$\RHom(Q,L_0)\simeq\RHom(L_3,L_0)$を得る．
$\SSer L_3=L_0[2]$と$L_0$の例外性により，これは$\C[-2]$である．
\end{proof}

\begin{proposition}[純粋性から生成障害の消滅へ]\label{prop:puritykillsdefect}
基礎公理と三乗根に加え，次の二条件を仮定する：
\[
 \chi(\OO,\OO(n))=\frac{(n+1)(n+2)}2\quad(n\in\Z),\qquad
 \Ext^{>0}(L_{-n},F)=0\quad(n\gg0,\ F\in\Cat).
\]
このとき
\[
 \PP=0\quad\Longleftrightarrow\quad Q[-2]\in\Cat.
\]
特に右辺の純粋性を示せば，固定した軌道環のAS-regular認識が従う．
ここでは純粋性自体が公理から従うとは主張していない．
\end{proposition}
\begin{proof}
Euler式の$n=-1$から$E_{-1}=0$であり，上の補題が使える．
$\PP=0$なら$Q\in\thick(L_0)$である．
$Q$はexceptionalなので$Q\simeq L_0[a]$となる整数$a$があり，
式\eqref{eq:mutationcomparison}から$a=2$である．従って右辺を得る．

逆に$F=Q[-2]\in\Cat$とする．補題から$F=H^{-2}(Q)$はtorsion-freeである．
またmutationのEuler公式と$q_1=3,q_2=6$から
\[
 [F]=[Q]=[L_3]-3[L_2]+3[L_1].
\]
$\dim\Hom(L_0,L_1)=3+\dim E_1>0$なので，$L_0,L_1,L_2,L_3$のgeneric quotientでの像は
全て同じ単純対象である．$F$の像は$L_1$の像の有限直和の部分対象なので有限長であり，
mutation三角形をgeneric quotientへ移し，有限長cohomologyの長さの交代和を取れば，
上の等式からその長さは$1-3+3=1$となる．
従って$F$はgeneric rank oneである．

式\eqref{eq:mutationcomparison}は$\Hom(F,L_0)\simeq\C$を与える．
非零射$f:F\to L_0$を取る．$L_0$のtorsion-free性から$f$のgeneric quotientでの像は非零であり，
単純対象間の同型になる．従って$\ker f\in\Cat_1$であり，
$F$のtorsion-free性から$\ker f=0$である．
余核を$C$として$0\to F\to L_0\to C\to0$を得る．

$P(n)=(n+1)(n+2)/2$と書くと，全整数$n$に対して
\[
 \chi(L_{-n},F)=P(n+3)-3P(n+2)+3P(n+1)=P(n).
\]
従って$\chi(L_{-n},C)=0$である．
全対象のSerre消滅を$C$に適用すれば，十分大きな$n$で
$\Hom(L_{-n},C)=0$となる．
AZ生成条件を$\tau^M C$に適用し，再び$\tau^{-M}$で移すと，
$M$を任意に大きくして$L_{-3(M+a_j)}$の有限直和から$C$への全射を得る．
その全ての成分射が零なので$C=0$である．
よって$F\simeq L_0$，$Q\simeq L_0[2]$となり，$\PP=0$が従う．
\end{proof}

\paragraph{最初のmutationで何が必要か．}
$U=\Hom(L_2,L_3)$，$V=\Ext^1(L_2,L_3)=E_1$とし，
$\mathrm{ev}:U\otimes L_2\to L_3$をevaluationとする．
最初のmutation三角形から
\[
 H^{-1}(Q_2)=\ker(\mathrm{ev}),\qquad
 0\to\operatorname{coker}(\mathrm{ev})\to H^0(Q_2)
 \to V\otimes L_2\to0
\]
を得る．従って$Q_2[-1]$がheartに入ることは，
$E_1=0$とevaluationの全射性の二条件に等しい．
全整数のEuler二次性はこの二条件を単独では与えない．
二重mutation$Q[-2]$の純粋性を示すためには，このような実際のevaluationの核・余核を制御する必要がある．

\section{全対象のSerre消滅を制御する障害加群}\label{sec:cohomology}
\begin{lemma}\label{lem:topvanishing}
全ての$F\in\Cat$に対し$\Ext^2(\OO(-n),F)=0$ for $n\gg0$である．
\end{lemma}
\begin{proof}
AZ生成条件により固定した全射$P=\bigoplus_i\OO(-a_i)\twoheadrightarrow F$を取る．
Serre双対性と前合成により
\[
 \Ext^2(\OO(-n),F)^*\simeq\Hom(F,\OO(-n-3))
 \hookrightarrow\bigoplus_i\Hom(\OO,\OO(a_i-n-3))
\]
となる．右辺は$n$が十分大きければ補題\ref{lem:mono}により零である．
\end{proof}

\begin{proposition}\label{prop:serrevanishing}
$E=\bigoplus_{n\in\Z}E_n$と置く．次は同値である．
\begin{enumerate}[label=(\roman*)]
\item $E_n=0$ for $n\gg0$．
\item 全ての$F\in\Cat$に対して$\Ext^i(\OO,s^nF)=0$ for $i>0,n\gg0$．
\item $\dim_\C E<\infty$．
\item $E_{3m}=0$ for $m\gg0$．
\end{enumerate}
\end{proposition}
\begin{proof}
(ii)$\Rightarrow$(i)は明らかであり，(i)$\Leftrightarrow$(iii)は\eqref{eq:duality}と各次数の有限次元性による．
(i)を仮定し，$0\to K\to P\to F\to0$を$P=\bigoplus_i\OO(-a_i)$で取る．
\[
 \Ext^1(\OO(-n),P)\to\Ext^1(\OO(-n),F)
 \to\Ext^2(\OO(-n),K)
\]
の左端は(i)，右端は補題\ref{lem:topvanishing}により消える．
$i=2$も同補題，$i>2$はホモロジー次元による．

(iii)$\Rightarrow$(iv)は明らかである．(iv)を仮定し，$r=1,2$について
元の$\tau$-ample性から固定した完全列
\[
 0\to K_r\to\bigoplus_j\OO(-3a_{rj})\to\OO(r)\to0
\]
を取る．$s^{3m}$で移して長完全列を取ると
\[
 \bigoplus_jE_{3(m-a_{rj})}\longrightarrow E_{3m+r}
 \longrightarrow\Ext^2(\OO,s^{3m}K_r)
\]
を得る．左端は(iv)，右端は
$\Ext^2(\OO,s^{3m}K_r)=\Ext^2(\OO(-3m),K_r)$と
補題\ref{lem:topvanishing}により$m\gg0$で零になる．
三つの剰余類を合わせて(i)が従う．
\end{proof}

AZのコホモロジー判定\cite[Corollary 7.5]{AZ}により，これらは$A$のfull right $\chi$条件とも同値である．
再構成だけで得られるのは$\chi_1$であり，命題\ref{prop:serrevanishing}の条件を公理に紛れ込ませてはならない．
本命題の四条件の直接証明は，この文献上の同定を使っていない．

\begin{proposition}[各元のtorsion性]\label{prop:torsion}
$E_n\simeq\Ext^1(\OO(-n),\OO)$と同定し，斉次切断の前合成で$E$を次数付き右$A$加群とする．
各斉次元$e\in E$は，十分高次の全ての元によって消される．
\end{proposition}
\begin{proof}
$e\in E_n$を$0\to\OO\to M\to\OO(-n)\to0$で表す．
AZの全射性条件から
\[
 \Hom(\OO(-n-r),M)\longrightarrow\Hom(\OO(-n-r),\OO(-n))
\]
は$r\gg0$で全射である．その連結射は拡大の引戻しだから，$eA_r=0$ for $r\gg0$となる．
\end{proof}

各元に対するこの証明だけでは共通の消滅指数は得られない．
しかし，次の左有限生成性と組み合わせると非負次数部分について一様化できる．

\subsection{正次数部分の左有限生成と一様な右消滅}

右torsion性の一様性を調べるには，$E$の左右の作用を区別する必要がある．
$a\in A_d$，$b\in A_m$，$e\in E_n$を，それぞれ
$a\colon\OO\to\OO(d)$，$b\colon\OO\to\OO(m)$，
$e\colon\OO\to\OO(n)[1]$という導来圏の射と見なす．
軌道環の積と両側の作用を
\begin{equation}\label{eq:E-biaction}
 a*b=s^m(a)\circ b,\qquad
 a*e=s^n(a)[1]\circ e,\qquad
 e*b=s^m(e)\circ b
\end{equation}
で定める．いずれも内部次数は加法的である．
自己同値の冪の自然な同一視を固定し，以下ではその結合同型を省略する．

\begin{lemma}\label{lem:E-bimodule}
式\eqref{eq:E-biaction}は$E$を次数付き$A$双加群にする．
特に，左右の作用は可換であり，
\[
 (a*e)*b=a*(e*b)
\]
が成り立つ．
\end{lemma}
\begin{proof}
上記の次数で，三つの結合則はそれぞれ次の同じ合成を表す．
\begin{align*}
 (a*b)*e=a*(b*e)
   &=s^{m+n}(a)[1]\circ s^n(b)[1]\circ e,\\
 (e*a)*b=e*(a*b)
   &=s^{d+m}(e)\circ s^m(a)\circ b,\\
 (a*e)*b=a*(e*b)
   &=s^{m+n}(a)[1]\circ s^m(e)\circ b.
\end{align*}
単位は$\id_{\OO}$である．$a,b$のコホモロジー次数は0なので，ここにはKoszul符号は現れない．
\end{proof}

\begin{proposition}[正次数部分の左有限生成]\label{prop:E-left-finite}
基礎公理とexactな三乗根の下で，
\[
 E_{\geq0}:=\bigoplus_{n\geq0}E_n
\]
は左$B=A^{(3)}$加群として有限生成であり，従って左$A$加群としても有限生成である．
ここでは$B\subset A$に$A$からの内部次数を入れる．
より具体的に，正の3の倍数$d_1,\dots,d_r$，$a_i\in A_{d_i}\subset B$，整数$N_0$を取って
\begin{equation}\label{eq:E-left-recursion}
 E_n=\sum_{i=1}^r a_i*E_{n-d_i}
 \qquad(n\geq N_0)
\end{equation}
とできる．
\end{proposition}
\begin{proof}
まず全射の全成分を真に負の反標準twistに取れることを確認する．
元の$(\OO,\tau)$に関するAZ生成条件を$\tau\OO=\OO(3)$に適用すると
\[
 \bigoplus_{i=1}^r\tau^{-m_i}\OO\twoheadrightarrow\tau\OO,
 \qquad m_i\geq0
\]
を得る．これを$\tau^{-1}$でtwistし，$d_i=3(m_i+1)>0$と置けば
\begin{equation}\label{eq:E-negative-cover}
 0\longrightarrow K\longrightarrow
 P:=\bigoplus_{i=1}^r\OO(-d_i)
 \xrightarrow{\ p\ }\OO\longrightarrow0
\end{equation}
が得られる．$p_i\colon\OO(-d_i)\to\OO$を成分とし，
$a_i=s^{d_i}(p_i)\in A_{d_i}$とする．

式\eqref{eq:E-negative-cover}に$\Ext^\bullet(\OO(-n),-)$を適用すると
\[
 \bigoplus_i\Ext^1(\OO(-n),\OO(-d_i))
 \longrightarrow\Ext^1(\OO(-n),\OO)
 \longrightarrow\Ext^2(\OO(-n),K)
\]
が完全である．右端は固定した$K$に補題\ref{lem:topvanishing}を適用して，
ある$N_0$以上の全ての$n$で零になる．
最初の二項を$s^n$でtwistして同定すると，$i$番目の成分の写像は
\[
 E_{n-d_i}\longrightarrow E_n,
 \qquad u\longmapsto s^{n-d_i}(a_i)[1]\circ u=a_i*u
\]
である．従って式\eqref{eq:E-left-recursion}が成り立つ．
ここで使用したのは最上位Extの既証明の漸近消滅であり，$E$自身の消滅ではない．

$D=\max_i d_i$，$b_0=\max\{N_0,D,1\}$と置く．
$n\geq b_0$なら，各$n-d_i$は$0\leq n-d_i<n$を満たす．
式\eqref{eq:E-left-recursion}による$n$についての帰納法から，
$E_{\geq0}$は有限次元部分空間$\bigoplus_{0\leq j<b_0}E_j$で左$B$生成される．
全ての$a_i$が$B$に属することを使った．各$E_j$の有限次元性はExt-finitenessによる．
\end{proof}

\begin{proposition}[一様な右消滅指数]\label{prop:E-uniform-right}
ある正整数$N$が存在し，
\begin{equation}\label{eq:E-uniform-right}
 E_{\geq0}*A_{\geq N}=0
\end{equation}
が成り立つ．$N$は$E_{\geq0}$の元とその内部次数に依存しない．
\end{proposition}
\begin{proof}
命題\ref{prop:E-left-finite}により有限個の斉次左生成元
$u_1,\dots,u_t$を取る．命題\ref{prop:torsion}の右作用は
式\eqref{eq:E-biaction}の右作用と一致する．
実際，$u_j\in E_{r_j}$を
$0\to\OO(r_j)\to X_j\to\OO\to0$で表し，
AZの全射性条件を$X_j\twoheadrightarrow\OO$に適用すると，
十分大きい$m$に対して全ての$\OO(-m)\to\OO$への引戻しが零になる．
$s^m$でtwistしたこの引戻しが，ちょうど$u_j*A_m$である．

従って各$j$に$u_j*A_{\geq N_j}=0$となる$N_j$がある．
$N=\max\{1,N_1,\dots,N_t\}$とする．
$x\in E_{\geq0}$を有限和$x=\sum_j a_j*u_j$で表せば，
任意の$b\in A_{\geq N}$に対し補題\ref{lem:E-bimodule}から
\[
 x*b=\sum_j(a_j*u_j)*b
     =\sum_j a_j*(u_j*b)=0
\]
となる．
\end{proof}

\begin{corollary}[残る有限性の左右差]\label{cor:E-left-torsion}
次は同値である．
\begin{enumerate}[label=(\roman*)]
\item $\dim_\C E<\infty$．
\item $E_{\geq0}$は左torsionである：全ての斉次元$u\in E_{\geq0}$について，
ある$q(u)$が存在して$A_{\geq q(u)}*u=0$となる．
\end{enumerate}
また，基礎公理だけで$\dim E_n=O(n^2)$ as $n\to+\infty$が成り立つ．
\end{corollary}
\begin{proof}
(i)なら$E$は有限個の内部次数にしか成分を持たないので(ii)が従う．
(ii)を仮定し，命題\ref{prop:E-left-finite}の斉次左生成元$u_j$を取る．
$A_{\geq q_j}*u_j=0$とすれば，$A*u_j$は
$A/A_{\geq q_j}$の次数シフトの商であり，有限次元である．
有限個の和を取って$E_{\geq0}$は有限次元となる．
Serre対称性$E_n^*\simeq E_{-n-3}$と各次数の有限次元性から$E$全体も有限次元である．

増大度について，$r_j=\deg u_j$とすると
\[
 \dim E_n\leq\sum_{j=1}^t\dim A_{n-r_j}.
\]
右辺は命題\ref{prop:growth}により$O(n^2)$である．
\end{proof}

\begin{remark}\label{rem:E-sides-gap}
命題\ref{prop:E-left-finite}は正次数部分の\emph{左}有限生成であり，
命題\ref{prop:E-uniform-right}はその\emph{右}作用の一様消滅である．
系\ref{cor:E-left-torsion}(ii)の左torsion性は，これらからはまだ導かれていない．
従って，ここでは命題\ref{prop:serrevanishing}のSerre消滅やfull right $\chi$を証明したとはしない．
この言換えは新しい公理の追加ではなく，元の公理からの有限性の証明に残る箇所の同定である．
\end{remark}


\subsection{局所コホモロジーと第二Bass空間}
標準のsection/local-cohomology比較を用いると，$\mathfrak m=A_{\geq1}$について，次数付きベクトル空間として
\begin{equation}\label{eq:localcohom}
 H^0_{\mathfrak m}(A)=H^1_{\mathfrak m}(A)=0,\quad
 H^2_{\mathfrak m}(A)\simeq E,\quad
 H^3_{\mathfrak m}(A)\simeq D(A)(3)
\end{equation}
となり，4次以上は零である．ここで$M(\ell)_n=M_{n+\ell}$，$D(M)_n=(M_{-n})^*$とする．
最後の式の成分は$(A_{-n-3})^*$であり，シフトは$+3$である．
$H^2_{\mathfrak m}(A)\simeq E$は命題\ref{prop:torsion}の自然な右$A$作用とも整合する．
$H^3_{\mathfrak m}(A)$については，この記述では双加群構造の同定は主張しない．
低次の二つの消滅は軌道環のsaturationと負次数Hom消滅，
残りは\eqref{eq:duality}とホモロジー次元から得られる．


\begin{proposition}[一つの次数付きBass空間による有限性判定]\label{prop:basscriterion}
$k=A/A_{\geq1}=\C$とし，下線付き$\underline{\Ext}$，$\underline{\Hom}$は
全内部次数の直和を表すものとする．自然な次数付きベクトル空間の同型
\begin{equation}\label{eq:basssocle}
 \underline{\Ext}^{2}_A(k,A)
 \simeq\underline{\Hom}_A(k,E)=\operatorname{soc}_A(E)
\end{equation}
がある．右辺は次数付きsocleである．従って次は同値である：
\[
 \dim_\C E<\infty
 \quad\Longleftrightarrow\quad
 \dim_\C\underline{\Ext}^{2}_A(k,A)<\infty.
\]
右辺は各内部次数の有限性ではなく，全内部次数を合わせた有限性である．
\end{proposition}
\begin{proof}
右noetherian connected graded環では，torsion関手$\Gamma_{\mathfrak m}$は
次数付きinjectiveを次数付きinjectiveへ送る．この標準的安定性と導来torsionの随伴性は，
例えば$\Z$-algebraへの一般化\cite[Corollary 3.16, Proposition 3.17]{Mizuno}にも明記されている．
$A$の次数付きinjective分解$I^\bullet$について
$\underline{\Hom}_A(k,\Gamma_{\mathfrak m}I^\bullet)
 =\underline{\Hom}_A(k,I^\bullet)$なので，
\[
 \mathbf R\underline{\Hom}_A(k,\mathbf R\Gamma_{\mathfrak m}A)
 \simeq\mathbf R\underline{\Hom}_A(k,A).
\]
従って第一象限のhyper-Extスペクトル系列
\[
 {}_2E^{p,q}=\underline{\Ext}^{p}_A(k,H^q_{\mathfrak m}(A))
 \ \Longrightarrow\ \underline{\Ext}^{p+q}_A(k,A)
\]
を得る．式\eqref{eq:localcohom}により，総次数2の唯一の項は
${}_2E^{0,2}=\underline{\Hom}_A(k,E)$である．
この項への高次微分は負の$p$から来るため零であり，出る高次微分は$q\leq1$へ向かうため零である．
これで式\eqref{eq:basssocle}を得る．

$E$が有限次元ならsocleも有限次元である．逆に$W=\operatorname{soc}_A(E)$を有限次元とする．
命題\ref{prop:torsion}により，$E$の任意の非零次数付き部分加群は$W$と非自明に交わる．
実際，その斉次元が生成する有限次元右加群の最高非零次数成分はsocleに入る．
特に$W=0$なら$E=0$である．

$W\ne0$のとき，斉次基底の次数を$r_1,\ldots,r_t$と置き，
$I=\bigoplus_{j=1}^tD(A)(-r_j)$とする．$D(A)$には左正則加群の双対として右作用を入れる．
任意の次数付き右加群$M$について
\[
 \Hom_{\operatorname{Gr}A}(M,D(A)(-r))\simeq (M_r)^*
\]
なので$I$は次数付きinjectiveであり，$\operatorname{soc}I\simeq W$である．
従って$W\hookrightarrow I$は$E\to I$へ延長できる．
その核が非零なら前段により$W$との交わりが非零となり，矛盾する．
よって$E\hookrightarrow I$であり，$E_n=0$ for $n>\max_jr_j$を得る．
Serre対称性\eqref{eq:duality}で左側も有界になり，各次数の有限次元性から$\dim_\C E<\infty$が従う．
\end{proof}

\begin{remark}[双対化複体に関する仮定の区別]
Serre関手の存在を，座標環がbalanced dualizing complexを持つという主張に置き換えてはならない．
Yekutieli--Zhang\cite[Theorem 2.3(2)]{YZ}の逆向きで得られるのは最上位コホモロジーの双対の有限生成であり，
同定理(1b)の下位コホモロジーに関する有限性はfull $\chi$を仮定する．
また，Mizuno\cite[Theorem 4.15]{Mizuno}の$\Z$-algebraに対する存在定理も，
両側の$\chi$，有限局所コホモロジー次元，左右のderived torsionの整合性を条件にする．
同文献Theorem 5.8は双対化複体からSerre関手を構成する向きであり，ここで必要な逆向きを与えない．
従って命題\ref{prop:basscriterion}の有限性自体は，依然として基礎公理からは未証明である．
\end{remark}

\subsection{normal切断と実局所コホモロジーの左右作用}\label{subsec:normal-vanishing}

\begin{lemma}[任意の非零切断の左作用]\label{lem:left-injective}
任意の非零斉次元$a\in A_d$，$d>0$について，ある$N(a)$が存在し，
\[
 a\cdot-:E_n\longrightarrow E_{n+d}
\]
は$n\geq N(a)$で単射となる．特に$e_{n+d}\geq e_n$ for $n\gg0$である．
この主張にnormal性は不要である．
\end{lemma}
\begin{proof}
$a:\OO\to\OO(d)$は単射なので，その余核を$C_a$として
\[
 0\longrightarrow\OO\xrightarrow{a}\OO(d)\longrightarrow C_a\longrightarrow0
\]
を得る．固定した全射$\OO(d)\twoheadrightarrow C_a$にAZ2を適用すると，十分大きい$n$について
\[
 \Hom(\OO(-n),\OO(d))\longrightarrow\Hom(\OO(-n),C_a)
\]
が全射になる．同じ短完全列に$\Hom(\OO(-n),-)$を適用した長完全列により
\[
 \Ext^1(\OO(-n),\OO)\longrightarrow\Ext^1(\OO(-n),\OO(d))
\]
が単射である．$s^n$で移すとこれは，先に定義した左作用
$\xi\mapsto (s^na)[1]\circ\xi$に一致する．
\end{proof}


\begin{lemma}[normal元に沿う左局所冪零性]\label{lem:normal-local}
$R$をlocally finite connected $\Z_{\geq0}$-graded右noetherian環，
$g\in R_d$を正次数の正則normal元とし，
$ag=g\sigma(a)$と置く．$U=H^i_{\mathfrak m}(R)$には，右正則加群$R_R$の自己射に
由来する左$R$作用と，本来の右$R$作用がある．
各$r\geq1$について次数0のベクトル空間同型$\Phi_r:U_{\sigma^r}\to U$が存在し，
\begin{equation}\label{eq:normal-local-action}
 g^r u=\Phi_r(u g^r)\qquad(u\in U)
\end{equation}
となる．特に$U$が右torsionなら，$g$の左作用も各元について冪零である．
\end{lemma}
\begin{proof}
任意の右$R$加群$M$について，右乗法は自然な次数付き$R$線形写像
\[
 \rho_{g^r}:M\longrightarrow M_{\sigma^r}(rd),\qquad m\longmapsto mg^r
\]
である．$\sigma$は次数を保つので，twistはtorsion関手と可換し，exactでinjectiveを保つ．
従ってinjective分解に$\rho_{g^r}$を項ごとに適用すると，導来torsion上の誘導写像は
通常の右乗法$u\mapsto ug^r$となる．

$M=R$のとき，$\sigma(g)=g$により
\[
 \psi_r:R_{\sigma^r}\xrightarrow{\sim}R,
 \qquad x\longmapsto\sigma^{-r}(x)
\]
は次数付き右$R$加群の同型であり，
$\psi_r(ag^r)=g^ra$である．従って左乗法$\lambda_{g^r}$は
$\psi_r\circ\rho_{g^r}$である．
$H^i_{\mathfrak m}$を適用し，$\Phi_r=H^i_{\mathfrak m}(\psi_r)$と置けば
式\eqref{eq:normal-local-action}を得る．
右torsion元$u$は十分大きい$r$で$ug^r=0$となるので，$g^ru=0$も従う．
\end{proof}


\begin{proposition}[一つのnormal切断からのSerre消滅]\label{prop:normal-vanishing}
$A$または反標準環$B=A^{(3)}$に正次数の非零斉次normal元が存在すると仮定する．
このとき$E=\bigoplus_{n\in\Z}E_n$は有限次元である．
従って任意の$F\in\Cat$について
\[
 \Ext^i(\OO,s^nF)=0\qquad(i>0,\ n\gg0)
\]
となり，$A$はfull right $\chi$を満たす．
normal元の次数，商環のdomain性・滑らかさ・楕円性について追加条件はない．
\end{proposition}
\begin{proof}
まず$g\in A_d$とする．section/local-cohomology比較
\[
 H^2_{\mathfrak m}(M)_n\simeq
 \Ext^1_{\qgr A}(\pi A,\pi M(n))
\]
は有限生成次数付き右加群$M$の射に自然である．
これを$M=A$および右$A$線形な左乗法$A(-d)\xrightarrow{g\cdot}A$に適用すると，
$H^2_{\mathfrak m}(A)\simeq E$は左$g$作用についても整合する．
すなわち局所コホモロジーの左作用は，補題\ref{lem:left-injective}の
postcompositionによる左作用と同じである．

$E$の各元は右torsionなので，補題\ref{lem:normal-local}より左$g$作用は局所冪零である．
一方，補題\ref{lem:left-injective}よりこの作用は$n\geq N(g)$で単射である．
$u\in E_n$，$n\geq N(g)$を取ると，ある$r$で$g^ru=0$であり，
\[
 E_n\xrightarrow{g}E_{n+d}\xrightarrow{g}\cdots\xrightarrow{g}E_{n+rd}
\]
の全ての矢が単射だから$u=0$となる．従って$E_n=0$ for $n\gg0$であり，
Serre対称性から$E$は有限次元である．

$g\in B=A^{(3)}$の場合は，最初に$\tau$による軌道を用い，同じ議論で
$\Ext^1(\OO,\tau^m\OO)=0$ for $m\gg0$を得る．
命題\ref{prop:serrevanishing}の反標準剰余類の消滅判定によって結論を得る．
\end{proof}
\begin{remark}[形式モデルとの差と反例候補への制約]
式\eqref{eq:normal-local-action}は，任意の抽象的双加群に成立する式ではない．
環$R$自体に対する導来torsionと，twistとの自然性から得た実現条件である．
特に中心元では$\sigma=\id$となり，実際の局所コホモロジー上で左右の作用が一致する．
既存の形式Yoneda代数は，正側の左正則作用と右augmentation作用を組み合わせているので，
多項式環に対するこの条件を満たさない．
従ってこの追加条件は，形式モデルに対して既に分かっていた実現不能性を具体的な作用の式で検出する．
また$E$が無限次元となる真の反例があれば，$A$と反標準環$B$のどちらにも
正次数の非零斉次normal元が存在してはならない．
\end{remark}


\begin{remark}[反標準環への適用と追加条件の限界]
$A^{(3)}$内のnormal性を$A$内のnormal性と同一視していない．
反標準環自身に議論を適用するので，その中心元でも十分である．
実際，三乗根を仮定しない元の六公理だけでも，反標準環のnormal元から
全対象の$\tau$方向のSerre消滅が従う．この場合のSerre対称性は
$\Ext^1(\OO,\tau^m\OO)^*\simeq\Ext^1(\OO,\tau^{-m-1}\OO)$である．
一方，元の公理からnormal元の存在を得たわけではなく，
ここでの$E$有限性だけでは$E_{-1}=0$や$\PP=0$までは結論しない．
研究ノートにはright $\chi_1$とRees補題による独立の代数的証明も記載した．
\end{remark}

\subsection{Normal切断上の構造対象の双対性}

\begin{lemma}[Normal切断の構造対象のExt]\label{lem:normal-boundary-ext}
$A$をlocally finite connected graded右noetherian代数，
$g\in A_e$を正次数の両側regularなnormal元，$B=A/(g)$とし，
$i_*:\qgr B\longrightarrow\qgr A$を忘却functorとする．
$D=\pi_BB$と置く．任意の$F\in\qgr B$と$j\geq0$について
\[
 \Ext^j_{\qgr A}(i_*D,i_*F)
 \simeq \Ext^j_{\qgr B}(D,F)
       \oplus \Ext^{j-1}_{\qgr B}(D,F(e)),
\]
ただし負のExtは零とする．この分解は$g$による自由表示を固定すれば$F$に自然である．
\end{lemma}
\begin{proof}
$-\otimes_A B$と忘却の随伴をgraded加群からtailsへ移す．
その導来functorがtailsへ降りることを確認する．
$B$は左$A$加群として長さ1の有限自由分解を持ち，
その両側加群表示で現れるnormal automorphismは次数を保つ．
従ってtorsion graded $A$加群$M$に対する
$\Tor_i^A(M,B)$もtorsionであり，$i>1$では零である．
よって導来tensorはtorsion cohomologyを持つ複体を同じ型の複体へ送り，
$D(\Gr A)$から$D(\QGr A)$への導来局所化を通じて
$Li^*:D(\QGr A)\to D(\QGr B)$を誘導する．
ここで導来局所化は，exactな商functor $\pi$とその右随伴$\omega$について
$\pi R\omega\simeq\mathrm{id}$となることによる．
実際，$\omega$はinjectiveを保ち，$\pi\omega=\mathrm{id}$である．
従って局所化の核はtorsion cohomologyを持つ複体であり，
上のtensorと忘却の双方がこの核を保つので，graded加群での随伴のunitとcounitも降りる．
忘却はexactでtorsionを保つため，導来随伴
\[
 \RHom_{\QGr A}(X,i_*F)
 \simeq \RHom_{\QGr B}(Li^*X,F)
\]
が得られる．有限生成対象のExtはlocally noetherianな大圏と一致する．

右$A$加群の自由分解
\[
 0\to A(-e)\xrightarrow{g\cdot}A\to B\to0
\]
に$-\otimes_A B$を施すと微分が零となるから，
\[
 Li^*i_*D\simeq D\oplus D(-e)[1].
\]
これを導来随伴へ代入する．normal automorphismは両側加群構造には現れるが，
ここで計算する右$B$加群$B$の構造対象はそのtwistと同型なので，表示されたshiftは$(-e)$である．
\end{proof}

\begin{proposition}[Normal境界の構造対象に対する曲線双対性]\label{prop:normal-boundary-duality}
上の設定で$\qgr A$がExt-finite，homological dimension $2$を持ち，
その有界導来圏のSerre functorが$(-\ell)[2]$であるとする．
このとき任意の$F\in\qgr B$について
\[
 \Ext^j_{\qgr B}(D,F)=0\quad(j\geq2),\qquad
 \Ext^1_{\qgr B}(D,F)
 \simeq\Hom_{\qgr B}(F,D(e-\ell))^*.
\]
これは構造対象を第一変数に取る双対性であり，
$\qgr B$の有限homological dimensionや$D^b(\qgr B)$のSerre functorの存在を主張しない．
\end{proposition}
\begin{proof}
前補題を$j\geq3$で用いれば，左辺の消滅から
$\Ext^{j-1}_{\qgr B}(D,F(e))=0$が従う．
$j=2$では，すでに消えた項を除いて
\[
 \Ext^1_{\qgr B}(D,F(e))
 \simeq \Ext^2_{\qgr A}(i_*D,i_*F)
 \simeq \Hom_{\qgr A}(i_*F,i_*D(-\ell))^*
 \simeq \Hom_{\qgr B}(F,D(-\ell))^*.
\]
最後に$F$を$F(-e)$で置き換える．
\end{proof}

\subsection{Normal切断の高次Ext周期性}
\label{subsec:normal-high-ext}

normal元のtwisted matrix factorizationはCassidyら\cite[Theorem 4.2]{CCKM}で扱われている．
同定理はambient代数の有限global dimensionを仮定する．
以下では商圏の有限homological dimensionを仮定するExtの形を，導来随伴から直接証明する．
構造対象以外の一般の対象では，二つのTor群から直和分解を結論しないことが要点である．

\begin{lemma}[Normal切断の高次Ext周期性]
\label{lem:v6-normal-ext-periodicity}
$A$をconnected locally finite graded右noetherian代数，
$g\in A_e$を両側regularなnormal元，$B=A/(g)$とする．
$ga=\nu(a)g$でnormal automorphismを定め，
$\Phi(F)=F_\nu$で$\qgr B$上のtwist functorを表す．
$T=(1)$を次数シフトとし，$\operatorname{hdim}\qgr A\leq d<\infty$を仮定する．
このとき任意の$F,G\in\qgr B$と$n\geq d$に対して
\[
 \Ext^{n+2}_{\qgr B}(F,G)
 \simeq \Ext^n_{\qgr B}(\Phi T^{-e}F,G)
\]
である．また$D=\pi_B B$について$\Phi D\simeq D$であり，$\Phi T=T\Phi$である．
\end{lemma}
\begin{proof}
$I=(g)$と置く．$a\mapsto ag$による両側加群同型は
\[
 I\simeq {}_1A_\nu(-e),\qquad
 I/I^2\simeq {}_1B_\nu(-e)
\]
を与える．従って$B$加群$F$について
\[
 \Tor^A_0(F,B)=F,\quad
 \Tor^A_1(F,B)=F_\nu(-e),\quad
 \Tor^A_j(F,B)=0\ (j\geq2).
\]
この計算はtorsion加群を保ち，forgetful functor $i_*$とderived tensorの随伴は
tailsへ降りる．具体的には，$B$の左$A$-flat dimensionが$1$で，
normal automorphismが次数を保つことからderived tensorはtorsion cohomologyを保つ．
exact quotientとその導来右随伴による局所化に随伴のunitとcounitを降ろせば，
$Li^*:D(\operatorname{QGr}A)\rightleftarrows D(\operatorname{QGr}B):i_*$を得る．
従って$Li^*i_*F$の標準切断はtriangle
\[
 \Phi T^{-e}F[1]\longrightarrow Li^*i_*F\longrightarrow F
 \longrightarrow\Phi T^{-e}F[2]
\]
を与える．ここでは一般の$F$についてこのtriangleが分裂するとは主張していない．
両側$A$加群の長さ$1$の表示をtensorして零微分に見えることは，
$D(\operatorname{QGr}B)$での分裂を意味しない．

このtriangleに$\RHom_B(-,G)$を適用し，導来随伴を用いると
\[
 \Ext_A^{n+1}(i_*F,i_*G)\longrightarrow
 \Ext_B^n(\Phi T^{-e}F,G)\longrightarrow
 \Ext_B^{n+2}(F,G)\longrightarrow
 \Ext_A^{n+2}(i_*F,i_*G)
\]
を得る．$n\geq d$なら両端が零であり，所要の同型が従う．
$B_\nu\simeq B$は右$B$加群として同型なので$\Phi D\simeq D$であり，
次数を保つtwistは$T$と可換である．
\end{proof}

\begin{proposition}[高次の偶数Extが対角的な単純対象への障害]
\label{prop:v6-diagonal-simple-obstruction}
前補題の設定で，$(D,T^e)$がAZの生成条件を満たすとする．
$\Sigma$を単純対象の同型類の集合で，$T$，$\Phi$およびそれらの逆functorで保たれるものとする．
十分大きい全ての$m$について
\[
 \Ext_B^{2m}(S,S)\ne0\quad(S\in\Sigma),\qquad
 \Ext_B^{2m}(S,S')=0\quad(S\ne S',\ S,S'\in\Sigma)
\]
と仮定する．このとき全ての$S\in\Sigma$について$\Hom_B(D,S)\ne0$である．
\end{proposition}
\begin{proof}
高次Ext周期性で$F=G=S$とすると
\[
 0\ne\Ext_B^{2m+2}(S,S)
 \simeq\Ext_B^{2m}(\Phi T^{-e}S,S).
\]
対角性より$\Phi T^{-e}S\simeq S$となり，$\Phi$と$T$の可換性から
$T^eS\simeq\Phi S$である．従って
\[
 \Hom_B(D,T^eS)\simeq\Hom_B(D,\Phi S)
 \simeq\Hom_B(D,S).
\]
$\Sigma$は$T$で保たれるので，この等式を反復すれば
$\Hom_B(D,T^{em}S)$の非零性は$m\in\Z$によらない．
一方AZ生成条件から，ある$n\geq0$に対して
$\Hom_B(T^{-en}D,S)=\Hom_B(D,T^{en}S)\ne0$である．
従って$\Hom_B(D,S)\ne0$となる．
\end{proof}

\subsection{飽和したnormal切断商からの全次数消滅}

\begin{proposition}[飽和境界による$E$の消滅]\label{prop:saturated-boundary}
六公理と固定したexactな三乗根の下で，$0\ne g\in A_e$は正次数のnormal元とする．
$B=A/(g)$が全次数で飽和している，すなわち標準写像
\begin{equation}\label{eq:boundary-saturation}
 B_n\longrightarrow\Hom_{\qgr B}(\pi_BB,\pi_BB(n))
\end{equation}
が全$n\in\Z$について同型であると仮定する．
このとき$E_n=\Ext^1_{\qgr A}(\OO,\OO(n))=0$は全$n\in\Z$で成り立つ．
ここで$B_n=0$ for $n<0$と約束する．
\end{proposition}
\begin{proof}
$D=\pi_AB$と置く．まず，商写像$\OO\to D$によるprecompositionは
\begin{equation}\label{eq:boundary-hom}
 \Hom_{\qgr A}(D,D(n))\xrightarrow{\ \sim\ }
 \Hom_{\qgr A}(\OO,D(n))
\end{equation}
を与える．単射性は商写像がepimorphismであることによる．
全射性を確認するため，$ga=\sigma(a)g$と書く．
$f:\OO\to D(n)$をtailsの射
$u:A_{\geq q}\to B(n)/T$で表す．ここで$T$は有限次元部分加群であり，
標的は$g$で消える．$a$の次数が十分大きければ
\[
 u(ga)=u(\sigma(a)g)=u(\sigma(a))g=0
\]
である．従って$f$は$g\OO(-e)$を殺し，一意に$D$を経由する．

忘却による$\qgr B\to\qgr A$は完全忠実なので，
式\eqref{eq:boundary-saturation}と\eqref{eq:boundary-hom}から
\[
 \Hom_{\qgr A}(\OO,D(n))\simeq B_n
 \qquad(n\in\Z)
\]
を得る．この同型は標準のsections mapと整合する．
短完全列
\[
 0\longrightarrow\OO(n-e)\xrightarrow{\;g\cdot\;}\OO(n)
 \longrightarrow D(n)\longrightarrow0
\]
に$\Hom(\OO,-)$を適用する．標準の$A_n\to B_n$は全射だから，
接続準同型$B_n\to E_{n-e}$は零であり，
\begin{equation}\label{eq:E-injection}
 E_{n-e}\lhook\joinrel\longrightarrow E_n
 \qquad(n\in\Z)
\end{equation}
を得る．負次数では$\Hom(\OO,\OO(n))=A_n=0$という基礎補題を用いた．
命題\ref{prop:normal-vanishing}から$E$は有限次元である．
もし$E\ne0$なら，$E_m\ne0$となる最大次数$m$を取れるが，
式\eqref{eq:E-injection}は$E_m\hookrightarrow E_{m+e}=0$を与えて矛盾する．
\end{proof}

\begin{remark}[局所コホモロジーによる同じ結論とconicへの適用]
上の因子化の証明は，$B$を右$A$加群と見ても
$B\simeq\Gamma_{*,A}(\pi_AB)$となることを示している．
従って$H^0_{\mathfrak m_A}(B)=H^1_{\mathfrak m_A}(B)=0$であり，
$0\to A(-e)\xrightarrow{g\cdot}A\to B\to0$の局所コホモロジー長完全列も
\eqref{eq:E-injection}を与える．
この説明では$A$上と$B$上の局所コホモロジーを無条件に同一視していない．
左乗法$g$は右$A$-linearな写像$A(-e)\to A$なので，この単射自体にmodule twistは現れない．

$B$が$\mathbb P^1$上の次数2の線束のtwisted homogeneous coordinate ringならば，
$\qgr B\simeq\operatorname{coh}\mathbb P^1$の下で全$n$のsectionsは
$H^0(\mathbb P^1,\mathcal L_n)$である．これは$n\geq0$で$B_n$，$n<0$で零なので，
式\eqref{eq:boundary-saturation}が成立し，上の命題を適用できる．
この段階では$A$の二次関係性やglobal dimensionを仮定する必要はない．
\end{remark}

\begin{corollary}\label{cor:eventual-ext-monotone}
$e_n=\dim E_n$とすると，$e_{n+3}\geq e_n$ for $n\gg0$である．
もし$E$が無限次元なら$e_{3m}\geq1$ for all $m\gg0$となる．
\end{corollary}
\begin{proof}
$A_3\ne0$の切断に補題\ref{lem:left-injective}を適用する．
$E$が無限次元なら，命題\ref{prop:serrevanishing}により$E_{3m}$はeventually zeroでない．
従って正方向の単調性から後半も従う．
\end{proof}

\section{Euler特性による最初の障害の除去}\label{sec:euler}
\begin{proposition}[全整数の多項式性による判定]\label{prop:polynomial}
全整数$n$に対して$f(n)=\chi(\OO,\OO(n))$が次数2以下の一つの多項式であると仮定する．
このとき
\[
 f(n)=\frac{(n+1)(n+2)}2,\qquad E_{-1}=E_{-2}=E_3=0.
\]
\end{proposition}
\begin{proof}
Serre双対性より$f(n)=f(-n-3)$であり，$f(0)=1$である．従って
$f(n)=1+c\,n(n+3)$と書ける．$e=\dim E_{-1}$とすると
\eqref{eq:negative}と\eqref{eq:duality}から$f(-1)=-e$，すなわち$c=(e+1)/2$である．
一方
\[
 10+9e=f(3)=10-\dim E_3.
\]
両方の次元が非負なので$e=\dim E_3=0$である．$E_{-2}$も双対性によって消える．
\end{proof}

仮定は全整数にわたる一つの多項式であることであり，正方向でのeventual polynomialityだけでは足りない．
また，この命題は$E_1,E_2$の消滅や$\PP=0$を単独では導かない．
定理\ref{thm:finite}と組み合わせれば，全整数の多項式性と$L_3\in\TT$は認識の十分条件となる．
\subsection{軌道上のEuler形式のrankによる補強}
$W\subset K_0(\Cat)\otimes_\Z\Q$を$[\OO(i)]$，$i\in\Z$の張る部分空間とし，$T=s_*$と置く．
Euler形式の$W$への制限の左右radicalは一致する．実際，$T$不変性と
\begin{equation}\label{eq:numericalserre}
 \chi(x,y)=\chi(y,T^{-3}x)
\end{equation}
から，一方のradicalは他方を$T^3$で移したものとなり，各radicalは$T$で保たれる．
共通radicalを$R$，$V=W/R$とし，$r=\dim_\Q V$を\emph{軌道上のEuler rank}と呼ぶ．
全$K_0(\Cat)$の有限生成性やrankについての条件ではない．

\begin{theorem}[小さいEuler rankからの消滅]\label{thm:numericalrank}
基礎公理と三乗根に加えて
\begin{equation}\label{eq:rankhyp}
 E_{3m}=0\quad(m\gg0),\qquad r\leq4
\end{equation}
を仮定する．このとき$r=3$であり，全整数$n$に対して
\[
 \chi(\OO,\OO(n))=\binom{n+2}{2}.
\]
特に$E_{-1}=E_{-2}=0$であり，$(\OO,\OO(1),\OO(2))$は例外列である．
\end{theorem}
\begin{proof}
$V$上の非退化なEuler形式を$B$，$[\OO]$の像を$v$と書く．
$V$は$T$と$v$に関してcyclicであり，$T$は$B$のisometryである．
$q_n=B(v,T^nv)$と置く．この両側数列の最小漸化式の特性多項式は$T$の最小多項式に一致する．
実際，多項式$p$について$p(\mathrm{shift})q_n=0$が全$n$で成立すれば
\[
 B(T^{-n}v,p(T)v)=0\quad(n\in\Z).
\]
cyclicityと非退化性より$p(T)v=0$，従って$p(T)=0$である．逆向きは直ちに従う．
特に最小漸化式の次数は$r$である．

$U=T^3$は可逆なので，列$q_{3m}$は定数項が非零の両側線形漸化式を満たす．
仮定により十分大きな$m$では$q_{3m}=P(3m)$，$P(n)=(n+1)(n+2)/2$である．
この漸化式の作用素を$P(3m)$に適用した結果は$m$の多項式で，十分大きな整数で零になるため恒等的に零である．
差の列は同じ可逆な漸化式に従い，十分大きな$m$で零なので，後方帰納法によって
\begin{equation}\label{eq:canonicalrecurrence}
 q_{3m}=P(3m)\quad(m\in\Z)
\end{equation}
を得る．この段階では$r$の有限性のみを使っている．

複素数体に拡大してJordan標準形を考える．$P(3m)$が真に二次であるため，
$U$の最小多項式は$(z-1)^3$で割り切れる．従って$T$には$\lambda^3=1$となる固有値の，
サイズ3以上のJordan blockがある．$T$がisometryであるため，一般固有空間$V_\lambda$は
$V_{\lambda^{-1}}$と完全に対になり，他の一般固有空間とは直交する．
$\lambda\ne1$なら$\lambda\ne\lambda^{-1}$なので，それぞれの次元が3以上となって$r\geq6$である．
よって$r\leq4$の下では，固有値1にサイズ3以上のblockがある．

固有値が全て1なら$q_n$は$n$の多項式である．\eqref{eq:canonicalrecurrence}によって$q_n=P(n)$となり，
最小漸化式の次数を比較して$r=3$を得る．
残る可能性は，サイズ3の1-blockと，別固有値$\mu$の一次元空間である．
$\mu\notin\{1,-1\}$なら逆数固有値の相手が必要となって次元不足である．
$\mu=1$の別blockはcyclicityに反する．$\mu=-1$の場合，その一次元空間は他と直交し，
\eqref{eq:numericalserre}により$B(x,y)=-B(y,x)$となる．
一次元の交代形式は退化するため，$B$の非退化性に反する．
従って$q_n=P(n)$のみが残る．最後の消滅は命題\ref{prop:polynomial}による．
\end{proof}

この定理の二つの追加条件は基礎公理からは未証明である．
命題\ref{prop:serrevanishing}により，\eqref{eq:rankhyp}の最初の条件は
$E$の有限次元性および全対象のSerre消滅と同値であり，それらより弱い追加仮定ではない．
また，任意の有限Euler rankと最初の条件だけでも，証明中の可逆漸化式により
$E_{3m}=0$が全整数$m$に延長する．

rank上界は数値データだけを考えると改善できない．例えば$\varepsilon_n$を周期6の列
$(0,1,1,0,-1,-1)$とし，$e>0$に対して
\[
 q_n=P(n)+e\varepsilon_n
\]
と置くと，$q_n=q_{-n-3}$，$q_{3m}=P(3m)$，$q_{-1}=-e$を満たす．
最小漸化式は$(z-1)^3(z^2-z+1)$であり，5対象のEuler型行列
$(q_{j-i})_{0\leq i,j\leq4}$の行列式は$e^2$である．
$e=1$ではunimodularにもなる．これらは圏・積・ample性を構成した反例ではなく，
追加の幾何学的議論が必要であることを示す数値モデルである．
詳しい二種類のモデルと行列は研究ノートに記す．


\subsection{積写像による低次数の制約}
この小節では，基礎公理とexactな三乗根から得られる軌道環
$A=\bigoplus_{n\geq0}\Hom(\OO,\OO(n))$が
$\C$上の次数付きdomainであることを使う．
$h_n=\dim_\C A_n$，$e_n=\dim_\C E_n$と置く．
従って$h_0=1$，$h_3=10$，$h_6=28$である．
以下の積写像の制約は，Euler形式の有限rankや$E$の有限次元性を仮定しない．

\begin{lemma}[有限次元部分空間の積]\label{lem:productdimension}
代数閉体$k$上のdomain $D$と，その非零有限次元部分空間$U,V$に対して，
$UV$を積$uv$の線形包とすると
\[
 \dim_k UV\geq\dim_k U+\dim_k V-1
\]
が成立する．$D$の可換性は仮定しない．
\end{lemma}
\begin{proof}
$a=\dim U$，$b=\dim V$，$c=\dim UV$とし，積写像
$\mu:U\otimes_k V\to UV$の核を$K$とする．
$K=0$なら$c=ab\geq a+b-1$である．
$K\ne0$とする．$D$がdomainなので，$K$は非零分解可能テンソル$u\otimes v$を含まない．
従って射影空間$\mathbb P(U\otimes V)$の中で，
線形部分空間$\mathbb P(K)$とSegre多様体
$\mathbb P(U)\times\mathbb P(V)$は交わらない．
代数閉体上の射影多様体の次元定理により，両者の次元の和は周囲の次元より小さいから，
\[
 (ab-c-1)+(a+b-2)<ab-1.
\]
よって$c>a+b-2$であり，整数性から結論を得る．
\end{proof}

\begin{proposition}[低次数の無条件制約]\label{prop:lowdegreeproduct}
基礎公理とexactな三乗根の下で
\[
 h_1\leq4,\qquad h_2\leq10
\]
が成立する．$h_1>0$の場合には，さらに
\[
 2h_1-1\leq h_2\leq11-h_1.
\]
特に$h_1=4$なら
\[
 h_2=7,\qquad A_1A_1=A_2,\qquad
 A_1A_2=A_2A_1=A_3=A_1^3.
\]
ここで各積は対応する積部分空間の線形包を表す．
\end{proposition}
\begin{proof}
$h_1>0$なら，補題\ref{lem:productdimension}を二度適用して
\[
 3h_1-2\leq\dim A_1^3\leq h_3=10
\]
となるので$h_1\leq4$である．$h_1=0$の場合にもこの上界は成立する．
同様に$h_2>0$なら
$3h_2-2\leq\dim A_2^3\leq h_6=28$より$h_2\leq10$であり，
$h_2=0$の場合も自明である．

$h_1>0$のときは$A_1^2\ne0$なので$h_2>0$であり，
\[
 h_2\geq\dim A_1^2\geq2h_1-1,
 \qquad
 10=h_3\geq\dim A_1A_2\geq h_1+h_2-1
\]
が両側の不等式を与える．
$h_1=4$なら$7\leq h_2\leq7$となり，
上記の各次元不等式は等号である．従って$A_1^2=A_2$と$A_1A_2=A_3$を得る．
$A_2A_1$にも同じ補題を適用すれば$A_2A_1=A_3$であり，
結合性から$A_1^3=A_3$である．
\end{proof}

\begin{proposition}[Euler二次性の下で残る四候補]\label{prop:fourlowdegreecases}
さらに，全整数$n$について
\[
 \chi(\OO,\OO(n))=P(n)=\frac{(n+1)(n+2)}2
\]
を仮定する．このとき
\[
 (e_1,e_2)\in\{(0,0),(0,1),(0,2),(1,1)\}.
\]
対応する$(h_1,h_2)$は，順に$(3,6),(3,7),(3,8),(4,7)$である．
\end{proposition}
\begin{proof}
$n\geq0$ではSerre双対性と負次数Hom消滅により
$\Ext^2(\OO,\OO(n))=0$なので，$h_n=P(n)+e_n$である．
特に$h_1=3+e_1>0$，$h_2=6+e_2$となる．
命題\ref{prop:lowdegreeproduct}から
\[
 e_1\leq1,\qquad e_1+e_2\leq2,\qquad e_2\geq2e_1-1
\]
を得る．非負整数解を分ければ，$e_1=0$のとき$e_2=0,1,2$，
$e_1=1$のとき$e_2=1$だけが残る．
\end{proof}

\begin{proposition}[位数6のrank5数値モデルの縮約]\label{prop:rankfiveproduct}
正整数$e$と周期6の列$\varepsilon_n=(0,1,1,0,-1,-1)$に対して，
実際の圏のEuler列が全整数$n$で
\[
 q_n=P(n)+e\varepsilon_n
\]
を満たすと仮定する．このとき必ず
\[
 e=1,\qquad E_1=E_2=0,\qquad(h_1,h_2,h_3)=(4,7,10)
\]
であり，$A_1^2=A_2$，$A_1^3=A_3$となる．
\end{proposition}
\begin{proof}
正次数の$\Ext^2$は零なので
$h_1=3+e+e_1$，$h_2=6+e+e_2$である．
$h_1\leq4$，$e\geq1$，$e_1\geq0$より$e=1$かつ$e_1=0$が従い，$h_1=4$となる．
命題\ref{prop:lowdegreeproduct}から$h_2=7$なので$e_2=0$である．
積部分空間の等式も同命題から従う．
\end{proof}

位数3の補正$q_n=P(n)-e\mathbf1_{3\nmid n}$に対しても，同じ方法は具体的な制約を与える．
$x=e_1$，$y=e_2$と置けば
\[
 h_1=3-e+x,\qquad h_2=6-e+y.
\]
$h_1>0$の場合には
\[
 0\leq x\leq e+1,\qquad
 0\leq y\leq e+4,\qquad
 y\geq2x-e-1,\qquad x+y\leq2e+2.
\]
これらには$3-e+x>0$という場合分けの条件も課される．
$h_1=0$の場合には$x=e-3$，$0\leq6-e+y\leq10$となる．
特にこの計算だけでは位数3の補正の係数$e$に一様な上界は得られない．

命題\ref{prop:rankfiveproduct}は数値モデル全体を排除するものではない．
残る$(h_1,h_2,h_3)=(4,7,10)$を基礎公理と三乗根から除外することは未解決である．
また，次数3までの積による生成から$A$全体の一次生成性を推論してはならない．



\subsection{次数6からの無条件制約}
\begin{proposition}\label{prop:h6bounds}
基礎公理と三乗根の下で$h_1=4$とする．このとき
\[
 13\leq h_4\leq22,\qquad h_4+3\leq h_5\leq25.
\]
さらに全整数Euler二次性$q_n=P(n)$があれば
\[
 0\leq e_4\leq7,\qquad0\leq e_5\leq4,
 \qquad e_5\geq e_4-3.
\]
ここで$e_i=\dim E_i$であり，この計算には$E$の有限次元性を用いない．
\end{proposition}
\begin{proof}
domainの積に対する次元不等式を使うと
\[
 h_4\geq h_3+h_1-1=13,\qquad
 h_5\geq h_4+h_1-1=h_4+3.
\]
他方$h_6=28$と$h_2=7,h_1=4$から
\[
 h_4+7-1\leq28,\qquad h_5+4-1\leq28
\]
を得る．$q=P$なら正次数の$\Ext^2$消滅から
$h_4=15+e_4,h_5=21+e_5$なので，代入すれば後半が従う．
\end{proof}

\subsection{四点の最小和集合}
以下で使う最小増大の逆向きは，四点の場合に限って初等的に証明できる．
非可換division ring内の一般の部分空間に対する逆定理を引用するものではない．

\begin{lemma}\label{lem:fourpoints}
格子$L$の四元部分集合$U$が$|U+U|=7$を満たせば，
ある$u,a\in L$，$a\ne0$について
\[
 U=\{u,u+a,u+2a,u+3a\}
\]
と書ける．
\end{lemma}
\begin{proof}
$U$の元に異なる値を取る実線形汎関数$\ell$を選び，その値で並べる．
最小の元を引いて$U=\{0,a,b,c\}$，
$0<\ell(a)<\ell(b)<\ell(c)$としてよい．
次の七つの和は$\ell$の値が厳密に増加するため互いに異なる：
\[
 0,\ a,\ 2a,\ a+b,\ 2b,\ b+c,\ 2c.
\]
仮定からこれらが$U+U$の全てである．
$b=0+b$の値は$\ell(a)$と$\ell(a+b)$の間にあるので，この列の中では$b=2a$しかあり得ない．
次に$c=0+c$は$\ell(2a)$と$\ell(2a+c)$の間にあるので，$c=3a$または$c=4a$である．
後者なら$a+c=5a$は上記の列
$\{0,a,2a,3a,4a,6a,8a\}$に属さず矛盾する．従って$c=3a$である．
\end{proof}

\subsection{normalな半群を台とする単項式型代数の排除}
ここでいうnormal affine semigroupは，有限生成格子$L$内の有限生成部分半群$S$で，
その生成する群が$L$，生成する実凸錐が$C$であり，$S=C\cap L$となるものをいう．
次数関数$\deg:L\to\Z$は$S\setminus\{0\}$上で正とする．
特に次数切断は有限である．

通常の半群環$\C[S]$に加えて，基底$z^u$，$u\in S$を持ち，
\[
 z^u z^v=\alpha(u,v)z^{u+v},\qquad\alpha(u,v)\in\C^*,
\]
という結合的な積を備えた次数付き代数も考える．
すなわち$\alpha$は単位と整合するmultiplicative $2$-cocycleである．
以下の証明は指数の和と各次数の次元だけを使うため，このねじれた場合も含む．

\begin{theorem}[normal半群型の最小増大の排除]\label{thm:toricexclusion}
上記の単項式型代数$B$について$\operatorname{rank}L=3$とする．
次の条件は両立しない：
\[
 \dim B_1=4,\qquad\dim B_2=7,\qquad\dim B_3=10,
 \qquad\dim B_{3m}=P(3m)\quad(m\gg0).
\]
従って，基礎公理と三乗根を満たし$h_1=4$となる軌道環は，
このようなnormal affine semigroupを台とする単項式型代数ではあり得ない．
\end{theorem}
\begin{proof}
$U=\{u\in S:\deg u=1\}$とする．$|U|=4$であり，
$U+U\subset S_2$なので$|U+U|\leq7$である．
補題\ref{lem:fourpoints}の証明中の七つの和は任意の四点に対して相異なるので，
$|U+U|\geq7$である．従って等号となり，同補題から
\[
 U=\{u,u+a,u+2a,u+3a\}
\]
となる．$\deg a=0$であり，$a\ne0$である．
その三重和は相異なる十点$3u+ia$，$0\leq i\leq9$からなる．
$|S_3|=10$だから，$S_3$はこの一直線上の十点そのものである．

一方，$\Delta=(C\cap\{\deg=3\})-3u$とする．
これは$L_0:=\ker(\deg:L\to\Z)$を格子とする二次元の有理凸多角形である．
$S=C\cap L$というnormalityから，そのEhrhart関数は
\[
 L_\Delta(m)=|m\Delta\cap L_0|=\dim B_{3m}
\]
である．$L_\Delta(m)$は準多項式であり，十分大きな全整数$m$で$P(3m)$に一致する．
準多項式の各剰余類に対応する多項式が無限個の整数で一致するため，
全整数へ延長した準多項式も$L_\Delta(m)=P(3m)$である．
Ehrhart--Macdonald reciprocityより
\[
 |\operatorname{int}\Delta\cap L_0|=L_\Delta(-1)=P(-3)=1.
\]
従って$\Delta$の十個の格子点のうち，一つが内部点で，九つが境界点である．

ところが，これら十点は全て同じ直線上にある．
内部点を含む直線と二次元凸多角形の境界との交わりは，切断線分の二端点だけである．
境界上に相異なる九点を持つことはできず，矛盾する．
\end{proof}

ここで引用した有理多角形の準多項式性とreciprocityは既知の結果であり，
Stanley\cite[\S6, Proposition 6.1, p.~208]{StanleyReciprocity}と
Beck--Develin\cite[\S4, Theorem 5]{BeckDevelin}で一次資料を確認した．
四点の和集合補題および本定理への適用には上記の直接証明を与えたが，新規性は未確認である．

\begin{corollary}\label{cor:toricdegeneration}
$h_1=4$の軌道環は，同じ次数付きHilbert関数を持つ上記の単項式型代数をspecial fiberとする
次数付きflat degenerationを持たない．
例えば$\C[t]$上の次数付き代数で，各次数が有限自由$\C[t]$加群となり，
generic fiberが$A\otimes_\C\C(t)$，special fiberが定理\ref{thm:toricexclusion}の$B$
であるようなdegenerationは存在しない．
\end{corollary}
\begin{proof}
各次数の有限自由性からspecial fiberとgeneric fiberのHilbert関数は等しい．
定理\ref{thm:toricexclusion}を適用すればよい．
\end{proof}

\emph{この系は任意のvaluationの値半群が有限生成またはnormalとなることを主張しない．}
一般の非可換軌道環の積を，normal半群の格子点の和へ置換する根拠はまだない．


\section{次数1の整なnormal切断による認識}
\label{sec:normal-boundary}

この節では，元の六公理と固定したexactな三乗根$s$を仮定する．
$A=\bigoplus_{n\geq0}\Hom(\OO,s^n\OO)$，$h_n=\dim_\C A_n$とする．
既に証明した再構成により，$A$はconnected graded right noetherian domainであり，
$(\Cat,\OO,s)\simeq(\qgr A,\pi A,(1))$である．特に
\[
 h_{3m}=\binom{3m+2}{2}\quad(m\geq0),\qquad
 \SSer=(-3)[2]\quad\text{on }\Db(\qgr A).
\]
以下の認識定理は，全整数Euler多項式性，$E$の有限性，
$\dim A_1$の下界，評価射の全射性を追加仮定としない．
切断商がdomainであることは実質的な仮定である．

\subsection{一次normal元と剰余次数の障害}

\begin{lemma}[Serreシフトと切断商の次数support]\label{lem:residue-obstruction}
$A$をlocally finite connected graded右noetherian代数，
$g\in A_e$を正次数の両側regularなnormal元とし，$B=A/(g)$は無限次元とする．
$\qgr A$がExt-finiteで，その有界導来圏のSerre functorが$(-\ell)[2]$であると仮定する．
ある$d\geq1$について$B_n=0$ whenever $d\nmid n$ならば
\begin{equation}\label{eq:residue-divisibility}
 d\mid(e-\ell)
\end{equation}
である．ここでは$B$のdomain性を仮定しない．
\end{lemma}
\begin{proof}
まずSerre双対性そのものから
$\Ext^j_{\qgr A}(X,Y)^*\simeq\Ext^{2-j}_{\qgr A}(Y,X(-\ell))=0$
for $j>2$となるので，命題\ref{prop:normal-boundary-duality}を適用できる．
$D=\pi_BB\ne0$と置き，$F=D(e-\ell)$に同命題を適用すると
\[
 \Ext^1_{\qgr B}(D,D(e-\ell))\simeq\End_{\qgr B}(D)^*\ne0
\]
である．しかし$B$のsupportの仮定により，$\qgr B$は次数の剰余類modulo $d$で直和分解される．
$d\nmid(e-\ell)$なら二対象は異なる直和成分に属し，その間の全Ext群が零となるため矛盾する．
\end{proof}

\subsection{AZ生成条件によるnormal商の剰余障害}

\begin{lemma}[AZ1と切断商の次数support]\label{lem:normal-quotient-AZ-residue}
$A$をlocally finite connected graded右noetherian代数とし，
$g\in A_e$を正次数の両側regularなnormal元とする．
$B=A/(g)$は無限次元とし，$\OO=\pi_AA$，$D=\pi_AB$と置く．
ある整数$r>0$について$(\OO,(r))$がAZ1を満たすと仮定する．
すなわち任意の$F\in\qgr A$は，$\OO(-rm_i)$，$m_i\geq0$の有限直和の商であるとする．
もし$d\geq1$について$B_n=0$ whenever $d\nmid n$ならば，
\[
 \gcd(d,r)=1
\]
である．この主張に$B$の飽和性，domain性，Serre双対性は不要である．
\end{lemma}
\begin{proof}
まず全$n\in\Z$について，商写像$\OO\twoheadrightarrow D$は同型
\begin{equation}\label{eq:normal-quotient-hom-factor}
 \Hom_{\qgr A}(D,D(n))\xrightarrow{\sim}
 \Hom_{\qgr A}(\OO,D(n))
\end{equation}
を誘導する．単射性はepimorphismによるprecompositionから従う．
全射性のため$ga=\sigma(a)g$と書き，任意の射$\OO\to D(n)$を
tailsの射$u:A_{\geq q}\to B(n)/T$で表す．$T$は有限次元部分加群である．
標的は$g$で消え，$\sigma$は次数を保つので，十分高い次数の$a$について
\[
 u(ga)=u(\sigma(a)g)=u(\sigma(a))g=0.
\]
従って射は$g\OO(-e)$を殺し，$D$を経由する．
この議論にsections map $B\to\Gamma_*D$の同型性は使っていない．

忘却functor $\qgr B\to\qgr A$は完全忠実である．
実際，$g$で消える加群の部分加群・商加群も$g$で消えるので，
有限次元部分を除いて計算するtailsのHomが一致する．
また$B$の次数supportの仮定により，$\qgr B$は次数の剰余類modulo $d$で直和分解される．
従って$d\nmid n$ならば
\[
 \Hom_{\qgr A}(\OO,D(n))
 \simeq\Hom_{\qgr B}(\pi_BB,\pi_BB(n))=0.
\]

$c=\gcd(d,r)>1$と仮定する．全ての$m\geq0$について
$1+rm\equiv1\pmod c$なので$d\nmid(1+rm)$であり，
\[
 \Hom_{\qgr A}(\OO(-rm),D(1))
 \simeq\Hom_{\qgr A}(\OO,D(1+rm))=0.
\]
しかし$B$は無限次元なので$D(1)\ne0$である．
これはAZ1による有限直和$\bigoplus_i\OO(-rm_i)\twoheadrightarrow D(1)$に矛盾する．
\end{proof}


\begin{remark}[normal商へのample pairの降下]
上の因子化は，$\qgr B$から来る任意の標的$F$について成立する．
したがって$(\pi_AA,(r))$がAZ1とAZ2を満たせば，$(\pi_BB,(r))$もample pairである．
実際，AZ1の全射の各成分は$\pi_BB(-rm_i)$を経由し，その因子化後も全射である．
AZ2は，$\qgr B$の全射を完全忠実なexact functorで$\qgr A$へ移し，
両側のHomをこの因子化で同一視すれば従う．
\end{remark}

\subsection{整な切断商の成長}

\begin{lemma}[稠密な次数supportにおける成長の上界]
\label{lem:dense-support-growth}
$B=\bigoplus_{n\geq0}B_n$をlocally finite connected graded $\C$-domainとし，
$b_n=\dim B_n$とする．$B_n\ne0$ for all $n\gg0$を仮定する．
ある無限列$N\to\infty$に沿って
\[
 \sum_{j=0}^{N}b_j=\tfrac12N^2+O(N)
\]
ならば，任意の$n\geq0$について$b_n\leq n+1$である．
\end{lemma}

\begin{proof}
$n>0$で$b_n>0$の場合だけを考えればよい．
Domain上の有限次元部分空間についてのHopf不等式
$\dim UV\geq\dim U+\dim V-1$を反復すると，
\[
 \dim (B_n)^q\geq q(b_n-1)+1.
\]
各剰余類$r\bmod n$に対し，十分大きい$j_r\equiv r\bmod n$と
$0\ne w_r\in B_{j_r}$を選ぶ．
$w_r(B_n)^q\subset B_{j_r+qn}$より，すべての十分大きい$\ell$について
\[
 b_\ell\geq\frac{b_n-1}{n}\ell-C
\]
となる定数$C$が存在する．これを加えると
\[
 \sum_{\ell=0}^N b_\ell
 \geq \frac{b_n-1}{2n}N^2-O(N).
\]
仮定した無限列に沿って主要項を比較し，$b_n-1\leq n$を得る．
\end{proof}

\begin{theorem}[次数1の整なnormal切断による認識]
\label{thm:normal-domain-recognition}
元の六公理と固定したexactな三乗根$s$を仮定する．
軌道環$A$に非零normal元$g\in A_1$が存在し，
$B=A/(g)$がdomainならば，$A$は三次元二次Koszul AS-regular代数で，
Gorenstein parameterは$3$，Hilbert級数は$(1-t)^{-3}$である．
また$B$は二次元二次AS-regular代数であり，Hilbert級数は$(1-t)^{-2}$である．
従って$(L_0,L_1,L_2)$はfull strong exceptional collectionであり，
\[
 Q=L_{L_1}L_{L_2}(L_3)\simeq L_0[2],\qquad \PP=0.
\]
特に，この結論に$E$の有限性や全整数Euler多項式性を別途仮定する必要はない．
\end{theorem}

\begin{proof}
\emph{Hilbert関数の決定．}
$A$はdomainなので$g$は両側regularである．短完全列
\[
 0\longrightarrow A(-1)\xrightarrow{\,g\,}A\longrightarrow B\longrightarrow0
\]
から
\begin{equation}
\label{eq:hilbert-sum}
 b_n=h_n-h_{n-1},\qquad
 \sum_{j=0}^{3m}b_j=h_{3m}=\binom{3m+2}{2}
\end{equation}
を得る．ただし$h_{-1}=0$である．
$B$は無限次元である．そうでなければconnected graded domainであることから$B=\C$となり，
\eqref{eq:hilbert-sum}に反する．

$B$の非零な正次数の集合の最大公約数を$d$とする．
Domainなので，非零成分の次数は加法について閉じている．
最大公約数が$d$となる有限個の次数を選び，それらの生成する数値半群を考えると，
$B_n\ne0$ for all sufficiently large multiples $n$ of $d$である．
逆に$d\nmid n$では$B_n=0$である．
$d\geq4$なら$h_3=\sum_{j=0}^3b_j=1$となって矛盾する．
$d=3$は補題\ref{lem:residue-obstruction}に反する．
$d=2$なら$b_1=b_3=0$なので$h_3=10$から$b_2=9$である．
Hopf不等式により$b_4\geq17$，$b_6\geq25$を得る．
従って
\[
 28=h_6=b_0+b_2+b_4+b_6\geq1+9+17+25=52,
\]
という矛盾が生じる．よって$d=1$である．

補題\ref{lem:dense-support-growth}を
\eqref{eq:hilbert-sum}の列$N=3m$に適用すると，$b_n\leq n+1$である．
しかも$\sum_{j=0}^{3m}(j+1)=\binom{3m+2}{2}$であるから，
\eqref{eq:hilbert-sum}の等号と各項の非負性により
\begin{equation}
\label{eq:all-hilbert}
 b_n=n+1,\qquad h_n=\binom{n+2}{2}\qquad(n\geq0)
\end{equation}
がすべての次数で成り立つ．

\emph{$B$の一次生成と二次表示．}
$R=\langle B_1\rangle\subset B$と置く．$\dim B_1=2$なので，
Hopf不等式の反復から$\dim R_n\geq n+1$である．
\eqref{eq:all-hilbert}より$R=B$となる．
$V=B_1$とすると$V\otimes V\to B_2$の核は一次元であり，その生成元を$r$とする．
$r$のtensor rankが$1$ならば，非零な次数$1$の元の積が$B$で零になるのでdomain性に反する．
従って$r$の係数行列は可逆である．
$\C$上で基底を選び直すと，$r$は
\[
 yx-qxy-cx^2\qquad(q\ne0)
\]
の形にできる．実際，対応する二変数二次形式の等方ベクトルを選んで$y^2$の係数を消し，
係数行列の可逆性から両交差項の係数が非零であることを使えばよい．
この一関係による書換えは重なりの曖昧性をもたず，
$\C\langle x,y\rangle/(r)$の基底は$x^iy^j$である．
従ってそのHilbert級数は$(1-t)^{-2}$であり，$B$への自然な全射は同型である．
これは量子平面またはJordan平面であり，二次元二次AS-regularである．
この最後の標準的な事実は，$r=xu+yv$と書いたときの線形分解
\[
 0\longrightarrow B(-2)\xrightarrow{(u,v)^{\mathsf t}}
 B(-1)^2\xrightarrow{(x,y)}B\longrightarrow\C\longrightarrow0
\]
からも確認できる．左端の単射はdomain性による．中項の完全性は
$b_n=n+1$からの次元比較で従い，$u,v$が$B_1$の基底であることから
双対複体のcohomologyは最終項$\C(2)$だけである．
反対側にも同じ線形分解がある．

\emph{$A$の一次生成と二次表示．}
$x,y\in A_1$を$B_1$の基底のliftとする．
$B$が一次生成されることと$A/(g)=B$から，次数に関する帰納法により
$A$は$g,x,y$で生成される．
$ga=\sigma(a)g$を満たすgraded automorphismを上と同様に$\sigma$とする．
自由代数$T=\C\langle G,X,Y\rangle$上で，$\sigma|_{A_1}$と同じ係数を持つ
一次成分$T_1$のautomorphismを$\widetilde\sigma$と書く．
$\widetilde\sigma(G)=G$である．
$B$の唯一の二次関係を$T$へliftすると，その$A$での値は$gA_1$に属するので，
ある一次式$\ell\in T_1$について$r(X,Y)-G\ell$は$A$の関係式となる．
次の三つの二次式で定めた代数を$A'$と置く：
\[
 GX-\widetilde\sigma(X)G,\qquad
 GY-\widetilde\sigma(Y)G,\qquad
 r(X,Y)-G\ell.
\]
$G$は$A'$でnormalであり，$A'/(G)\simeq B$である．
実際，二つのnormality関係は$GA'\subset A'G$を与える．
$\widetilde\sigma(G),\widetilde\sigma(X),\widetilde\sigma(Y)$が一次成分の基底であるため，
逆の包含も一次成分から全積へ帰納して従う．
ここで$G$が$A'$でregularであるとはまだ仮定しない．
それでも各次数における右完全列から
\[
 \dim A'_n\leq\dim A'_{n-1}+\dim B_n
 \leq\sum_{j=0}^n(j+1)=\binom{n+2}{2}
\]
を帰納的に得る．一方，構成から$A'\twoheadrightarrow A$であり，
\eqref{eq:all-hilbert}より両辺の次元は等しい．
従って$A'=A$であり，$A$は二次代数である．

これで$A$は三生成元の二次domainで$H_A(t)=(1-t)^{-3}$である．
Iyudu--Shkarin\cite[\S12.1, p.~78]{IS}の分類の帰結
（同論文のクラス$\Omega$ではdomainと三次元AS-regular性が同値）を適用する．
これにより三次元二次Koszul AS-regular性とparameter $3$が従う．
最後に，AS-regular代数の標準helixとAZのtriple equivalenceを用いる
既証の有限認識定理の逆向きにより，fullness，$Q\simeq L_0[2]$，$\PP=0$を得る．
対応するKoszul resolutionを$\qgr A$へ移して三回shiftすると，
\[
 0\longrightarrow L_0\longrightarrow L_1^{\oplus3}
 \longrightarrow L_2^{\oplus3}\longrightarrow L_3\longrightarrow0
\]
というheart内の完全列も得られる．
\end{proof}

\begin{remark}[仮定・証明の依存関係]
数値的な部分では，$A$のdomain性，normal元によるHilbert差分，
$h_{3m}=\binom{3m+2}{2}$，および$B$のdomain性だけを用いた．
Serre関手の正確な三乗根は，残る剰余次数$d=3$を
補題\ref{lem:residue-obstruction}で排除する際に使っている．
この使い方は，単なるHilbert級数の計算では代替できない．
また，normal extensionの一般的なglobal dimension上昇定理は使っていない．
一次生成・二次関係を直接証明した後に，既知の三次元二次domainの認識を適用した．
元の公理からそのような$g$とdomain商の存在を証明したわけではない．
本定理の新規性については別途調査を要する．
\end{remark}


\section{次数2の整なnormal切断による認識}
\label{sec:degree2-hilbert}

元の六公理と固定した完全三乗根$s$を仮定する．
$A=\bigoplus_{n\geq0}\Hom(\OO,s^n\OO)$とし，
$0\ne g\in A_2$がnormalで，$B=A/gA$がdomainであると仮定する．
$A$がdomainであることは既証なので，$g$は両側regularである．
$h_n=\dim_\C A_n$，$b_n=\dim_\C B_n$と書き，負次数では両者を零とする．
既証の再構成により$(\Cat,\OO,s)\simeq(\qgr A,\pi A,(1))$であり，
\[
 \SSer=(-3)[2],\qquad h_{3m}=P(3m),\qquad
 P(n)=\frac{(n+1)(n+2)}2.
\]


\begin{theorem}[次数2の整なnormal切断による認識]
\label{thm:quadratic-normal-recognition}
六公理と固定したexactな三乗根$s$の下で，$0\ne g\in A_2$をnormal元とし，
$B=A/(g)$がdomainであると仮定する．このとき
\[
 H_A(t)=\frac1{(1-t)^3},\qquad
 B\simeq B(\mathbb P^1,\OO(2),\sigma)
\]
となり，$A$は三次元二次Koszul AS-regular代数，Gorenstein parameterは3である．
特に$A$は一次生成され，$E_n=0$ for all $n\in\Z$であり，
$(\OO,\OO(1),\OO(2))$はfull exceptional collectionとなる．
ここで$\sigma$は$\mathbb P^1$のある自己同型である．
\end{theorem}

証明を，Hilbert関数の決定，切断商の識別，独立した次数2生成元の排除，
Koszul性の持上げに分ける．切断商の曲線分類は既知結果を用いるが，
ambientな$A$の一次生成性や二次性をその分類の仮定に含めない．

\subsection{商の次数support}

次数supportの制御には，一般のnormal切断の剰余次数障害補題
\ref{lem:residue-obstruction}を用いる．
同補題は，$\deg g=e$，$\SSer=(-\ell)[2]$，
$B_n=0$ whenever $d\nmid n$の下で
\[
 d\mid(e-\ell)
\]
が必要であることを述べる．本節では$(e,\ell)=(2,3)$なので，$d=1$が必要である．
同補題の証明はnormal境界の構造対象に対する双対性を使い，
Euler多項式性や$E$の有限性に依存しない．

\subsection{二つの累積和による全Hilbert関数の決定}

\begin{theorem}[二次の整なnormal切断のHilbert関数]
\label{thm:degree2-hilbert}
元の六公理と固定した完全三乗根$s$を仮定し，
$0\ne g\in A_2$がnormalで$B=A/gA$がdomainであるとする．
このとき全ての$n\geq0$について
\[
 \boxed{\quad b_n=2n+1,\qquad h_n=P(n).\quad}
\]
従って
\[
 H_B(t)=\frac{1+t}{(1-t)^2},\qquad H_A(t)=\frac1{(1-t)^3}.
\]
さらに$B$は次数$1$で生成され，$A$は$A_1$と$g$で生成される．
この主張には，全整数Euler多項式性や$E$の有限性を先行して仮定しない．
\end{theorem}

\begin{proof}
$g$による短完全列
\[
 0\longrightarrow A(-2)\xrightarrow{\,g\,}A
 \longrightarrow B\longrightarrow0
\]
より
\begin{equation}
\label{eq:parity-sum}
 b_n=h_n-h_{n-2},\qquad
 h_N=\sum_{\substack{0\leq j\leq N\\j\equiv N\ (\mathrm{mod}\ 2)}} b_j.
\end{equation}
特に$b_3+b_1=h_3=10$である．従って$B$の正次数部分は非零であり，
domain性から$B$は無限次元である．

$B$の非零な正次数の最大公約数を$d$とする．
非零成分の次数は加法について閉じており，
最大公約数が$d$となる有限個の次数の生成する数値半群を考えれば，
$B_n\ne0$ for all sufficiently large multiples $n$ of $d$である．
逆に$d\nmid n$なら$B_n=0$である．
補題\ref{lem:residue-obstruction}を$(e,\ell)=(2,3)$に適用すると，
$d\mid-1$または$d\mid1$が必要であり，従って$d=1$である．
よって$B_n\ne0$ for all $n\gg0$となる．

固定した$n>0$について$b_n>0$の場合を考える．
Domainの積に対するHopf不等式を反復すると
\[
 \dim(B_n)^q\geq q(b_n-1)+1\qquad(q\geq0)
\]
である．各剰余類$r\bmod n$について，十分大きい
$j_r\equiv r\bmod n$と$0\ne w_r\in B_{j_r}$を選ぶ．
$w_r(B_n)^q\subset B_{j_r+qn}$より，ある定数$C$について
\[
 b_\ell\geq\frac{b_n-1}{n}\ell-C\qquad(\ell\gg0)
\]
となる．\eqref{eq:parity-sum}では一方の偶奇の次数だけを加えるので，
\[
 h_N\geq\frac{b_n-1}{4n}N^2-O(N).
\]
$N=3m\to\infty$に沿って$h_N=P(N)=N^2/2+O(N)$と比較すると，
$b_n\leq2n+1$を得る．$b_n=0$の場合と$n=0$にも同じ上界が成り立つ．

標準列は
\[
 \sum_{\substack{0\leq j\leq N\\j\equiv N\ (\mathrm{mod}\ 2)}}(2j+1)=P(N)
\]
を満たす．従って$N=3m$について
\[
 \sum_{\substack{0\leq j\leq3m\\j\equiv3m\ (\mathrm{mod}\ 2)}}
 \bigl((2j+1)-b_j\bigr)=0.
\]
各項は非負である．任意の$j\geq0$に対し，$3m\geq j$かつ$m\equiv j\pmod2$となる$m$を選べるので，
$b_j=2j+1$がすべての次数で従う．
\eqref{eq:parity-sum}へ戻れば$h_n=P(n)$も従う．

$R=\langle B_1\rangle$と置くと$\dim B_1=3$であり，
Hopf不等式により$\dim R_n\geq2n+1$である．
ところが$R_n\subset B_n$かつ$\dim B_n=2n+1$なので$R=B$となる．
最後に，$B$の一次生成元を$A_1$へliftし，
$A/gA=B$と次数に関する帰納法を使えば，$A$は$A_1$と$g$で生成される．
\end{proof}

\subsection{生成元の二分}

\begin{proposition}[一次生成に残る二つの場合]
\label{prop:degree2-generation-gap}
定理\ref{thm:degree2-hilbert}の仮定の下で
\[
 \dim A_1=3,\qquad \dim A_2=6,\qquad
 \dim(A_1A_1)\in\{5,6\}.
\]
$A$が一次生成されること，$g\in A_1A_1$，$\dim(A_1A_1)=6$は同値である．
\end{proposition}

\begin{proof}
最初の二式はHilbert関数から従う．Hopf不等式から
$\dim(A_1A_1)\geq2\dim A_1-1=5$であり，上界は$\dim A_2=6$である．
$A_1A_1\to B_2$は$B$の一次生成性により全射で，$\dim B_2=5$，
$\ker(A_2\to B_2)=\C g$である．
従って$g\in A_1A_1$は$\dim(A_1A_1)=6$と同値である．
この条件の下では，$A$が$A_1$と$g$で生成されることから一次生成性が従う．
逆は定義から明らかである．
\end{proof}


\subsection{線形成長からの整なconicの識別}
\label{subsec:domain-conic-identification}

以下では$k=\C$とする．Artin--Staffordの曲線分類
\cite[Theorem~0.2]{ASquadratic}のうち，用いる部分を明記する．
一次生成のconnected graded domain $B$が$\GKdim B=2$を満たすならば，
integral projective curve $Y$，$\sigma\in\operatorname{Aut}(Y)$，
ample invertible sheaf $\mathcal L$とgraded単射
\[
 B\lhook\joinrel\longrightarrow
 T=B(Y,\mathcal L,\sigma)
   =\bigoplus_{n\geq0}H^0(Y,\mathcal L_n),
 \qquad
 \mathcal L_n=\mathcal L\otimes\sigma^*\mathcal L\otimes\cdots
                   \otimes(\sigma^{n-1})^*\mathcal L
\]
があり，$\dim_k(T/B)<\infty$である．
この主張は共著者による\cite[Theorem~5.2]{SVdBsurvey}にも再掲されている．
ambientなAS-regular代数や，$\qgr B$の滑らかさはこの定理の仮定ではない．

\begin{proposition}[整なHilbert conicの識別]
\label{prop:domain-conic-identification}
$B$をconnected graded $\C$-domainとし，$B$は$B_1$で生成され，
\[
 \dim_\C B_n=2n+1\qquad(n\geq0)
\]
を満たすと仮定する．このときある$\sigma\in\operatorname{Aut}(\mathbb P^1)$について
\[
 B\simeq B(\mathbb P^1,\OO_{\mathbb P^1}(2),\sigma)
\]
がgraded代数として成り立つ．特に$B$は
\[
 C=\C[u,v,w]/(uw-v^2),\qquad\deg u=\deg v=\deg w=1,
\]
のgraded automorphismによるZhang twistであり，二次Koszul代数である．
その自明加群のminimal graded free resolutionの第$n$項は
\[
 P_0=B,\qquad P_1=B(-1)^{\oplus3},\qquad
 P_n=B(-n)^{\oplus4}\quad(n\geq2)
\]
となる．
\end{proposition}

\begin{proof}
まず$\sum_{i=0}^n\dim B_i=(n+1)^2$なので$\GKdim B=2$である．
上記Artin--Stafford定理を適用し，$B\hookrightarrow T=B(Y,\mathcal L,\sigma)$を取る．
$d=\deg\mathcal L>0$，$p_a(Y)=\dim H^1(Y,\OO_Y)$と置く．
automorphismのpullbackは次数を保つので$\deg\mathcal L_n=nd$である．
integral projective curve上では十分大きい次数のinvertible sheafの$H^1$は零となり，
Riemann--Rochにより
\[
 h^0(Y,\mathcal L_n)=nd+1-p_a(Y)\qquad(n\gg0)
\]
である．ここでは$\mathcal L_n$が一つのsheafの通常の冪である必要はない．
次数だけに依存する消滅上界を使えばよい．
$T/B$が有限次元であることから左辺は大きい$n$で$2n+1$に等しく，
従って
\[
 d=2,\qquad p_a(Y)=0.
\]

この算術種数の消滅は，$Y$の滑らかさも含意する．
実際，normalization $\nu:X\to Y$は有限で，$X$はsmooth projective curveである．
完全列
\[
 0\longrightarrow\OO_Y\longrightarrow\nu_*\OO_X
   \longrightarrow\mathcal T\longrightarrow0
\]
において$\mathcal T$は有限長である．$Y,X$はintegral projectiveなので
$H^0(Y,\OO_Y)=H^0(X,\OO_X)=\C$であり，
\[
 p_a(Y)=g(X)+\operatorname{length}(\mathcal T)
\]
を得る．両項は非負なので$g(X)=0$かつ$\mathcal T=0$である．
有限射$\nu$について$\nu_*\OO_X=\OO_Y$となるため$\nu$は同型であり，
$Y$自体がsmooth genus-zero curveである．
代数閉体上のsmooth projective genus-zero curveは$\mathbb P^1$なので
$Y\simeq\mathbb P^1$である．

$\operatorname{Pic}(\mathbb P^1)\simeq\Z$より
$\mathcal L\simeq\OO(2)$，$\mathcal L_n\simeq\OO(2n)$であり，
今度は\emph{すべての}$n\geq0$について
\[
 \dim T_n=h^0(\mathbb P^1,\OO(2n))=2n+1=\dim B_n.
\]
従って，もともと全次数で定義された単射$B\hookrightarrow T$は同型である．
漸近的な同型だけから低次数の同型を推論しているのではなく，
全次数の単射と各次数の次元一致を用いている．

twistの記述を具体化する．$\mathcal L=\OO(2)$とし，同型
$\psi:\sigma^*\mathcal L\xrightarrow{\sim}\mathcal L$を一つ選ぶ．
これにより通常のsection ring
\[
 C=\bigoplus_{n\geq0}H^0(\mathbb P^1,\mathcal L^{\otimes n})
   \simeq\C[x,y]^{(2)}
   \simeq\C[u,v,w]/(uw-v^2)
\]
上にgraded automorphism
\[
 \phi_n=H^0(\psi^{\otimes n})\circ\sigma^*:C_n\longrightarrow C_n
\]
が定まる．$u=x^2,v=xy,w=y^2$とし，Veroneseの次数は半分に取り直した．
$\psi$を反復して各$\sigma^{i*}\mathcal L$を$\mathcal L$と同一視すると，
THCRの積は同じgraded vector space $C$上で
\[
 a*b=a\phi^{m}(b)\qquad(a\in C_m)
\]
となる．これは$\phi$によるZhang twist $C^\phi$の定義そのものである．

最後にKoszul性とBetti数を直接確認する．
$S=\C[u,v,w]$，$f=uw-v^2$，$C=S/(f)$とし，
$K=K_C(u,v,w)$を三つの変数のKoszul complexとする．
$S$上の完全列$0\to S(-2)\xrightarrow{f}S\to C\to0$を使うと
\[
 H_0(K)=\C,\qquad H_1(K)=\C(-2),\qquad H_i(K)=0\quad(i\geq2).
\]
$H_1(K)$の生成元はcycle
$c=w e_u-v e_v$で表せる．これは$S$上へ持ち上げると$d(c)=f$となることから分かる．
homological degree $2$，internal degree $2$のdivided-power変数$\eta$を加え，
\[
 F=K\otimes_\C\Gamma(\eta),\qquad
 d(\eta^{[j]})=c\eta^{[j-1]}\quad(j\geq1)
\]
と定める．$d(c)=0$かつ$c^2=0$より$d^2=0$である．
$\eta$の次数でfilterすると，最初のhomologyは
$(\C\oplus\C c)\otimes\Gamma(\eta)$であり，
正次数の項は$\eta^{[j]}\mapsto c\eta^{[j-1]}$により対になって消える．
このfiltrationは各homological degreeで有限なので，$F$は$\C$のfree resolutionである．
基底$e_I\eta^{[j]}$のhomological degreeとinternal degreeはともに$|I|+2j$であり，
微分の全係数は$C_1$に属する．従ってこれはminimal linear resolutionである．
free rankの母関数は
\[
 \sum_{n\geq0}\beta_n t^n
 =\frac{(1+t)^3}{1-t^2}
 =\frac{(1+t)^2}{1-t}
 =1+3t+4t^2+4t^3+\cdots.
\]

graded automorphismによるtwistingはgraded module圏のexact同値を与え，
自明加群とfree moduleの次数を保つ．
具体的にはhomogeneousな$m$に対して右作用を
$m*a=m\phi^{\deg m}(a)$と変えればよく，逆functorは$\phi^{-1}$で与えられる．
従って上記resolutionは$B=C^\phi$のminimal linear resolutionに移り，
主張したKoszul性とBetti数を得る．
特に第2syzygyの生成次数は2なので$B$は二次代数である．
同値な記述として，$C$のtensor-algebra表示における四つの二次関係
\[
 uv-vu,\quad uw-wu,\quad vw-wv,\quad uw-v^2
\]
が張る空間を$R\subset C_1\otimes C_1$とすれば，
$B$の関係空間は$(\id\otimes\phi^{-1})(R)$である．
\end{proof}


\subsection{周期的な頂点対象を使う分裂Ore拡大の排除}

以下の補題はnormal extension一般のglobal dimensionの上昇を主張しない．
環が実際にskew polynomial extensionとして分裂した場合だけを扱う．

\begin{lemma}[分裂拡大の頂点対象のExt]\label{lem:split-vertex}
$B$をlocally finite connected graded右noetherian代数，
$\nu$を次数を保つ自己同型とし，
\[
 A=B[z;\nu],\qquad zb=\nu(b)z,\qquad \deg z=e>0
\]
が右noetherianであるとする．$\Cat=\qgr A$，$\OO=\pi A$と置く．
$M=A/B_{\geq1}A\simeq\C[z]$，$N=\pi M$とする．
$k_B$の最小次数付き自由分解の第$i$項を
$P_i=\bigoplus_jB(-a_{ij})$と書くと，任意の$r\in\Z$について
\begin{equation}\label{eq:vertex-ext}
 \Ext^i_{\Cat}(N,N(r))
 \simeq\bigoplus_j\C[z,z^{-1}]_{a_{ij}+r}
 \qquad(i\geq0)
\end{equation}
というベクトル空間の同型がある．
特に$P_i\ne0$なら，ある$r\in\{0,\ldots,e-1\}$について左辺は非零である．
\end{lemma}
\begin{proof}
$D=A[z^{-1}]$，$\widetilde M=M[z^{-1}]\simeq\C[z,z^{-1}]$と置く．
次数付きOre局所化
\[
 L:\operatorname{Gr}A\longrightarrow\operatorname{Gr}D
\]
はexactであり，exactな制限functor $\operatorname{res}$を右随伴に持つ．
従って$\operatorname{res}$はinjectiveをinjectiveに送り，任意のtorsion加群$T$と
$D$加群$V$について
\[
 \Ext^i_{\operatorname{Gr}A}(T,\operatorname{res}V)
 \simeq\Ext^i_{\operatorname{Gr}D}(T[z^{-1}],V)=0
 \qquad(i\geq0)
\]
である．ここでtorsion加群は局所化で零になる．
従って$\operatorname{res}V$はtorsion-localであり，$L$が誘導するexactなfunctor
$\overline L:\operatorname{QGr}A\to\operatorname{Gr}D$の右随伴
$J=\pi\operatorname{res}$はexactかつ完全忠実である．
また$\overline L$のexactnessから$J$はinjectiveを保つ．よって
\begin{equation}\label{eq:localized-ext}
 \Ext^i_{\operatorname{QGr}A}(JV,JW)
 \simeq\Ext^i_{\operatorname{Gr}D}(V,W)
\end{equation}
が全次数で成り立つ．

$\widetilde M/M$は$A$-torsionである．実際，$B_{\geq1}$は零作用し，
各元はある正の$z$冪で消える．従って$N\simeq J\widetilde M$であり，
式\eqref{eq:localized-ext}とexactな局所化の随伴性から
\begin{align*}
 \Ext^i_{\operatorname{QGr}A}(N,N(r))
 &\simeq\Ext^i_{\operatorname{Gr}D}(\widetilde M,\widetilde M(r))\\
 &\simeq\Ext^i_{\operatorname{Gr}A}(M,\widetilde M(r)).
\end{align*}
$A$は左$B$加群として自由なので，$P_\bullet\otimes_B A$は$M$の自由分解となる．
そのminimalな微分の行列成分は$B_{\geq1}$に入り，標的$\widetilde M(r)$には零作用する．
従って$\Hom_{\operatorname{Gr}A}(-,\widetilde M(r))$を適用した複体の微分は全て零で，
その第$i$項は$\bigoplus_j\C[z,z^{-1}]_{a_{ij}+r}$である．

最後に，$\operatorname{QGr}A$はlocally noetherianであり，そのnoetherian対象の充満部分圏は
$\qgr A$である．この包含はnoetherian対象間のYoneda Extを全次数で保つ．
これは有限長の拡大図式の各項を，右端への全射を保ち，左端を含むnoetherian部分対象で
順次置き換えることにより従う．同様の置換を有限な比較図式にも行えば単射性も得る．
従って上の計算は$\qgr A$における式\eqref{eq:vertex-ext}を与える．
最後に，非零summandの$a_{ij}$について$r\equiv-a_{ij}\pmod e$を選べばよい．
\end{proof}

\begin{corollary}[有限ホモロジー次元からの分裂境界の制約]\label{cor:split-pd}
前補題の仮定の下で$\operatorname{hdim}\Cat\leq d$ならば，
$\operatorname{pd}_{\operatorname{Gr}B}k\leq d$である．
従って$1/H_B(t)$はLaurent多項式でなければならない．
\end{corollary}
\begin{proof}
式\eqref{eq:vertex-ext}より$i>d$では$P_i=0$である．
有限な自由分解のHilbert級数の交代和を取ると
\[
 1=H_B(t)\sum_{i,j}(-1)^it^{a_{ij}},
\]
となる．
\end{proof}

\begin{proposition}[Conic型のskew polynomial分裂は不可能]\label{prop:conic-split-exclusion}
六公理と固定したexactな三乗根の下で，
$A\simeq B[z;\nu]$，$\deg z=2$，
$H_B(t)=(1+t)/(1-t)^2$となることはない．
\end{proposition}
\begin{proof}
$\operatorname{hdim}\Cat=2$は基礎公理の帰結である．
従って系\ref{cor:split-pd}が適用できるが，
\[
 \frac1{H_B(t)}=\frac{(1-t)^2}{1+t}
\]
は$t=-1$に極を持ち，Laurent多項式ではない．矛盾する．
\end{proof}



\begin{proposition}[独立した次数2生成元の排除]\label{prop:conic-degree-one-generation}
定理\ref{thm:quadratic-normal-recognition}の仮定の下で，$A$は$A_1$で生成される．
\end{proposition}
\begin{proof}
命題\ref{prop:degree2-generation-gap}の二分のうち，$\dim A_1^2=5$の場合を排除する．
$V=A_1\simeq B_1$と置くと，$A_1^2\to B_2$は同型である．
命題\ref{prop:domain-conic-identification}より$B$は二次代数であるから，
$B$の全ての定義関係は$V$のliftにおいても零になる．
従って商写像$p:A\to B$にはgraded section $j:B\hookrightarrow A$があり，
その像は$\langle A_1\rangle$である．

$ga=\nu(a)g$で定まる$A$のgraded automorphismは$A_1$を保つので，$j(B)$も保つ．
同じ記号$\nu$で$j(B)\simeq B$への制限を表せば，$A$が$A_1$と$g$で生成されることから
\[
 B[z;\nu]\twoheadrightarrow A,\qquad z\longmapsto g,\quad\deg z=2
\]
を得る．左辺のHilbert級数は
\[
 \frac{H_B(t)}{1-t^2}=\frac1{(1-t)^3}=H_A(t)
\]
なので，この全射は同型である．しかし命題\ref{prop:conic-split-exclusion}がこの分裂拡大を排除する．
従って$\dim A_1^2=6$であり，$g\in A_1^2$と一次生成性が従う．
\end{proof}

\begin{corollary}[二次境界における全Ext障害の消滅]\label{cor:conic-all-ext-vanishing}
定理\ref{thm:quadratic-normal-recognition}の仮定の下で，$E_n=0$ for all $n\in\Z$であり，
$q_n=P(n)$は全整数で成立する．
\end{corollary}
\begin{proof}
命題\ref{prop:domain-conic-identification}で得たTHCRは全次数で飽和している．
従って命題\ref{prop:saturated-boundary}より$E=0$である．
$h_n=P(n)$ for $n\geq0$とSerre双対性を用いれば，
$q_n=h_n+h_{-n-3}=P(n)$は負次数を含めて従う．
この証明は$A$の二次性やAS-regular性を先に用いない．
\end{proof}

\subsection{一次生成されたconic被覆の二次性}
\label{sec:conic-cover-koszul}

ここでは，次数2のnormal切断の商が滑らかなconicのtwistと同定され，
被覆環$A$の一次生成性も得られた段階から，$A$の二次性を導く．
一般の「Koszulな商を持つ環はKoszul」という逆命題は用いない．
必要な追加構造は，conicのYoneda代数がgraded-primeであることである．

\subsubsection{左右のchange-of-rings列からのnormal性}

ここでは軌道$\Ext^1$を表す既存の$E$と区別するため，
connected graded algebra $R$のYoneda代数を
\[
 \mathcal Y^{p,q}(R)
 =\Ext^p_{\operatorname{gr}R}(\C,\C(-q)),\qquad
 \mathcal Y^p(R)=\bigoplus_q\mathcal Y^{p,q}(R),\qquad
 \mathcal Y(R)=\bigoplus_p\mathcal Y^p(R)
\]
と書く．$q$は最小分解の生成次数を正に数える内部weightであり，
Koszul性は$p=q$以外での消滅を意味する．
$\langle2\rangle$はこの内部weightを$2$増やすshiftを表す．
$\operatorname{gr}R$は本稿の規約に従い右graded moduleの圏とする．
Cassidy--Sheltonの原典は左加群で記述されているので，
そのchange-of-rings列は$A^{\mathrm{op}}\to B^{\mathrm{op}}$へ適用して用いる．
左右のmodule列の呼び名が入れ替わっても，以下で用いる二つの像と共通kernelの等式は変わらない．
またKoszul双対とYoneda代数の同一視に伴う反対環の慣習は，
本節で用いるgraded primeness，Koszul対角性およびHilbert級数に影響しない．

\begin{lemma}[graded-primeなYoneda代数による持ち上げ]
\label{lem:prime-yoneda-lift}
$A$をconnected locally finite right noetherian graded $\C$-algebraとし，
$A_1$で生成されると仮定する．
$0\ne g\in A_2$が両側regularかつnormalで，$B=A/gA$とする．
$B$がKoszulであり，$\mathcal Y(B)$がcohomological gradingに関してgraded-primeならば，
$A$もKoszulである．より詳しく，bidegree $(2,2)$の非零normal regular元
$\eta\in\mathcal Y(B)$が存在して
\[
 \mathcal Y(A)\simeq\mathcal Y(B)/(\eta)
\]
となる．
\end{lemma}

\begin{proof}
$A$のright noetherian性と有限次元の$A_1$から，
augmentation idealの最初のsyzygyも有限生成である．
従って$A$は代数として有限個の斉次関係による表示を持つ．
$A^{\mathrm{op}}$も有限表示を持つので，
以下で用いるCassidy--Sheltonの原典の共通設定とも整合する．

$\phi:A\twoheadrightarrow B$に伴うrestrictionを
$\phi^*:\mathcal Y(B)\to\mathcal Y(A)$とする．
Cassidy--Shelton\cite[\S6, (6.1), (6.2)]{CassidyShelton}の二つのchange-of-rings長完全列を用いる．
原典は左加群で記述されているが，反対環に適用して右加群版を得る．
Yoneda積はKoszul代数に対して$\mathcal Y(R)\simeq R^!$となる規約を用いる．
積の逆の規約を選んだ場合も，以下のnormal性・regularity・graded primenessは反対環で不変である．
これらはともに$g$のregularityとnormalityの下で存在し，
$A$のKoszul性や$\mathcal K_2$性を仮定しない．
第一の列は右$\mathcal Y(B)$-moduleの列
\begin{equation}
\label{eq:right-cor}
 \cdots\longrightarrow
 \mathcal Y^{p-2}(B)\langle2\rangle
 \xrightarrow{d_2}\mathcal Y^p(B)
 \xrightarrow{\phi^*}\mathcal Y^p(A)
 \longrightarrow\mathcal Y^{p-1}(B)\langle2\rangle
 \longrightarrow\cdots
\end{equation}
である．第二の列は同じ$\phi^*$を持つ左$\mathcal Y(B)$-moduleの列
\begin{equation}
\label{eq:left-cor}
 \cdots\longrightarrow
 \mathcal Y^{p-2}(B)\langle2\rangle
 \xrightarrow{\bar d_2}\mathcal Y^p(B)
 \xrightarrow{\phi^*}\mathcal Y^p(A)
 \longrightarrow\mathcal Y^{p-1}(B)\langle2\rangle
 \longrightarrow\cdots
\end{equation}
である．二つの列そのものをbimoduleの一つの列と同一視する必要はない．
右および左のmodule線形性から，それぞれ
\[
 d_2(a)=\eta a,\qquad \bar d_2(a)=a\eta'
\]
と書ける．ここで$\eta,\eta'$はbidegree $(2,2)$の元である．

$A$は一次生成であり，$g$は次数2なので，
$\phi^*:\mathcal Y^1(B)\to\mathcal Y^1(A)$は同型である．
従って両方の長完全列で，$\mathcal Y^0(B)\langle2\rangle$上の$d_2$と$\bar d_2$は単射である．
特に$\eta\ne0$，$\eta'\ne0$であり，
$\ker(\phi^*:\mathcal Y^2(B)\to\mathcal Y^2(A))$は一次元である．
全次数におけるexactnessから
\[
 \eta\mathcal Y(B)=\ker\phi^*=\mathcal Y(B)\eta'
\]
を得る．次数2部分を比べれば$\eta'=c\eta$，$c\in\C^*$である．
従って$\eta$はnormalである．
このnormal性の導出は，原典のTheorem 9.1の逆向きではなく，
\S6の二つの長完全列だけを使っている．

Graded-prime ringの非零斉次normal元はregularである．
実際，斉次元$x\ne0$について$\eta x=0$ならば，normalityから
$\eta\mathcal Y(B)x=0$となり，graded primenessに反する．
$x\eta=0$の場合も同様である．
したがって\eqref{eq:right-cor}の$d_2$は全次数で単射であり，
$\phi^*$は全次数で全射となる．これにより
$\mathcal Y(A)=\mathcal Y(B)/(\eta)$を得る．
$B$のKoszul性から右辺はKoszul対角上に集中するので，$A$もKoszulである．
\end{proof}

\subsubsection{Conicの双対とそのtwistのgraded primeness}

\begin{lemma}[conicのKoszul双対]
\label{lem:conic-dual-prime}
$C=\C[x,y,z]/(xz-y^2)$を通常の次数1生成の滑らかなconicの座標環とする．
任意のgraded automorphism $\tau$に対し，Zhang twist $B=C^\tau$はKoszulであり，
$\mathcal Y(B)$はgraded-primeである．
\end{lemma}

\begin{proof}
$C$とそのZhang twistのKoszul性は，命題\ref{prop:domain-conic-identification}の
Tate型の線形自由分解で既に確認した．ここではその双対の積構造を調べる．

$\C$上の非退化な三変数二次形式は対角化できるので，$C^!$は
\[
 R=\C\langle\xi_1,\xi_2,\xi_3\rangle/
 \bigl(\xi_i\xi_j+\xi_j\xi_i\ (i<j),\
 \xi_1^2-\xi_2^2,\ \xi_1^2-\xi_3^2\bigr)
\]
と書ける．共通の平方を$w=\xi_i^2$とすると，$w$は次数2の中心元であり，
書換えにより$R$は$\C[w]$上で
$\xi_1^{\epsilon_1}\xi_2^{\epsilon_2}\xi_3^{\epsilon_3}$，
$\epsilon_i\in\{0,1\}$の八元で生成される．

Pauli行列
\[
 P_1=\begin{pmatrix}0&1\\1&0\end{pmatrix},\quad
 P_2=\begin{pmatrix}0&-\mathrm i\\\mathrm i&0\end{pmatrix},\quad
 P_3=\begin{pmatrix}1&0\\0&-1\end{pmatrix}
\]
を使い，$\xi_i\mapsto tP_i$により$R\to M_2(\C[t])$を定める．
$P_i^2=I$，$P_iP_j=-P_jP_i$であるから，これはwell-definedである．
上の八元による次数ごとの上界と，行列の像による下界が一致するので，この写像は単射であり，
像は
\begin{equation}
\label{eq:clifford-matrix}
 R_0=\C I,\qquad R_1=t\,\mathfrak{sl}_2(\C),\qquad
 R_n=t^nM_2(\C)\quad(n\geq2)
\end{equation}
である．特に八元は実際に$\C[w]$-自由基底である．

さらに，$R$の任意のgraded automorphism $\rho$によるZhang twist $R^\rho$が
graded-primeであることを直接示す．
積の規約を$a*b=a\rho^{\deg a}(b)$とする．
非零斉次元$a\in R_m$，$b\in R_n$を取り，$k\geq2$を固定する．
$c\in R_k$について
\[
 a*c*b=a\rho^m(c)\rho^{m+k}(b)
\]
である．\eqref{eq:clifford-matrix}により，
$\rho^m(R_k)=t^kM_2(\C)$である．
左右の$a$と$\rho^{m+k}(b)$は非零行列なので，
両者の間に適切なmatrix unitを入れれば積は非零になる．
従ってある$c$について$a*c*b\ne0$である．
これにより$R^\rho$はgraded-primeである．
一般のZhang twistがprime性を保つとは仮定していない．

最後に$B^!$がこの形のtwistであることを確認する．
$V=C_1$，二次関係空間を$J\subset V\otimes V$とすると，
$C^\tau$の関係空間は$(\id\otimes\tau^{-1})J$である．
その直交空間は$(\id\otimes\tau^*)J^\perp$であるから，
\[
 (C^\tau)^!\simeq(C^!)^\rho,
 \qquad \rho=(\tau^*)^{-1}.
\]
$\rho$は$C^!$のgraded automorphismである．
したがって$B^!$はgraded-primeである．
Koszul性により$\mathcal Y(B)$は$B^!$またはYoneda積の規約に対応するその反対環と同型であり，
graded primenessは反対環でも不変なので結論を得る．
\end{proof}

\subsubsection{次数2の被覆への適用}

\begin{proposition}[一次生成されたconic被覆の認識]
\label{prop:conic-cover-regular}
$A$をconnected locally finite right noetherian graded domainとし，
$A_1$で生成され，$H_A(t)=(1-t)^{-3}$であるとする．
$0\ne g\in A_2$がnormalで，$B=A/gA$が
滑らかなconicの座標環$C=\C[x,y,z]/(xz-y^2)$のZhang twistならば，
$A$は三次元二次Koszul AS-regular代数で，Gorenstein parameterは$3$である．
\end{proposition}

\begin{proof}
Domain性から$g$はregularである．
補題\ref{lem:conic-dual-prime}と補題\ref{lem:prime-yoneda-lift}により，
$A$はKoszulであり，特に二次代数である．
Yoneda代数のHilbert級数も
\[
 H_{\mathcal Y(A)}(t)
 =(1-t^2)H_{\mathcal Y(B)}(t)
 =\frac{1-t^2}{H_B(-t)}
 =(1+t)^3
\]
と決まる．ここで$H_B(t)=(1+t)/(1-t)^2$を使った．
最後に$A$は三生成元の二次domainで$H_A(t)=(1-t)^{-3}$なので，
既証の認識と同じくIyudu--Shkarin\cite[\S12.1, p.~78]{IS}のクラス$\Omega$に対する
domain--AS-regular認識を適用する．
\end{proof}



\begin{proof}[定理\ref{thm:quadratic-normal-recognition}の証明]
定理\ref{thm:degree2-hilbert}により，全Hilbert関数と$B$の一次生成性を得る．
命題\ref{prop:domain-conic-identification}で$B$をconicのZhang twistと同定し，
命題\ref{prop:conic-degree-one-generation}で$A$の一次生成性を得る．
従って命題\ref{prop:conic-cover-regular}が適用でき，$A$は主張したAS-regular代数である．
$E$の全消滅は系\ref{cor:conic-all-ext-vanishing}，
full exceptional collectionの結論は定理\ref{thm:finite}による．
\end{proof}

\begin{remark}[使用した公理と境界の次数]
Hilbert関数の決定には，domainの積不等式，反標準次数のHilbert値，
Serre双対性からの剰余障害を用いた．
有限ホモロジー次元は，独立した次数2生成元が作る分裂拡大を排除する箇所で不可欠である．
二次性の持上げにはconic双対のgraded primenessを実際に証明した．
これらの段階を省き，normal extension一般のglobal dimension上昇やKoszul性の逆命題で代用することはできない．
次数3については，節\ref{sec:elliptic}で一次生成または可換な整境界の場合の認識を得る．
一次生成を仮定しない一般の非可換境界の認識は未解決である．
\end{remark}


\section{一次生成で双有理的に幾何的な場合}\label{sec:birational}

\begin{theorem}[双有理的に可換な部分クラスの認識]\label{thm:birational-recognition}
基礎公理と三乗根の下で，$A=A_s$がさらに次を満たすとする．
\begin{enumerate}[label=(\roman*)]
\item $A$は左noetherianである．
\item $A$は$A_1$で生成される．
\item 次数付き商環は$Q_{\mathrm{gr}}(A)\simeq\C(Y)[t,t^{-1};\sigma]$であり，
$Y$はintegral projective surface，$\sigma$は$Y$の自己同型から誘導される．
\end{enumerate}
このとき，ある$\alpha\in\operatorname{Aut}(\mathbb P^2)$，$\alpha^3=\id$について
\[
 (\Cat,\OO,s)\simeq
 \bigl(\coh\mathbb P^2,\OO_{\mathbb P^2},
       \alpha^*(-)\otimes\OO_{\mathbb P^2}(1)\bigr).
\]
従って固定した軌道環$A$は三変数多項式環のZhang twistであり，
定理\ref{thm:finite}のAS-regular認識が成立する．
また反標準環は，次数を取り直した三次Veronese
$A^{(3)}\simeq\C[x_0,x_1,x_2]^{(3)}$である．
\end{theorem}

\begin{proof}
命題\ref{prop:reconstruction}の右noetherian性と追加仮定を合わせ，
Rogalski--Stafford\cite[Theorem 1.1]{RSsurface}を適用する．
有限次元の差を除いて，$A$はprojective surface $X$上の
twisted homogeneous coordinate ring $B(X,\mathcal L,\alpha)$，または
非自明なnaive blowup algebraである．後者では
$\Ext^1_{\qgr A}(\pi A,\pi A)$が無限次元になる
\cite[Corollary 1.2(II)(f)]{RSsurface}ので，Ext-finite性に反する．
従って前者となり，THCRの標準同値は
\[
 (\Cat,\OO,s)\simeq(\coh X,\OO_X,t),
 \qquad t=\alpha^*(-)\otimes\mathcal L
\]
を与える．自己同型を逆にする慣習の差は$\alpha$へ吸収する．
有限次元の差で得られた同値は構造対象とshiftを保つので，
$A_n=\Hom(\OO,s^n\OO)$という定義から，全次数の座標環もこのtripleの軌道環に一致する．

まず$X$は滑らかである．任意の閉点$x$について，local-to-global Extスペクトル系列が，局所Ext層の点支持性により退化することから
\[
 \Ext^i_X(k(x),k(x))\simeq
 \Ext^i_{\OO_{X,x}}(k(x),k(x)).
\]
補題\ref{lem:dimension}により$i>2$で左辺は零である．
局所環上の極小自由分解と正則局所環の判定
\cite[Tag 065U, Proposition 10.110.5]{Stacks}から$\OO_{X,x}$は正則であり，
$\C$上なので$X$は滑らかである．
Serre関手の一意性により$\tau=(-)\otimes\omega_X^{-1}$となる．
元のAZ生成条件とample line bundleの判定
\cite[Tag 01PR, Proposition 28.27.13]{Stacks}より$D=-K_X$はampleである．
この判定で正の指数だけを使うには，AZ生成を対象$F\otimes\omega_X^{-1}$へ適用してから
$\omega_X$をtensorすればよい．
Riemann--RochとSerre消滅より，$m\gg0$で
\[
 h^0(X,\OO_X(mD))=\chi(\OO_X)+\tfrac12D^2m(m+1).
\]
反標準Hilbert関数との比較から$D^2=9$を得る．

次に$t^3=\tau$をskyscraper sheafに適用すると$\alpha^3=\id_X$となる．
また$H=c_1(\mathcal L)$と置けば
\begin{equation}\label{eq:birational-norm}
 D=H+\alpha^*H+(\alpha^2)^*H
\end{equation}
がPicard群で成立する．$\mathcal L$は大域生成される．実際，
$A$の一次生成性により$A_n=A_{n-1}A_1$なので，$A_1=H^0(X,\mathcal L)$に基点があれば
全ての$A_n=H^0(X,\mathcal L_n)$にも同じ基点が残る．
これは$\mathcal L_{3m+r}\simeq\mathcal L_r\otimes\OO_X(mD)$と$D$のample性から，
$\mathcal L_n$が$n\gg0$で大域生成されることに反する．
従って$H$はnefである．

$a=H^2\geq0$，$b=H\cdot\alpha^*H\geq0$と置くと，$\alpha^3=1$と交点数の不変性から
\[
 3a+6b=D^2=9,\qquad D\cdot H=a+2b=3.
\]
Hodge指数定理\cite[V, Theorem 1.9]{Hartshorne}より
$9a\leq(D\cdot H)^2=9$である．$a+2b=3$と$a,b\in\Z$から$a$は奇数なので，
$a=1$，$b=1$となる．さらに$(D-3H)\cdot D=(D-3H)^2=0$であり，
同じ指数定理から$D\equiv3H$である．従ってNakai--Moishezon判定
\cite[V, Theorem 1.10]{Hartshorne}により$H$はampleとなる．
大域生成された$\mathcal L$が定める射は有限で，その像の次数と射の次数の積は$H^2=1$である．
従って像は次数1の非退化曲面，すなわち線形な$\mathbb P^2$であり，射は有限双有理である．
像の正規性により$X\simeq\mathbb P^2$，$\mathcal L\simeq\OO(1)$となる．

以上でtripleの記述を得た．標準Koszul完全列
\[
 0\longrightarrow\OO\longrightarrow\OO(1)^{\oplus3}
 \longrightarrow\OO(2)^{\oplus3}\longrightarrow\OO(3)\longrightarrow0
\]
と$H^1(\mathbb P^2,\OO(-1))=0$から定理\ref{thm:finite}(i)が成立する．
軌道環の積は$\alpha$の線形持ち上げによる多項式環のtwisted積であり，Zhang twistの記述も得る．
最後に，反標準tripleは$(\coh\mathbb P^2,\OO,(-)\otimes\OO(3))$なので，
その軌道環は可換な三次Veroneseである．
\end{proof}

\begin{remark}[分類の適用範囲]\label{rem:birational-scope}
この定理の主要な外部入力はRogalski--Staffordの分類であり，本稿の新しい一般分類を主張するものではない．
本稿ではその仮定を明示し，Ext-finite性によるnaive blowupの除外と，
固定した三乗根・偏極の識別を最後まで行った．
左noetherian性，一次生成性，双有理的に可換・幾何的であることは，基礎公理からは未証明である．
一般のSklyanin型はこの双有理的可換性を仮定した定理の対象ではない．
高VeroneseがTHCRになるという結論だけで，元の$qgr$のheartが$\coh X$に同値と推論してはならない．
研究ノートでこの点を具体的に区別する．
\end{remark}

\section{次数3のnormal境界}\label{sec:elliptic}

\subsection{一般の整な三次境界のsupportとHilbert上界}

\begin{corollary}[三次の整なnormal切断の稠密support]\label{cor:cubic-normal-dense-support}
六公理と固定したexactな三乗根$s$の下で，$0\ne g\in A_3$をnormal元とし，
$B=A/(g)$がdomainであると仮定する．
このとき
\[
 \gcd\{n>0:B_n\ne0\}=1,
 \qquad B_n\ne0\quad(n\gg0).
\]
\end{corollary}
\begin{proof}
$A$はdomainなので$g$は両側regularである．$h_n=\dim A_n$，$b_n=\dim B_n$とすると，
$0\to A(-3)\xrightarrow{g\cdot}A\to B\to0$より
\[
 b_3=h_3-h_0=10-1=9.
\]
従って$B_3\ne0$であり，domain性から$B$は無限次元である．
非零な正次数の最大公約数を$d$とすると$d\mid3$である．
元のAZ1は$(\OO,(3))$に対して成立するので，
補題\ref{lem:normal-quotient-AZ-residue}から$\gcd(d,3)=1$となり，$d=1$である．

Domain性から非零成分の次数は加法について閉じている．
最大公約数1となる有限個の非零次数を選べば，それらの生成する数値半群は
十分大きい全整数を含む．従って後半も従う．
\end{proof}


\begin{proposition}[三次の整なnormal切断のHilbert上界]\label{prop:cubic-normal-hilbert-bound}
六公理と固定したexactな三乗根の下で，$0\ne g\in A_3$がnormalで
$B=A/(g)$がdomainであるとする．$h_n=\dim A_n,b_n=\dim B_n$と置くと
\[
 b_{3m}=9m\ (m\geq1),\qquad b_n\leq3n\ (n>0),\qquad
 h_n\leq P(n)\ (n\geq0).
\]
特に$h_1\leq3,h_2\leq6$である．
さらに$B$が一次生成なら$\dim B_1=3$であり，
Artin--Staffordの定理による境界のintegral curve $Y$と線束$\mathcal L$は
\[
 \deg\mathcal L=3,\qquad p_a(Y)=1
\]
を満たす．この命題自体は曲線の数値だけを決定する．全次数の識別と$A$の認識は以下で証明する．
\end{proposition}
\begin{proof}
$g$による短完全列から$b_n=h_n-h_{n-3}$であり，反標準次数では
\[
 b_{3m}=P(3m)-P(3m-3)=9m\qquad(m\geq1)
\]
である．固定した$n>0$について$b_n>0$なら，domainのHopf不等式を三重積へ適用して
\[
 3b_n-2\leq\dim(B_n)^3\leq b_{3n}=9n
\]
を得る．整数性から$b_n\leq3n$である．$b_n=0$の場合にも同じ上界が成立する．
この上界にはsupportの稠密性を必要としない．

$b_0=1$と
$h_N=\sum_{0\leq j\leq N,\ j\equiv N\pmod3}b_j$を使う．
$N=3m$では公理により$h_N=P(N)$であり，他の二つの剰余類では
\begin{align*}
 h_{3m+1}&\leq\sum_{j=0}^m3(3j+1)
       =\frac{(3m+2)(3m+3)}2=P(3m+1),\\
 h_{3m+2}&\leq\sum_{j=0}^m3(3j+2)
       =\frac{(3m+3)(3m+4)}2=P(3m+2).
\end{align*}
$g$の次数は3なので$h_1=b_1,h_2=b_2$である．

$B$が一次生成なら$B_1^{\otimes3}\twoheadrightarrow B_3$より
$9=b_3\leq b_1^3$であり，上界$b_1\leq3$と合わせて$b_1=3$となる．
$b_n\leq3n$と$b_{3m}=9m$から$\GKdim B=2$である．
Artin--Stafford\cite[Theorem 0.2]{ASquadratic}を適用し，
$B\hookrightarrow B(Y,\mathcal L,\sigma)$と有限次元余核を得る．
$d=\deg\mathcal L$と置けば，大きい$m$についてRiemann--Rochから
\[
 9m=b_{3m}=3md+1-p_a(Y)
\]
なので，$d=3,p_a(Y)=1$である．
\end{proof}

\begin{remark}[残る二つの剰余類]
次数2の場合には反標準次数$3m$が二つの偶奇を通るため，累積和の等号が全項の等号を強制した．
次数3の場合には確定した累積和は剰余類0だけを通るので，
上の$h_n\leq P(n)$から剰余類1・2の等号を結論することはできない．
また，Serreによる剰余補題は$(e,\ell)=(3,3)$で$d\mid0$を許すため，
supportの稠密性には上のAZ1による補題を使った．
\end{remark}

\subsection{三次境界の双対性とsection欠損}

\begin{corollary}[三次の整な境界と反標準剰余類の消滅]\label{cor:cubic-boundary-duality}
元の六公理とexactな三乗根の下で，$0\ne g\in A_3$がnormalで，
$B=A/(g)$がdomainであるとし，$D=\pi_BB$と置く．
\[
 \widehat B_n=\Hom_{\qgr B}(D,D(n)),\qquad E_n=\Ext^1_{\qgr A}(\OO,\OO(n))
\]
と置くと，
\[
 \widehat B_0=\C,\qquad \widehat B_n=0\ (n<0),\qquad
 \Ext^1_{\qgr B}(D,D(n))\simeq \widehat B_{-n}^*.
\]
特に
\[
 \Ext^1_{\qgr B}(D,D)=\C,\qquad
 \Ext^1_{\qgr B}(D,D(n))=0\ (n>0),\qquad E_{3m}=0\ (m\in\Z).
\]
さらに$n>0$に対し，切断の長完全列は完全列
\[
 0\to B_n\to \widehat B_n\to E_{n-3}\to E_n\to0
\]
を与える．
\end{corollary}
\begin{proof}
Right noetherian domain $B$はOre domainなので，tailsの射
$B_{\geq q}\to B(n)$はgraded quotient division ringにおける左乗法として表される．
従って$\widehat B=\bigoplus_{n\in\Z}\widehat B_n$はgraded domainである．
完全忠実な$i_*$とambient Ext-finitenessから$\widehat B_0$は有限次元であり，
代数閉体上の有限次元domainだから$\widehat B_0=\C$である．
もし$0\ne u\in \widehat B_{-r}$，$r>0$なら，$u^3B_{3r}\subset \widehat B_0$は
$B_{3r}$と同次元である．しかし$\dim B_{3r}=9r>1$なので矛盾する．

$e=\ell=3$を命題\ref{prop:normal-boundary-duality}へ代入すると双対性と正次数の消滅を得る．
また$Li^*\OO=D$だから
\[
 \Ext^j_{\qgr A}(\OO,i_*D(n))\simeq
 \Ext^j_{\qgr B}(D,D(n)).
\]
この等式と$0\to\OO(n-3)\to\OO(n)\to i_*D(n)\to0$から，
$n>0$で表示された完全列を得る．
特に$E_{n-3}\twoheadrightarrow E_n$であり，$E_0=0$から$E_{3m}=0$ for $m\geq0$を得る．
負の$m$は$E_n^*\simeq E_{-n-3}$による．
\end{proof}


\begin{corollary}[三次境界の飽和後のHilbert上界]\label{cor:cubic-saturated-hilbert}
上の三次の整な境界の設定で，$\widehat b_n=\dim_\C \widehat B_n$と置く．
このとき
\[
 \widehat B_{3m}=B_{3m},\quad \widehat b_{3m}=9m\quad(m>0),
 \qquad b_n\leq \widehat b_n\leq3n\quad(n>0).
\]
特に$n>0$に対して
\[
 0\leq e_{n-3}-e_n=\widehat b_n-b_n\leq3n-b_n,
 \qquad e_n=\dim_\C E_n.
\]
\end{corollary}
\begin{proof}
前系の完全列と$E_{3m}=0$から，標準射$B_{3m}\to \widehat B_{3m}$は全ての$m>0$で同型である．
$\widehat B$はgraded domainなので，$\widehat b_n>0$ならHopfの次元不等式を三重積へ適用して
\[
 3\widehat b_n-2\leq\dim_\C(\widehat B_n)^3\leq \widehat b_{3n}=9n
\]
を得る．整数性から$\widehat b_n\leq3n$である．$\widehat b_n=0$の場合にもこの上界は成立する．
最後に$B_n\hookrightarrow \widehat B_n$と前系の完全列を用いる．
\end{proof}

\subsection{有限section欠損と最初の認識障害}

引き続き六公理とexactな三乗根を仮定し，$0\ne g\in A_3$はnormalで
$B=A/(g)$はdomainとする．$\widehat B$，$E_n$は系\ref{cor:cubic-boundary-duality}の記号を用い，
\[
 C=\widehat B/B,\qquad
 \delta_n=\dim_\C C_n,\qquad
 \epsilon=\dim_\C E_{-1}=\dim_\C E_{-2}
\]
と置く．$C$は正次数の剰余類$1,2$だけにsupportを持つ有限次元graded vector spaceである．
有限次元性は$E$の有限次元性と同系の完全列から従う．
ここで$C$の零性と$B$の一次生成性は別の条件である．

\begin{proposition}[Section欠損の質量とExtの復元]\label{prop:cubic-defect-mass}
上の設定で
\begin{equation}\label{eq:cubic-defect-mass}
 \sum_{m\geq0}\delta_{3m+1}
 =\sum_{m\geq0}\delta_{3m+2}
 =\epsilon,
 \qquad \dim_\C C=2\epsilon.
\end{equation}
また$r\in\{1,2\}$と$m\geq0$に対して
\begin{equation}\label{eq:cubic-E-tail}
 e_{3m+r}=\sum_{j>m}\delta_{3j+r},\qquad e_n=\dim_\C E_n.
\end{equation}
従って正次数の$E$の次元はsection欠損の末尾和から一意に復元され，
負次数は$e_n=e_{-n-3}$により決まる．
特に次の四条件は同値である：
\[
 E_{-1}=0;\qquad E=0;\qquad B=\widehat B;\qquad
 \text{剰余類$1$または$2$の一方で全section欠損が零}.
\]
\end{proposition}
\begin{proof}
系\ref{cor:cubic-boundary-duality}の完全列より，$n>0$で
$\delta_n=e_{n-3}-e_n$である．従って$r=1,2$について
\[
 \sum_{j=0}^{M}\delta_{3j+r}
 =e_{r-3}-e_{3M+r}.
\]
$E$は有限次元なので，$M$を十分大きくすると最後の項は零となる．
Serre双対性から$e_{-2}=e_{-1}=\epsilon$であるため，式\eqref{eq:cubic-defect-mass}を得る．
同じtelescopingを$j=m+1$から始めれば式\eqref{eq:cubic-E-tail}となる．

$E_{-1}=0$なら両剰余類の欠損は非負項の和が零なので全て零であり，
式\eqref{eq:cubic-E-tail}から正次数の$E$も零である．
反標準剰余類の$E_{3m}$は既に零であり，残る負次数もSerre双対性により零となる．
逆向きと他の同値は式\eqref{eq:cubic-defect-mass}から従う．
次数$0$と負次数では$B$と$\widehat B$は既に一致していることを用いた．
\end{proof}

\begin{proposition}[欠損多項式とExtのLaurent多項式]\label{prop:cubic-defect-polynomial}
\[
 c(t)=\sum_{n>0}\delta_nt^n,
 \qquad \mathcal E(t)=\sum_{n\in\Z}e_nt^n,
 \qquad \mathcal E_+(t)=\sum_{n>0}e_nt^n
\]
と置くと，いずれも有限和であり，
\begin{align}
 c(t)&=\epsilon(t+t^2)+(t^3-1)\mathcal E_+(t),\label{eq:cubic-defect-plus}\\
 (t^3-1)\mathcal E(t)&=c(t)-c(t^{-1}),\label{eq:cubic-defect-duality}\\
 \dim_\C E&=\frac23\sum_{n>0}n\delta_n=\frac23c'(1).
 \label{eq:cubic-defect-moment}
\end{align}
\end{proposition}
\begin{proof}
$n>0$の係数に$\delta_n=e_{n-3}-e_n$を用いると，
正次数への最初の流入は$e_{-2}t+e_{-1}t^2=\epsilon(t+t^2)$であるから，
式\eqref{eq:cubic-defect-plus}を得る．

$\mathcal E(t)$の$(t^3-1)$倍の$t^n$係数は$d_n=e_{n-3}-e_n$である．
Serre対称性から
\[
 d_{-n}=e_{-n-3}-e_{-n}=e_n-e_{n-3}=-d_n,
 \qquad d_0=e_{-3}-e_0=0.
\]
正次数では$d_n=\delta_n$なので，式\eqref{eq:cubic-defect-duality}となる．
このLaurent多項式恒等式を微分して$t=1$を代入すれば
$3\mathcal E(1)=2c'(1)$を得る．
\end{proof}

\begin{corollary}[飽和後のHilbert列が標準の場合のEuler列]\label{cor:cubic-standard-saturation-euler}
さらに$\widehat b_n=3n$ for every $n>0$と仮定する．
このとき全$n\in\Z$に対して
\begin{equation}\label{eq:cubic-period-three-euler}
 q_n=\chi(\OO,\OO(n))
 =P(n)-\epsilon\,\mathbf1_{3\nmid n}
\end{equation}
である．特に$n\geq0$では
\[
 h_n=P(n)-\epsilon\,\mathbf1_{3\nmid n}+e_n.
\]
この数値計算は標準section列を条件として行う．その条件が一般の三次domain境界で自動的に成立することを，
後の定理\ref{thm:v8-cubic-standard-sections}で証明する．通常の曲線上のTHCR表示はここでは必要ない．
\end{corollary}
\begin{proof}
切断完全列と$h_n=0$ for $n<0$より
\[
 h_N=\sum_{\substack{0\leq j\leq N\\j\equiv N\pmod3}}b_j.
\]
$r=1,2$と$N=3m+r$について$\widehat b_j=3j$を用いると
\[
 h_{3m+r}
 =P(3m+r)-\sum_{j=0}^m\delta_{3j+r}
 =P(3m+r)-\epsilon+e_{3m+r}.
\]
$n\equiv0\pmod3$では既知の$h_n=P(n)$と$e_n=0$を用いる．
$n\geq0$で$\Ext^2(\OO,\OO(n))=0$だから$q_n=h_n-e_n$である．
$n=-1,-2$では負次数Homと$\Ext^2$がともに零であり，$q_n=-\epsilon$となる．
$n\leq-3$には$q_n=q_{-n-3}$と$P(n)=P(-n-3)$を用いる．
\end{proof}

\begin{corollary}[最初の障害とEuler orbit rank]\label{cor:cubic-defect-Euler-rank}
前系の追加仮定の下で，軌道$\{[\OO(n)]:n\in\Z\}$に張られる空間上の
Euler形式のrankは，$\epsilon=0$なら$3$，$\epsilon>0$なら$5$である．
後者の最小漸化式の多項式は
\[
 (z-1)^3(z^2+z+1)
\]
である．従ってこの範囲では，$B$の飽和性，$E_{-1}=0$，Euler orbit rank $\leq4$は同値となる．
\end{corollary}
\begin{proof}
式\eqref{eq:cubic-period-three-euler}は，二次多項式と周期$3$の列の和である．
$\epsilon>0$のとき原始三乗根$\zeta,\zeta^2$の両成分が非零なので，
最小漸化式は表示された次数$5$の多項式である．
より直接には，$0\leq i,j\leq4$のEuler行列$(q_{j-i})$の行列式は
\[
 \det\begin{pmatrix}
 1&3-\epsilon&6-\epsilon&10&15-\epsilon\\
 -\epsilon&1&3-\epsilon&6-\epsilon&10\\
 -\epsilon&-\epsilon&1&3-\epsilon&6-\epsilon\\
 1&-\epsilon&-\epsilon&1&3-\epsilon\\
 3-\epsilon&1&-\epsilon&-\epsilon&1
 \end{pmatrix}
 =243\epsilon^2.
\]
漸化式からrankは高々$5$であり，この行列式から少なくとも$5$である．
$\epsilon=0$なら二次列なのでrankは高々$3$であり，
$0\leq i,j\leq2$のEuler行列は対角成分が$1$の上三角行列なのでrankは$3$である．
\end{proof}

\subsection{Veroneseへの有限性と\texorpdfstring{$\chi_1$}{chi1}の継承}

\begin{lemma}[有限巡回次数による有限拡大]\label{lem:v8-veronese-finite}
$R$をlocally finiteな$\mathbb N$次数付き代数とし，$d>0$，$S=R^{(d)}$とする．
$R$が右noetherianなら，$S$も右noetherianであり，$R$は有限生成右$S$加群である．
左側についても同じ主張が成立する．
\end{lemma}
\begin{proof}
ここでは$S$の次数を一時的に取り直さず，$R$を$\Z/d\Z$次数付き環と見て
$R=\bigoplus_{r\in\Z/d\Z}R_{[r]}$，$S=R_{[0]}$と書く．
次数$0$への射影$p_0:R\to S$は$S$両側線形である．右イデアル$I\subset S$について
$p_0(IR)=I$だから，$IR\cap S=I$となる．従って$S$の右イデアルの真の昇鎖は
$R$の右イデアルの真の昇鎖を与え，$S$は右noetherianである．

各$r$について，右イデアル$R_{[r]}R$は有限個の元$x_1,\ldots,x_t\in R_{[r]}$で生成される．
$x\in R_{[r]}$を$x=\sum_i x_i a_i$と表し，剰余次数$r$の成分を取ると
$x=\sum_i x_i p_0(a_i)$となる．従って$R_{[r]}=\sum_i x_iS$であり，
有限個の剰余類を合わせて$R_S$の有限生成性を得る．左側は反対環へ適用する．
\end{proof}

\begin{lemma}[$\chi$条件のVeroneseへの継承]\label{lem:v8-chi-veronese}
$R$をconnected locally finiteな両側noetherian次数付き代数とする．
$R$がright $\chi_j$を満たせば，任意の$d>0$について$S=R^{(d)}$もright $\chi_j$を満たす．
特にright $\chi_1$は全てのVeroneseへ継承する．
\end{lemma}
\begin{proof}
前補題より$S$は両側noetherianで，$R$は左右有限$S$加群である．
$M$を有限生成graded右$S$加群とし，次数を$d$倍に膨らませて
$N=M\otimes_S R$をgraded右$R$加群とする．$N$は有限生成である．
$R=S\oplus\bigoplus_{r\ne0}R_{[r]}$は$S$両側加群の分解だから，
$M=M\otimes_SS$は$N$を$S$へ制限した加群の直和因子である．

$R$の有限生成自由左$S$加群による分解を用いると，
全ての$q\ge0$について$\Tor_q^S(\C,R)$は有限次元である．
これらはgraded右$R$加群なので，有限個の$\C$の次数shiftを組成因子に持つ．
環の変更のスペクトル系列
\[
 \underline\Ext_R^p\bigl(\Tor_q^S(\C,R),N\bigr)
 \ \Longrightarrow\
 \underline\Ext_S^{p+q}(\C,N)
\]
を用いる．ここでは$N$を$S$へ制限し，全体を元の$\Z$次数で扱っている．
このスペクトル系列は$\C$の右$S$自由分解と$N$の右$R$injective分解の二重複体から得られ，
$p,q\ge0$の第一象限にある．$R$のright $\chi_j$と有限長加群に関する帰納法から，
$p\le j$の左辺は有限次元である．従って$i\le j$で
$\underline\Ext_S^i(\C,N)$は有限次元となる．
その剰余次数$0$の直和因子$\underline\Ext_S^i(\C,M)$も有限次元であり，
最後に次数を$d$で割れば主張を得る．
\end{proof}

\begin{proposition}[整な三次境界の両側noetherian性と高Veronese]
\label{prop:v8-cubic-high-veronese}
六公理，exactな三乗根，非零normal元$g\in A_3$，domain商$B=A/(g)$を仮定する．
このとき$B$は両側noetherianであり，ある$3$の倍数$N>0$について
\[
 B^{(N)}\simeq B(C,\mathcal L,\sigma),\qquad
 p_a(C)=1,\quad\deg\mathcal L=3N
\]
という全次数での同型が存在する．ここで$C$はintegral projective curveであり，
右辺は全sectionを取ったTHCRで，次数$1$で生成される．
\end{proposition}
\begin{proof}
$B$はconnected locally finiteな右noetherian domainで，$\GKdim B=2$である．
Artin--Stafford\cite[Theorem 0.1]{ASquadratic}
（共著者による再掲\cite[Theorem 5.1]{SVdBsurvey}）により，
\[
 Q_{\mathrm{gr}}(B)=K[t,t^{-1};\rho]
\]
であり，$K$は一変数関数体である．次数supportの最大公約数は$1$なので$\deg t=1$とする．
$\rho$が無限位数なら，同分類\cite[Theorem 0.4]{ASquadratic}
（\cite[Theorem 5.5(i)]{SVdBsurvey}）を$B$と$B^{\mathrm{op}}$へ適用して，
$B$は両側noetherianである．

$\rho$の位数が有限$d$なら，$S=B^{(d)}$は可換である．
補題\ref{lem:v8-veronese-finite}から$S$はnoetherianである．
非零な剰余成分$B_{[r]}$に$x\ne0$を選び，$a=x^{d-1}\in B_{[-r]}$と置く．
右乗法$y\mapsto ya$は，domain性から左$S$加群の単射
$B_{[r]}\hookrightarrow S$を与える．従って$B_{[r]}$は有限生成左$S$加群である．
これを全剰余類について合わせると${}_SB$はnoetherianであり，特に$B$は左noetherianである．
よっていずれの場合も$B$は両側noetherianとなる．

$B$の有限生成右加群$M$について，Yoneda拡大の忘却は
\[
 \underline\Ext_B^1(\C,M)\hookrightarrow
 \underline\Ext_A^1(\C,M)
\]
を与える．$B$線形な拡大の同型を与える$A$線形写像は$B$線形だからである．
$\underline\Hom$も同様なので，$A$のright $\chi_1$から$B$のright $\chi_1$が従う．
補題\ref{lem:v8-chi-veronese}により全ての$B^{(e)}$もright $\chi_1$を満たす．

どのVeroneseも一次生成でないと仮定する．$\rho$が有限位数なら
Artin--Staffordの上記分類の有限位数の場合により$B$がnoetherianであることに反する．
$\rho$が無限位数なら，同分類のTheorem 0.5
（\cite[Theorem 5.6]{SVdBsurvey}）は，あるVeroneseがright $\chi_1$を満たさないことを与え，
やはり矛盾する．従ってある$e>0$で$B^{(e)}$は一次生成である．
$N=3e$と取り直しても一次生成性は保たれる．

$R=B^{(N)}$と置く．全$m>0$で$\dim R_m=3Nm$である．
Artin--Staffordの一次生成の場合の定理\cite[Theorem 0.2]{ASquadratic}により，
$R\hookrightarrow B(C,\mathcal L,\sigma)$という有限次元余核を持つgraded単射を得る．
大きい$m$のRiemann--Rochから
\[
 3Nm=m\deg\mathcal L+1-p_a(C)
\]
なので$\deg\mathcal L=3N$，$p_a(C)=1$となる．
算術種数$1$のintegral curveはGorensteinで$\omega_C\simeq\OO_C$である
（補題\ref{lem:integral-genus-one-sections}の証明を参照）．
従って正次数の線束の$H^1$は零であり，全$m>0$で
$\dim B(C,\mathcal L,\sigma)_m=3Nm$となる．
上の単射の各次数の次元が等しいので，全次数で同型である．
\end{proof}


\begin{lemma}[正次数normal元によるnoetherian性の持上げ]
\label{lem:v8-normal-noetherian-lift}
$A$を非負次数付き代数，$g\in A_e$を正次数のregular normal元とし，$B=A/(g)$とする．
$B$が左noetherianなら$A$も左noetherianである．右側も同様である．
\end{lemma}
\begin{proof}
$ga=\nu(a)g$と書き，$F^mA=g^mA=Ag^m$とする．$g$のregular性から
\[
 \operatorname{gr}_F A\simeq B[z;\overline\nu],\qquad
 zb=\overline\nu(b)z,\qquad \overline a z^m\longmapsto[ag^m]
\]
となる．これは$g$-進次数と元の全次数を持つ二重次数付き同型で，$z$の全次数は$e$である．
$\overline\nu$は自己同型なので，skew polynomial環のHilbert基底定理により右辺は左noetherianである．

$J\subset A$をgraded左idealとする．その$g$-進初期idealの有限個の二重斉次生成元を，
斉次元$j_1,\ldots,j_t\in J$へ持ち上げる．任意の全次数$D$の斉次元$j\in J$について，
初期項の表示の係数を同じ二重次数で持ち上げれば，
$j-\sum_i a_i j_i$の全次数は$D$のままで$g$-進位数が真に増える．
$F^mA_D=0$ for $me>D$なので，この除去は有限回で停止する．
従って$J=\sum_iAj_i$であり，全てのgraded左idealは有限生成である．

一般の左ideal $I$には最高全次数による初期idealを取る．これはgraded左idealなので，
有限個の初期生成元を$I$の元へ持ち上げることができる．
最高全次数を順次下げる除去により，それらが$I$を生成する．従って$A$は左noetherianである．
\end{proof}

\begin{corollary}[整三次normal境界を持つambientの両側noetherian性]
\label{cor:v8-cubic-ambient-noetherian}
六公理とexactな三乗根の下で，$0\ne g\in A_3$がnormal，$A/(g)$がdomainなら，
$A$は両側noetherianである．特に$A$は各Veronese上で左右有限生成である．
\end{corollary}
\begin{proof}
命題\ref{prop:v8-cubic-high-veronese}により$B=A/(g)$は両側noetherianであり，
$A$はdomainなので$g$はregularである．前補題を適用する．
最後の主張は補題\ref{lem:v8-veronese-finite}による．
\end{proof}

\begin{remark}[高Veroneseと元のheart]
命題\ref{prop:v8-cubic-high-veronese}は，元の$\qgr B$と$\coh C$の同値を主張しない．
Veronese直像が非忠実な場合，元のheartに余分な単純対象が残り得る．
以下では圏同値を用いず，構造対象からのcohomologyだけを比較する．
\end{remark}

\subsection{Veroneseによる構造対象のcohomologyの比較}

\begin{lemma}[有限Veroneseとcohomology]\label{lem:v8-veronese-cohomology}
$R$をconnected locally finiteな両側noetherian次数付き代数とし，$S=R^{(d)}$とする．
Veronese直像を
\[
 V:\QGr R\longrightarrow\QGr S,
 \qquad V(\pi_RM)=\pi_S(M^{(d)})
\]
と書く．この関手はexactであり，全ての$F\in\QGr R$と$i\ge0$について
自然な同型
\[
 \Ext^i_{\QGr R}(\pi_RR,F)
 \simeq
 \Ext^i_{\QGr S}(\pi_SS,VF)
\]
が成立する．この主張に$V$の忠実性は必要ない．
\end{lemma}
\begin{proof}
Veroneseはgraded加群上でexactで，torsion加群をtorsion加群へ送るので，
商圏上でもexactである．両側の有限生成性は補題\ref{lem:v8-veronese-finite}による．
以下，quotient functorの右随伴を$\omega_R,\omega_S$と書く．

$J$を$\QGr R$のinjective対象とし，$I=\omega_RJ$と置く．
$\pi_R$はexactなので$I$はgraded右$R$加群としてinjectiveである．
さらに$I$は$R$-torsion-freeであり，$R$-saturatedである．
$S$の次数は$d$倍して扱い，$X$をgraded右$S$加群とすると，
gradedな導来随伴と$I$のinjectivityから
\[
 \Ext^q_{\Gr S}(X,I^{(d)})
 \simeq
 \Hom_{\Gr R}\bigl(\Tor_q^S(X,R),I\bigr).
\]
ここで左辺・右辺とも次数$0$の射を用い，誘導加群の次数は$d$倍にする．

$X$が有限次元なら，全ての$q$について$\Tor_q^S(X,R)$は有限次元である．
実際，${}_SR$は有限生成で$S$は左noetherianなので，有限生成自由左$S$加群の分解を取れる．
$I$の$R$-torsion-free性から，上の右辺は全て零である．
特に$X=(S/S_{\ge\ell})(j)$とし，
$0\to S_{\ge\ell}\to S\to S/S_{\ge\ell}\to0$へ$\Hom(-,I^{(d)})$を適用すると，
全ての次数shiftを含めて
\[
 \underline\Hom_S(S,I^{(d)})
 \simeq\underline\Hom_S(S_{\ge\ell},I^{(d)})
\]
となる．section functorのtruncation公式から$I^{(d)}$は$S$-saturatedである．

次に$X=S_{\ge\ell}$とする．$q>0$では
\[
 \Tor_q^S(S_{\ge\ell},R)
 \simeq \Tor_{q+1}^S(S/S_{\ge\ell},R)
\]
が有限次元なので，$\Ext^q_{\Gr S}(S_{\ge\ell},I^{(d)})=0$である．
構造対象のcohomologyをtruncationからのExtの直極限で表す公式
\[
 \Ext^i_{\QGr S}(\pi_SS,\pi_SM)
 =\varinjlim_{\ell}
   \Ext^i_{\Gr S}(S_{\ge\ell},M)
\]
を用いる．これは\cite[\S4.1, (4.1), Lemmas 4.1.4--4.1.6]{BvdBGenerators}
に証明されている．同文献は左加群の規約なので，ここでは反対環へ適用する．
同節の条件$\dim\underline\Ext^i_S(\C,\C)<\infty$はnoetherian性から得られる
有限生成自由加群の分解により満たされ，この公式に有限cohomological dimensionの仮定は不要である．
以上から$\pi_S I^{(d)}=VJ$は
$\Hom_{\QGr S}(\pi_SS,-)$についてacyclicである．
またsaturationから
\[
 \Hom_{\QGr S}(\pi_SS,VJ)
 =(I^{(d)})_0=I_0
 =\Hom_{\QGr R}(\pi_RR,J).
\]
一般の$F$のinjective分解へexactな$V$を適用すると，右側のcohomologyを計算する
acyclic分解になる．次数$0$の上の自然同型を分解全体へ適用すれば結論を得る．
なお，noetherian対象間のExtは小さい$\qgr$と大きい$\QGr$で一致する．
これはYoneda拡大とその比較図式を，必要な有限個の元を含むnoetherian部分対象へ制限することによる．
\end{proof}

\begin{theorem}[整な三次境界の標準section Hilbert関数]
\label{thm:v8-cubic-standard-sections}
六公理，exactな三乗根，非零normal元$g\in A_3$，domain商$B=A/(g)$を仮定する．
このときEuler orbit rankの有限性を追加せず，全ての$n>0$について
\[
 \dim\widehat B_n=3n
\]
が成立する．より正確には，全整数$n$について
$\chi_{\qgr B}(D,D(n))=3n$である．
\end{theorem}
\begin{proof}
命題\ref{prop:v8-cubic-high-veronese}により，$3\mid N$と
$R=B^{(N)}=B(C,\mathcal L,\sigma)$，$p_a(C)=1$，$\deg\mathcal L=3N$を取る．
$0\le r<N$について
\[
 M_r=\bigoplus_{m\in\Z}B_{Nm+r}
\]
をgraded右$R$加群とし，$\qgr R\simeq\coh C$の下で$\pi_RM_r$に対応する層を$F_r$と書く．
$M_r$は補題\ref{lem:v8-veronese-finite}により有限生成である．
$B$の正次数supportは余有限なので，十分大きい$k$を選んで
$0\ne a\in B_{Nk-r}$を取れる．左乗法による写像
\[
 M_r\longrightarrow R(k),\qquad y\longmapsto ay
\]
はgraded右$R$線形であり，$B$のdomain性から単射である．
$M_r$は有限次元でないから$F_r\ne0$である．
従って$F_r$はintegral curve上の線束に含まれる非零coherent sheafであり，
rank $1$のtorsion-free sheafとなる．

$n=Nm+r$を任意の整数とする．前補題と次数shiftの整合性により
\[
 \chi_{\qgr B}(D,D(n))
 =\chi\bigl(C,T_R^mF_r\bigr)
 =3Nm+\deg F_r,
\]
ここで$T_R$は$R$の次数shiftに対応するautoequivalenceであり，
線束のtensorと曲線のautomorphismからなる．従ってrankを保ち，
rank $1$の層のdegreeを$3N$増加させる．
最後のRiemann--Rochは負の$m$も含む．
一方，境界双対性から
\[
 \chi_{\qgr B}(D,D(n))=
 \begin{cases}
 \widehat b_n,&n>0,\\
 0,&n=0,\\
 -\widehat b_{-n},&n<0.
 \end{cases}
\]
正次数の上界$\widehat b_n\le3n$を$m\gg0$へ適用すると$\deg F_r\le3r$である．
同じ上界を$-n>0$へ適用し，$m\ll0$で比較すると
$3Nm+\deg F_r=-\widehat b_{-n}\ge3n=3Nm+3r$となる．
従って$\deg F_r=3r$であり，全整数で境界Euler列は$3n$である．
特に全$n>0$で$\widehat b_n=3n$となる．
\end{proof}

\begin{corollary}[三次domain境界でのEuler列と有限rank]
\label{cor:v8-cubic-automatic-euler-rank}
定理\ref{thm:v8-cubic-standard-sections}の仮定の下で，
$\epsilon=\dim E_{-1}=\dim E_{-2}$と置く．全整数$n$について
\[
 q_n=P(n)-\epsilon\mathbf1_{3\nmid n}
\]
である．特にEuler orbit rankは自動的に有限であり，
$\epsilon=0$なら$3$，$\epsilon>0$なら$5$となる．
\end{corollary}
\begin{proof}
定理\ref{thm:v8-cubic-standard-sections}により，
系\ref{cor:cubic-standard-saturation-euler}および
系\ref{cor:cubic-defect-Euler-rank}の追加条件が満たされる．
\end{proof}

\begin{remark}[なお残る障害]
この定理は$\epsilon=0$を証明しない．有限section欠損$\widehat B/B$の総次元は
$2\epsilon$であり，その消滅と$B$の一次生成性は引き続き別の問題である．
また，高Veroneseの曲線$C$から元のheartを復元したという主張も含まない．
\end{remark}


\subsection{滑らかな高Veroneseから元の境界heartを復元する}

ここでは右加群の規約を固定する．関数体 $K=k(C)$ の自己同型を
曲線自己同型のpullbackと同じ記号 $\sigma$ で表し，
\[
 (f t^n)(h t^m)=f\sigma^n(h)t^{n+m}
\]
とする．可逆層を $K$ の部分層として実現するとき，その積と商も
$K$ 内で取る．このような部分層を以下では分数可逆層と呼ぶ．
THCRの積とSerre同値は\cite[\S2.3, Proposition 2.2, Theorem 2.4]{CKSTHCR}で確認できる．
同文献は左加群を使うため，以下では右加群への規約変換を明示する．

\begin{lemma}[余有限な可逆層cocycleの一次延長]
\label{lem:v10-cofinite-line-cocycle}
$C$ をintegral projective curve，$\sigma\in\operatorname{Aut}(C)$ とする．
ある $n_0>0$ 以上の全整数 $n$ に分数可逆層 $G_n\subset K$ が与えられ，
\[
 G_{n+m}=G_n\,\sigma^{n*}G_m\qquad(n,m\ge n_0)
\]
とする．このとき分数可逆層 $L\subset K$ が一意に存在して，
\[
 G_n=L_n:=L\,\sigma^*L\cdots\sigma^{(n-1)*}L
 \qquad(n\ge n_0)
\]
となる．さらに $\deg G_n=dn$ なら $\deg L=d$ である．
\end{lemma}
\begin{proof}
$a\ge n_0$ を固定し，
\[
 L=G_{a+1}(\sigma^*G_a)^{-1}
\]
と置く．$G_{n+a+1}$ を $(n,a+1)$ と $(n+1,a)$ に分けると，
\[
 G_n\sigma^{n*}G_{a+1}
 =G_{n+1}\sigma^{(n+1)*}G_a
\]
なので，可逆層を消去して
$G_{n+1}=G_n\sigma^{n*}L$ を得る．
従って $H=G_{n_0}L_{n_0}^{-1}$ と置けば，全ての $n\ge n_0$ で
$G_n=HL_n$ である．元のcocycleへ代入すると
\[
 H L_{n+m}=H\sigma^{n*}H L_{n+m}
\]
となり，$\sigma^{n*}H=\mathcal O_C$，従って $H=\mathcal O_C$ を得る．
一意性は $L=\sigma^{-n*}(G_{n+1}G_n^{-1})$ による．
次数は最初の $L$ の式から直ちに従う．
\end{proof}

\begin{proposition}[可逆な剰余層によるheart復元判定]
\label{prop:v10-invertible-residue-reconstruction}
$B$ をconnected locally finiteなnoetherian次数付きdomainとし，
\[
 Q_{\mathrm{gr}}(B)=K[t,t^{-1};\sigma],\qquad \deg t=1
\]
とする．ある $N>0$ について，$R=B^{(N)}$ が全section環
\[
 R=B(C,\mathcal M,\rho),\qquad \rho=\sigma^N,
 \qquad\deg\mathcal M=3N
\]
であると仮定する．ここで $C$ は算術種数 $1$ のintegral projective curveであり，
分数環の同一視はこの表示と整合するものとする．さらに
\begin{enumerate}
\item $K$ 上の $\sigma$ が $C$ の自己同型へ延長する．
\item $M_r=\bigoplus_m B_{Nm+r}$ に対応する
      $F_r\in\coh C$ は，$0\le r<N$ の全てについて可逆層であり，
      $\deg F_r=3r$ である．
\end{enumerate}
このとき次数 $3$ の可逆層 $L$ と，有限次元余核を持つgraded単射
\[
 B\hookrightarrow T:=B(C,L,\sigma)
\]
が存在する．特に元の次数shiftを含めて
\[
 (\qgr B,\pi B,(1))
 \simeq
 \bigl(\coh C,\mathcal O_C,
       \sigma^{-1*}((- )\otimes L)\bigr)
\]
となり，$\widehat B\simeq T$ である．
\end{proposition}
\begin{proof}
$B_n=V_nt^n$ と書く．$\mathcal M$ を分数可逆層として実現し，
\[
 R_m=H^0(C,\mathcal M_m)t^{Nm},\qquad
 \mathcal M_m=\mathcal M\rho^*\mathcal M\cdots\rho^{(m-1)*}\mathcal M
\]
とする．$R$ の分数環の次数 $1$ 元が $ut^N$ で与えられる場合も，
分数可逆層 $\mathcal M$ を $u\mathcal M$ に取り直すことでこの形になる．

各 $r$ について，$M_r$ の元を
\[
 V_{Nm+r}t^{Nm+r}\longrightarrow
       \sigma^{-r}(V_{Nm+r})t^{Nm}
\]
へ送る．これは右 $R$ 加群の同型である．実際，$h\in R_\ell$ の係数を同じく
$h\in K$ と書けば，
\[
 \sigma^{-r}\bigl(v\sigma^{Nm+r}(h)\bigr)
 =\sigma^{-r}(v)\sigma^{Nm}(h)
\]
である．THCRの右加群に対するSerre対応では，$F\in\coh C$ のsection加群を
\[
 \Gamma_*(F)_m=H^0(C,F\otimes\mathcal M_m),\qquad
 f\cdot h=f\rho^m(h)
\]
と取れる．有限生成加群とそのsection回復は十分高い次数で一致する．
従って，各 $F_r$ のgeneric fiberを上の分数加群と整合するように自明化すると，
十分大きい $m$ について実際の $K$ の部分空間として
\[
 \sigma^{-r}(V_{Nm+r})=H^0(C,F_r\mathcal M_m)
\]
となる．ここで右 $R$ 加群 $M_r$ の有限生成性は
補題\ref{lem:v8-veronese-finite}による．

従って，全ての十分大きい $n=Nm+r$ に対し，分数可逆層
\[
 G_{Nm+r}=\sigma^{r*}(F_r\mathcal M_m)
\]
を取れば
\[
 V_n=H^0(C,G_n),\qquad \deg G_n=3n
\]
となる．有限個の剰余類しかないので，一つの $n_0$ を全てに使える．
$n_0$ を増やして，各 $G_n$ が大域切断で生成されるようにする．
算術種数 $1$ のintegral curveでは，十分大きい次数の可逆層は大域生成される．

$B$ の積の閉性と大域生成性から，$n,m\ge n_0$ に対し
\[
 G_n\sigma^{n*}G_m\subseteq G_{n+m}
\]
を得る．両辺は同じ次数 $3(n+m)$ の可逆層である．従ってその余核は
長さ $0$ のtorsion層であり，この包含は等号である．
補題\ref{lem:v10-cofinite-line-cocycle}により，次数 $3$ の分数可逆層 $L$ が存在して
$G_n=L_n$ for $n\ge n_0$ となる．

低次数についても $B_n\subset H^0(C,L_n)t^n$ を示す．
$b\in V_n$，$n\ge0$ を固定し，$m$ を十分大きく取る．積の閉性から
\[
 b\,\sigma^n H^0(C,L_m)\subseteq H^0(C,L_{n+m})
\]
である．$\sigma^{n*}L_m$ は大域生成されるので，分数可逆層の包含
\[
 b\,\sigma^{n*}L_m\subseteq L_n\sigma^{n*}L_m
\]
を得る．可逆層を消去すると $b\mathcal O_C\subseteq L_n$，すなわち
$b\in H^0(C,L_n)$ となる．これで全次数のgraded単射 $B\hookrightarrow T$ と
高次数での等号が得られた．

有限余次元の変更がtailsを変えないことをここで直接確認する．
十分大きい $q$ に対し $I=B_{\ge q}=T_{\ge q}$ は両環の共通両側idealであり，
$T$ は左右有限 $B$ 加群である．restrictionと $-\otimes_B T$ を考える．
有限次元加群は有限次元へ送られる．また有限生成 $M$ に対し，
$\operatorname{Tor}_1^B(M,T)$ は有限次元加群 $\operatorname{Tor}_1^B(M,T/B)$ に含まれるので，
$-\otimes_B T$ はtails上でexactになり，両関手はtails上の関手を定める．
ここで後者のTorの有限次元性には，right noetherian性による有限生成自由加群からなる分解と，$T/B$の有限次元性を用いる．
unit $\eta_M:M\to M\otimes_B T$ の余核は $I$ で消える．
その核については，各 $i\in I$ に対し
\[
 M\otimes_B T\longrightarrow M,\qquad m\otimes t\longmapsto m(ti)
\]
が定義でき，$\eta_M(m)=0$ なら $mi=0$ となる．
またcounit $\varepsilon_P:P_B\otimes_B T\to P$ は全射である．
$\sum_jp_j\otimes t_j$ がその核にあれば，任意の $i\in I$ に対し
\[
 \sum_jp_j\otimes t_ji
 =\left(\sum_jp_jt_ji\right)\otimes1=0
\]
なので，その核も $I$ で消える．有限生成なこれらの核・余核は有限次元であり，
unitとcounitはtails上で同型となる．この同値は $\pi B$ と次数shiftを保つ．

最後にTHCRの通常の同値を合成する．上記の右加群規約では
\[
 \Gamma_*(F)(1)_m=H^0(C,F L_{m+1})
 \xrightarrow{\ \sigma^{-1}\ }
 H^0\bigl(C,\sigma^{-1*}(F L)L_m\bigr)
\]
だから，次数shiftは $F\mapsto\sigma^{-1*}(F\otimes L)$ に対応する．
これが主張したtriple同値であり，構造対象のsection環は $T$ である．
\end{proof}

\begin{lemma}[一つの原始剰余類の可逆性で十分]
\label{lem:v10-one-invertible-residue}
命題\ref{prop:v10-invertible-residue-reconstruction}において，$F_r$ は全て
rank $1$ torsion-freeで $\deg F_r=3r$ と仮定する．
条件(2)の代わりに，$\gcd(r,N)=1$ を満たす一つの $r$ について
$F_r$ が可逆と仮定すれば十分である．特に $F_1$ の可逆性で足りる．
\end{lemma}
\begin{proof}
前命題の証明の前半は，$F_r$ をtorsion-freeとしても成立し，
十分高い全次数に分数torsion-free層 $G_n$ を与える．
各 $G_n$ の次数は $3n$ であり，十分高いところでは大域生成される．
これは，各剰余類で固定された $F_r$ に正次数の $\rho$-ampleな可逆層列 $\mathcal M_m$ をtensorし，automorphismでpullbackする構成による．
可逆な剰余類に属する十分大きい $n$ を一つ固定すると，$G_n$ は可逆である．
$G_{jn}$ が可逆なら，積は可逆層からtorsion-free層への非零射
\[
 G_n\sigma^{n*}G_{jn}\longrightarrow G_{(j+1)n}
\]
を与える．これは単射であり，両辺の次数は等しいので同型である．
従って帰納法により全ての $G_{jn}$ は可逆となる．
$n\equiv r\pmod N$ が巡回群を生成するため，これらは全剰余類を通る．
各 $G_{jn}$ は対応する $F_s$ のautomorphism pullbackと可逆層のtensorなので，
全ての $F_s$ が可逆となる．
\end{proof}

\begin{theorem}[滑らかな高Veroneseによる元の境界heartの復元]
\label{thm:v10-smooth-veronese-heart}
六公理，exactな三乗根，非零normal元 $g\in A_3$，domain商 $B=A/(g)$ を仮定する．
命題\ref{prop:v8-cubic-high-veronese}の高Veronese表示の曲線 $C$ が滑らかなら，
次数 $3$ の可逆層と曲線自己同型による通常の曲線triple表示
\[
 (\qgr B,D,(1))\simeq
 (\coh C,\mathcal O_C,\mathcal L\otimes\alpha^*(-))
 \qquad(\deg\mathcal L=3)
\]
が成立する．具体的には前命題の記号で
$\alpha=\sigma^{-1}$，$\mathcal L=\sigma^{-1*}L$ と取れる．
従って，系\ref{cor:v8-geometric-one-obstruction}の同値条件は，この滑らかさの追加仮定の下で全て適用できる．
\end{theorem}
\begin{proof}
滑らかな射影曲線の関数体自己同型は曲線自己同型へ一意に延長する．
また滑らかな曲線上のrank $1$ torsion-free coherent sheafは可逆層である．
定理\ref{thm:v8-cubic-standard-sections}の証明で得た
$\deg F_r=3r$ と合わせ，前命題を適用する．
\end{proof}

\begin{corollary}[Veroneseの核の消滅と有限巡回gradingのdensity]
\label{cor:v10-smooth-veronese-dense}
前定理，あるいは命題\ref{prop:v10-invertible-residue-reconstruction}の仮定の下で，
元のVeronese直像
\[
 V:\qgr B\longrightarrow\qgr B^{(N)}
\]
は同値である．特に $M^{(N)}$ が有限次元なら，有限生成graded右 $B$ 加群 $M$ 自体も有限次元である．
また $B_{[r]}=\bigoplus_{n\equiv r\ (N)}B_n$ と置くと，全ての $r\in\mathbb Z/N\mathbb Z$ について
\[
 B_{[r]}B_{[-r]}\subseteq B_{[0]}
\]
は有限余次元を持つ．これは巡回gradingのdense性であり，分数環のstrong gradingより強い
局所的情報である．
\end{corollary}
\begin{proof}
$B\subset T$ は高次数で等号なので，$R\subseteq T^{(N)}$ でも高次数で等号である．
両辺は全section環で，各正次数 $m$ の次元は $3Nm$ だから
$R=T^{(N)}$ が全次数で成立する．特に大域生成される分数可逆層の切断空間が
一致するので，$\mathcal M=L_N$ も成り立つ．
前命題のrestriction同値と $T$ の通常のVeronese同値の間で
$(-)^{(N)}$ は可換する．従って $V$ は同値である．
ここで $T$ は次数 $3$ の種数 $1$ 曲線THCRなので一次生成であり，
通常のVeronese同値が適用できる．

densityを積からも確認する．$a\equiv r\pmod N$ を十分大きく固定する．
全ての十分大きい $c\equiv0\pmod N$ について $b=c-a$ と置く．
$B_a=T_a$，$B_b=T_b$，$B_c=T_c$ であり，$b>a$ としてよい．
$L_a$ の二つの大域切断を共通零点なしに選ぶと，basepoint-free pencilの完全列
\[
 0\longrightarrow L_a^{-1}\longrightarrow\mathcal O_C^{\oplus2}
 \longrightarrow L_a\longrightarrow0
\]
を得る．これを $\sigma^{a*}L_b$ でtensorする．
$\deg(\sigma^{a*}L_b L_a^{-1})=3(b-a)>0$ なので
$H^1(C,\sigma^{a*}L_b L_a^{-1})=0$ である．従って
\[
 H^0(C,L_a)\otimes H^0(C,\sigma^{a*}L_b)
 \longrightarrow H^0(C,L_c)
\]
は全射であり，$B_aB_b=B_c$ を得る．すなわち
$B_{[r]}B_{[-r]}$ は $B_{[0]}$ の十分高い次数の全体を含む．
\end{proof}

\begin{remark}[正確な到達点と特異曲線で残る障害]
\label{rem:v10-heart-status}
この結果はcohomology保存だけからheart同値を推論したものではない．
新たに使った情報は，各剰余層が可逆であることと，その次数が厳密に $3r$ であることから
積の分数層包含を等号へ強められる点である．
得られるのは $B\subseteq\widehat B=T$ の有限余次元性と元のtripleの幾何学的表示であり，
$B=\widehat B$ をこの証明だけでは得ない．従って $\epsilon=0$，一次生成性，
ambient $A$ のAS-regular性は別途の認識条件を要する．

特異な $C$ では，rank $1$ torsion-freeと次数の一致だけでは可逆性は従わない．
例えば算術種数 $1$ の有理節点曲線の正規化 $\nu:\mathbb P^1\to C$ に対して
$\nu_*\mathcal O_{\mathbb P^1}(3r-1)$ は次数 $3r$ のrank $1$ torsion-free層であるが，
節点で可逆ではない．また，$\sigma^N$ が特異モデル $C$ へ延長することから
$\sigma$ の延長を直ちには結論できない．従って一般の場合に残る課題は
この二つの局所条件の自動性，またはそれらを不要にする復元法である．

本節の命題の新規性は未確認である．証明は上記の通り直接与えたが，
Artin--Staffordの分類におけるorderやVeroneseの復元定理との完全な比較は未完了である．
\end{remark}



\subsection{三次境界の周期成分とEuler判別式}
\label{subsec:v10-cubic-discriminant}
以下では六公理，heart上exactな三乗根，非零normal元$g\in A_3$，
domain商$B=A/(g)$を仮定する．$\epsilon=\dim E_{-1}$と置くと，
定理\ref{thm:v8-cubic-standard-sections}から
\[
 q_n=\chi(L_0,L_n)=\frac{(n+1)(n+2)}2
          -\epsilon\mathbf1_{3\nmid n}\qquad(n\in\Z)
\]
である．このEuler列の周期部分は，rankだけでなく，Euler格子の判別式も制約する．

\begin{lemma}[周期成分の$A_2$型の形式と判別式]
\label{lem:v10-period-three-discriminant}
$x_i=[L_i]$，$0\leq i\leq4$とし，そのEuler行列を$G_\epsilon$とする．このとき
\[
 G_\epsilon=
 \begin{pmatrix}
 1&3-\epsilon&6-\epsilon&10&15-\epsilon\\
 -\epsilon&1&3-\epsilon&6-\epsilon&10\\
 -\epsilon&-\epsilon&1&3-\epsilon&6-\epsilon\\
 1&-\epsilon&-\epsilon&1&3-\epsilon\\
 3-\epsilon&1&-\epsilon&-\epsilon&1
 \end{pmatrix},\qquad
 \det G_\epsilon=3^5\epsilon^2.
\]
特に$\epsilon>0$なら$x_0,\ldots,x_4$は数値的に一次独立であり，
そのEuler形式の判別式の平方類は
$[3]\in\Q^\times/(\Q^\times)^2$である．

より具体的に
\[
 p_i=x_i+x_{i+1}+x_{i+2}\quad(0\leq i\leq2),\qquad
 u_0=x_3-3x_2+3x_1-x_0,\quad
 u_1=x_4-3x_3+3x_2-x_1
\]
と置くと，$\langle p_0,p_1,p_2\rangle_\Q$と$\langle u_0,u_1\rangle_\Q$は
Euler形式について左右両方に直交し，後者の行列は
\[
 9\epsilon\begin{pmatrix}2&-1\\-1&2\end{pmatrix}
\]
である．従って非零の周期成分は，Euler形式について正定値の二次元部分空間である．
\end{lemma}
\begin{proof}
exactな$s$はEuler形式を保つから$\chi(L_i,L_j)=q_{j-i}$であり，表示された行列を得る．
五つの生成元$(x_0,\ldots,x_4)$を$(p_0,p_1,p_2,u_0,u_1)$へ移す線形変換の行列は
\[
 T=\begin{pmatrix}
 1&0&0&-1&0\\
 1&1&0&3&-1\\
 1&1&1&-3&3\\
 0&1&1&1&-3\\
 0&0&1&0&1
 \end{pmatrix},\qquad \det T=27.
\]
直接積を取ると
\[
 T^{\mathsf t}G_\epsilon T=
 \begin{pmatrix}H_\epsilon&0\\0&9\epsilon A_2\end{pmatrix},
\]
\[
 H_\epsilon=
 \begin{pmatrix}
 15-6\epsilon&33-6\epsilon&60-6\epsilon\\
 6-6\epsilon&15-6\epsilon&33-6\epsilon\\
 6-6\epsilon&6-6\epsilon&15-6\epsilon
 \end{pmatrix},
 \quad A_2=\begin{pmatrix}2&-1\\-1&2\end{pmatrix}.
\]
ここで$\det H_\epsilon=3^6$，$\det(9\epsilon A_2)=3^5\epsilon^2$である．
従って$\det(T)^2\det G_\epsilon=3^{11}\epsilon^2$から
所望の行列式を得る．$\epsilon>0$なら行列式は非零なので，
数値的な一次関係は存在しない．基底変換は行列式を有理数の平方倍するため，
判別式の平方類は$[3^5\epsilon^2]=[3]$となる．
ここでEuler形式が対称である必要はない．同じ基底を始域と終域に使う双線形形式の行列は，
対称性の有無によらず$G\mapsto T^{\mathsf t}GT$で変換されるからである．
\end{proof}

\begin{proposition}[数値格子の判別式による欠損消滅]
\label{prop:v10-unimodular-cubic-defect}
Euler形式が，有限階数の自由abel群$N$上の整数値双線形形式を経由すると仮定する．
ここで全ての$[L_i]$を$N$へ写すことができ，$\operatorname{rank}N\leq5$とする．
さらに$\operatorname{rank}N=5$の場合は，$N$上のEuler行列が退化しているか，
またはその非零判別式の平方類が$[3]$と異なると仮定する．
このとき
\[
 \epsilon=0,\qquad E_n=0\ (n\in\Z),\qquad B=\widehat B.
\]
特に，$\operatorname{rank}N\leq5$でEuler形式がunimodularなら，この結論が成立する．
ここで$N$には，有限階数であることを別に確認できた数値Grothendieck格子を用いてよい．
\end{proposition}
\begin{proof}
$\epsilon>0$と仮定する．前補題により五つの像$x_0,\ldots,x_4\in N$は
一次独立なので，$\operatorname{rank}N=5$でなければならない．
$N$の整数基底でEuler行列を$G_N$，五つの像を列に持つ行列を$M$とすると，
\[
 G_\epsilon=M^{\mathsf t}G_NM,
 \qquad
 3^5\epsilon^2=(\det M)^2\det G_N.
\]
左辺は非零なので$G_N$は非退化であり，その判別式の平方類は$[3]$となる．
これは仮定に反する．特に$\det G_N=\pm1$なら，正の左辺は平方数でなければならないが，
$3^5\epsilon^2$の$3$-進付値は奇数なので不可能である．
従って$\epsilon=0$であり，命題\ref{prop:cubic-defect-mass}から全$E=0$と$B=\widehat B$を得る．
\end{proof}

\begin{corollary}[滑らかな高Veroneseまたは幾何学的境界による認識]
\label{cor:v10-five-exceptionals-recognition}
上の三次境界の設定に加え，次の幾何学的条件のいずれかを仮定する：
高Veronese表示の曲線が滑らかである，または構造対象とshiftを保つ同値
\[
 (\qgr B,D,(1))\simeq
 (\coh C,\OO_C,\mathcal L\otimes\sigma^*(-)),\qquad
 p_a(C)=1,\quad\deg\mathcal L=3
\]
があり，$C$がintegral projective curveであると仮定する．
次のいずれかが成り立てば，固定した$A$は三次元二次Koszul AS-regularである．
\begin{enumerate}[label=(\roman*)]
\item 前命題のEuler格子の仮定が成り立つ．
\item $D^b(\Cat)$が長さ$5$以下のfull exceptional collectionを持つ．
\end{enumerate}
(ii)の例外列にはstrong性やrank-one性を仮定しない．
結論が成り立つ場合，full exceptional collectionの長さは実際には$3$である．
\end{corollary}
\begin{proof}
高Veroneseの曲線が滑らかな場合には，定理\ref{thm:v10-smooth-veronese-heart}により
通常曲線表示を得る．従って両方の幾何学的条件を同じ場合へ帰着できる．
(i)では前命題から$\epsilon=0$である．
通常曲線表示の下での系\ref{cor:v8-geometric-one-obstruction}を適用すれば，
$B$の一次生成性と$A$のAS-regular性が従う．

(ii)で例外列の長さを$m\leq5$とする．full exceptional collectionは
$K_0(\Cat)\simeq\Z^m$を与え，そのEuler行列は対角成分が1の上三角行列である．
従ってEuler形式はunimodularであり，前命題と(i)を適用できる．
標準的な非可換射影平面の$K_0$は$\Z^3$なので，認識後には$m=3$となる．
\end{proof}

\begin{remark}[適用範囲と残る問題]
この判別式の議論は，五次元の軌道部分空間を\emph{全}数値空間と比較するものである．
全数値格子のrankが$6$以上である場合には，五次元部分格子の非平方判別式だけで
全格子のunimodularityに矛盾することはない．
六公理，両側$\chi$，有限torsion次元から，全数値格子のrankが$5$以下であることや
Euler形式のunimodularityを導いたわけではない．
また通常曲線表示を仮定しない場合の前命題の結論は，$E=0$と境界の飽和性までであり，
一次生成性や生成障害$\PP$の消滅を追加の証明なしに結論しない．
この節の格子計算は明示的なEuler列の初等的帰結である．
独立した文献上の新規性は確認していない．
\end{remark}

\subsection{Euler列の数値的剛性による別証明}

\begin{lemma}[整係数の有限rank軌道と準多項式]\label{lem:v7-integral-recurrence}
$q=(q_n)_{n\in\Z}$を整数列とし，その両側shiftの張る$\Q$ベクトル空間が有限次元であるとする．
さらに$|q_n|=O((1+|n|)^d)$と仮定する．
このとき，ある正整数$N$について，各$r\in\Z/N\Z$上の$q_n$は次数高々$d$の有理係数多項式である．
\end{lemma}
\begin{proof}
零列の場合は$N=1$で明らかである．以下では両側shiftを$U$，その軌道の$\Q$線形包を$V$とし，$r=\dim V>0$と置く．
有限個の整数$n_1,\ldots,n_r$を選んで評価写像
\[
 \operatorname{ev}:V\longrightarrow\Q^r,
 \qquad f\longmapsto(f(n_1),\ldots,f(n_r))
\]
を同型にできる．実際，有限次元の関数空間上では全ての点評価が双対空間を張る．
$\Lambda=\sum_{k\in\Z}\Z U^kq$と置けば，列の整係数性から
$\operatorname{ev}(\Lambda)\subset\Z^r$である．従って$\Lambda$は有限生成自由加群で，
$\Lambda\otimes_\Z\Q=V$，$U\Lambda=\Lambda$となる．
特に$U$はこの格子上の整数係数可逆行列である．

$V$の基底を有限個のshift $U^{k_j}q$から取る．評価写像を用いると，
$U^m$の各行列成分は有限個の$q_{m+k_j+n_i}$の固定した有理線形結合となる．
従って$m\to\pm\infty$の両方向で$\|U^m\|=O((1+|m|)^d)$である．
全ての固有値の絶対値は$1$であり，Jordan blockのサイズは高々$d+1$となる．

各$m>0$について$U^m$の特性多項式は整係数であり，その根の絶対値は全て$1$なので，
各係数は$r$だけで決まる定数で有界である．従ってこれらの特性多項式は有限種類しかない．
固定した固有値$\lambda$の全ての正の冪$\lambda^m$は，この有限個の多項式の根の有限集合に属する．
よって$\lambda$は1の冪根である．全ての固有値の位数の公倍数を$N$とすれば，
$(U^N-1)^{d+1}=0$である．従って全整数$\ell$について
\[
 (U^N)^\ell=\sum_{j=0}^d\binom{\ell}{j}(U^N-1)^j.
\]
この式に点評価を施せば，各剰余類上での所要の多項式表示を得る．
\end{proof}

\begin{proposition}[有限Euler rankから三次境界の標準Hilbert列へ]
\label{prop:v7-finite-rank-cubic-euler}
六公理とexactな三乗根を仮定し，$0\ne g\in A_3$がnormalで，$B=A/(g)$はdomainとする．
さらに軌道$\{[\OO(n)]:n\in\Z\}$上のEuler形式のrankが有限であると仮定する．
そのrankに数値的な上界は課さない．$\epsilon=\dim E_{-1}$と置くと
\[
 \widehat b_n=3n\quad(n>0),\qquad
 q_n=P(n)-\epsilon\mathbf1_{3\nmid n}\quad(n\in\Z)
\]
が成立する．従って軌道上のEuler rankは$3$または$5$に限られ，
それぞれ$\epsilon=0$または$\epsilon>0$に対応する．
逆に$\widehat b_n=3n$ for all $n>0$ならEuler orbit rankは有限である．
従ってこの三次境界の設定では，Euler orbit rankの有限性と飽和後の標準Hilbert列は同値である．
\end{proposition}
\begin{proof}
有限Euler rankから，整数列$q_n=\chi(\OO,\OO(n))$の両側shiftの線形包は有限次元である．
これは定理\ref{thm:numericalrank}で用いた軌道Euler行列の有限rankと同じ条件である．
また，$h_n=n^2/2+O(n)$，normal元による$E$の有限性，Serre対称性から
\[
 q_n=P(n)+O(1+|n|)\qquad(n\in\Z)
\]
である．前補題に$d=2$を適用する．必要なら周期を3倍して$3\mid N$とし，
$f_n=q_n-P(n)$と置くと，各$r\in\Z/N\Z$について
\[
 f_n=a_rn+b_r\qquad(n\equiv r\pmod N)
\]
と書ける．二次の係数が消えることには上の一様な$O(1+|n|)$評価を用いた．

切断完全列と導来随伴より
\[
 q_n-q_{n-3}
 =\chi_{\qgr B}(D,D(n))
 =\widehat b_n\qquad(n>0).
\]
最後の等式は系\ref{cor:cubic-boundary-duality}の正次数Ext消滅による．
系\ref{cor:cubic-saturated-hilbert}の上界と$P(n)-P(n-3)=3n$から
\begin{equation}\label{eq:v7-cubic-monotone}
 f_n-f_{n-3}\leq0\qquad(n>0)
\end{equation}
を得る．$n\equiv r\pmod N$を保って$n\to+\infty$とすると
$a_r\leq a_{r-3}$である．有限集合上の$r\mapsto r-3$の各巡回軌道に沿って比較すれば，
$a_r=a_{r-3}$となる．従って$a_r$は$r\bmod3$だけに依存する．

この等式を\eqref{eq:v7-cubic-monotone}に戻すと
\[
 3a_r+b_r-b_{r-3}\leq0.
\]
同じ巡回軌道上で和を取ることにより$a_r\leq0$となる．
一方，Serre対称性$f_n=f_{-n-3}$は
$a_{-r-3}=-a_r$を与える．両辺の添字について$a\leq0$であるから，全ての$a_r$は零である．
従って$b_r\leq b_{r-3}$であり，再び巡回軌道に沿って比較すると$b_r=b_{r-3}$となる．
すなわち$f$の周期は3を割る．

$q_{3m}=P(3m)$から剰余類0で$f=0$であり，
Serre対称性から剰余類1と2での値は一致する．
$q_{-1}=-\dim E_{-1}=-\epsilon$と$P(-1)=0$から，その値は$-\epsilon$である．
これで全整数における$q_n$の式を得る．その3刻みの差から
$q_n-q_{n-3}=3n$となり，正次数では$\widehat b_n=3n$である．
最後のrankの値と逆向きは系\ref{cor:cubic-defect-Euler-rank}による．
\end{proof}

\begin{remark}[有限rank条件の位置]
定理\ref{thm:v8-cubic-standard-sections}と系\ref{cor:v8-cubic-automatic-euler-rank}により，
三次domain境界ではこの有限rank条件は自動的に満たされる．
本命題は，曲線分類やVeroneseのcohomology比較を使わず，有限rankから同じ結論を導く独立の数値的証明である．
ただし定理\ref{thm:numericalrank}のrank $\leq4$を単に有限性へ置き換えて
$E_{-1}=0$と結論することはできず，rank $5$と$\epsilon>0$は残る．
\end{remark}

\begin{remark}[残る生成障害]
以上はsection欠損と認識の最初の障害$E_{-1}$を完全に対応させる．
一方，$B=\widehat B$でも$B$が一次生成されるとは限らず，
研究ノートの$B_{\mathrm{node}}$がその例である．
従ってsection欠損の除去と高次生成元の除去には，別のambientの論証が必要である．
可換境界では後の定理\ref{thm:v7-commutative-cubic-recognition}がその論証を与える．
$B_{\mathrm{node}}$は境界だけの数値と飽和性の不足を示す例であり，六公理を満たすnormal被覆の存在例ではない．
\end{remark}

\begin{corollary}[不均衡なsection欠損によるnormal持上げの障害]\label{cor:cubic-unbalanced-no-lift}
$T$を通常の滑らかな平面三次曲線の座標環とし，$V\subset T_1$を二次元部分空間とする．
\[
 B'=\C\oplus V\oplus\bigoplus_{n\geq2}T_n
\]
は，六公理とexactな三乗根を満たす軌道環$A$の，次数$3$の非零normal元による商にはなり得ない．
\end{corollary}
\begin{proof}
$B'$は$T$の有限余次元graded部分環であり，その飽和は$T$である．
従って$\widehat B'/B'$は次数$1$だけにある一次元空間であり，
命題\ref{prop:cubic-defect-mass}の二つの剰余類の質量の等式に反する．
\end{proof}

\subsection{一次生成された整な三次境界の全次数での識別}
\label{subsec:cubic-boundary-identification}

以下の曲線論は特異点を許す．特に，算術種数$1$から滑らかさを結論しない．

\begin{lemma}[算術種数1の整な曲線上の次数3線束]
\label{lem:integral-genus-one-sections}
$Y$を$\C$上のintegral projective curveとし，$p_a(Y)=1$とする．
このとき$Y$はGorensteinで$\omega_Y\simeq\OO_Y$である．
任意の正次数のinvertible sheaf $\mathcal M$について
\[
 H^1(Y,\mathcal M)=0,\qquad h^0(Y,\mathcal M)=\deg\mathcal M.
\]
次数$3$のinvertible sheaf $\mathcal L$はvery ampleであり，その完全線形系は
$Y$をintegral plane cubicとして埋め込む．
さらに$\deg\mathcal L=3$，$\deg\mathcal M\geq3$なら，積写像
\[
 H^0(Y,\mathcal L)\otimes H^0(Y,\mathcal M)
   \longrightarrow H^0(Y,\mathcal L\otimes\mathcal M)
\]
は全射である．
\end{lemma}

\begin{proof}
Integral curveはCohen--Macaulayである．そのdualizing sheaf $\omega_Y$は
rank one torsion-freeで，曲線の双対性により
\[
 h^0(Y,\omega_Y)=h^1(Y,\OO_Y)=1,
 \qquad h^1(Y,\omega_Y)=h^0(Y,\OO_Y)=1
\]
である\cite[Tags 0BS2, 0BS3]{Stacks}．
非零sectionによる射$\OO_Y\to\omega_Y$は単射で，余核$\mathcal Q$は有限長である．
両辺のEuler標数がともに零なので
$\operatorname{length}(\mathcal Q)=\chi(\omega_Y)-\chi(\OO_Y)=0$であり，
$\omega_Y\simeq\OO_Y$となる．特に$Y$はGorensteinである．

正次数の$\mathcal M$に対して$\mathcal M^{-1}$の次数は負なので非零sectionを持たない．
実際，非零sectionは$\OO_Y\hookrightarrow\mathcal M^{-1}$を与え，
その余核の長さが$\deg\mathcal M^{-1}<0$となって矛盾する．
双対性とRiemann--Rochから
\[
 H^1(Y,\mathcal M)^*\simeq H^0(Y,\mathcal M^{-1})=0,
 \qquad h^0(Y,\mathcal M)=\deg\mathcal M
\]
を得る．

Very amplenessも特異点を許して確認できる．長さ$1$または$2$の閉部分スキーム
$Z\subset Y$とそのideal $\mathcal I_Z$を取る．
$\mathcal F=\mathcal I_Z\mathcal L$はrank one torsion-freeで
\[
 \chi(\mathcal F)=3-\operatorname{length}(Z)>0.
\]
もし$\mathcal F\to\OO_Y$が非零なら単射であり，余核は有限長だから
$\chi(\mathcal F)\leq\chi(\OO_Y)=0$となって矛盾する．
従って双対性から$H^1(Y,\mathcal I_Z\mathcal L)=0$となり，
$H^0(Y,\mathcal L)\to H^0(Z,\mathcal L|_Z)$は全射である．
長さ$2$の部分スキームを分離する判定により$\mathcal L$はvery ampleである．
この判定を含む一般形は\cite[Tags 0E3A, 0H2V]{Stacks}を参照．
$h^0(Y,\mathcal L)=3$，$\deg\mathcal L=3$なので，埋め込み先は$\mathbb P^2$，
像はintegral cubicである．

次に積写像を調べる．$\mathcal L$はglobally generatedなので，
共通零点を持たない二次元部分空間$V\subset H^0(Y,\mathcal L)$を選べる．
具体的には非零sectionの有限個の零点で消えない第二のsectionを選べばよい．
従ってbasepoint-free pencilの完全列
\[
 0\longrightarrow\mathcal L^{-1}
 \longrightarrow V\otimes\OO_Y
 \longrightarrow\mathcal L\longrightarrow0
\]
がある．これに$\mathcal M$をtensorする．
$\deg\mathcal M>3$なら$\mathcal M\mathcal L^{-1}$は正次数なので
$H^1(Y,\mathcal M\mathcal L^{-1})=0$であり，
既に$V\otimes H^0(Y,\mathcal M)$から全射となる．

$\deg\mathcal M=3$かつ$\mathcal M\not\simeq\mathcal L$の場合には，
$\mathcal M\mathcal L^{-1}$は非自明な次数零の線束である．
次数零の線束に非零sectionがあれば，そのsectionの余核の長さが零なので線束は自明となる．
従って$H^1(Y,\mathcal M\mathcal L^{-1})=0$であり，同じ議論が使える．

残る$\mathcal M\simeq\mathcal L$の場合には，$\mathcal L$によるplane cubicの埋め込みを使う．
Integral plane cubicを含む非零二次曲線は存在しないので
\[
 \operatorname{Sym}^2 H^0(Y,\mathcal L)
   \longrightarrow H^0(Y,\mathcal L^{\otimes2})
\]
は単射である．両辺の次元は$6$なので同型であり，積写像の全射性が従う．
\end{proof}

\begin{proposition}[一次生成された三次normal境界の識別]
\label{prop:cubic-boundary-full-identification}
六公理と固定したexactな三乗根の下で，$0\ne g\in A_3$がnormalであり，
$B=A/(g)$がdomainかつ$B_1$で生成されると仮定する．
このとき，integral plane cubic $Y$，次数$3$の線束$\mathcal L$，
$\sigma\in\operatorname{Aut}(Y)$が存在して
\[
 B\simeq B(Y,\mathcal L,\sigma)
\]
が全次数のgraded代数として成り立つ．特に
\[
 \dim B_0=1,\qquad \dim B_n=3n\ (n>0),\qquad
 H_B(t)=\frac{1+t+t^2}{(1-t)^2},\qquad
 H_A(t)=\frac1{(1-t)^3}.
\]
さらに商$B$は全整数次数で飽和しており，
$E_n=\Ext^1(\OO,\OO(n))=0$ for all $n\in\Z$である．
\end{proposition}

\begin{proof}
命題\ref{prop:cubic-normal-hilbert-bound}から$b_1=3$を得る．
またArtin--Staffordの曲線分類により全次数のgraded単射
\[
 B\lhook\joinrel\longrightarrow T=B(Y,\mathcal L,\sigma)
\]
があり，余核は有限次元で，$\deg\mathcal L=3$，$p_a(Y)=1$である．
補題\ref{lem:integral-genus-one-sections}から
$h^0(Y,\mathcal L_n)=3n$ for $n>0$であり，
特に$B_1\hookrightarrow T_1$は次元$3$の同型である．

$\mathcal L_n=\mathcal L\otimes\sigma^*\mathcal L\otimes\cdots
\otimes(\sigma^{n-1})^*\mathcal L$と書く．
$T_nT_1\to T_{n+1}$は，$H^0(Y,\mathcal L)$を$\sigma^{n*}$で移した後の通常の積写像
\[
 H^0(Y,\mathcal L_n)\otimes H^0(Y,\sigma^{n*}\mathcal L)
  \longrightarrow H^0(Y,\mathcal L_{n+1})
\]
である．二つの線束の次数は$3n$と$3$なので，上の補題により全射である．
従って$T$は$T_1$で生成される．$B_1=T_1$なので，単射$B\hookrightarrow T$の像は$T$全体である．
これで低次数を含む$B=T$が示された．

上記の次元式から$H_B(t)=1+3t/(1-t)^2$を得る．
$g$はregularなので$H_A(t)=H_B(t)/(1-t^3)=(1-t)^{-3}$である．
曲線上のample線束は$\sigma$-ampleであり，THCRの標準同値
$\qgr B\simeq\coh Y$は$\pi B$を$\OO_Y$，その次数シフトを対応する線束へ移す
\cite[Theorem 3.2, Corollary 5.3]{SVdBsurvey}．
従って非負次数では$B$の定義そのものが標準section環であり，
負次数では対応する線束の次数が負なのでsectionがない．
すなわち，全$n\in\Z$で
\[
 B_n\xrightarrow{\sim}
 \Hom_{\qgr B}(\pi B,\pi B(n))
\]
である．命題\ref{prop:saturated-boundary}を適用して$E=0$を得る．
\end{proof}

\begin{remark}[適用範囲]
$Y$の滑らかさや$\sigma$の位数は仮定していない．
$\omega_Y\simeq\OO_Y$は，特異な$Y$について
$D^b(\coh Y)$全体がSerre関手$[1]$を持つという意味ではない．
Ambient代数$A$の一次生成性・二次性・AS-regular性は，次項のnormal被覆の二分と有限ホモロジー次元を使って導く．

識別の最後の部分は，一次生成性より弱いデータでも使える．
すなわち，ある次数$3$・算術種数$1$のTHCR $T$への全次数のgraded単射
$B\hookrightarrow T$が既に与えられ，$\dim B_1=3$なら，$T$の一次生成性から直ちに$B=T$となる．
ただし，Artin--Staffordの上記の全次数の単射を得る際には$B$の一次生成性を用いており，
元の六公理からその仮定を取り除いたとは主張しない．
\end{remark}

\subsection{整な三次境界の二次包絡とnormal被覆}

本項の曲線上の関係式とoverlapの幾何は，Artin--Tate--Van den Bergh\cite[Theorem 6.6, \S7]{ATV1990}に既知である．
ここではintegralな場合の適用条件を明らかにするため，評価核のstabilityを用いて必要な部分を証明する．

以下では$C$をintegral projective Gorenstein curve，
$p_a(C)=1$，$\omega_C\simeq\OO_C$とする．
滑らかさは仮定しない．$\chi(\OO_C)=0$なので，torsion-free sheafの次数は任意のrankで$\deg F=\chi(F)$となる．

\begin{lemma}[次数3の評価核]\label{lem:cubic-evaluation-kernel}
$L$を$C$上の次数3の可逆層とし，$V=H^0(C,L)$と置く．
評価の完全列
\[
 0\longrightarrow M_L\longrightarrow V\otimes\OO_C
 \longrightarrow L\longrightarrow0
\]
において，$M_L$はrank $2$，degree $-3$のstable vector bundleである．
また，任意の次数3の可逆層$L'$について$M_L\otimes L'$は全体生成され，
その大域切断空間は三次元である．
\end{lemma}
\begin{proof}
まずrank one torsion-free sheaf $F$について
$F^*=\mathcal Hom(F,\OO_C)$と置く．
$C$はGorenstein curveなので$F$は各局所環上でmaximal Cohen--Macaulayである．
具体的に，$R=\OO_{C,p}$，$d=\dim R\in\{0,1\}$とすると，
$R[d]$はnormalized dualizing complexであり，$F_p$は次元$d$のCohen--Macaulay加群である．
Cohen--Macaulay duality \cite[Tag 0B5A, Lemma 47.16.7]{Stacks} により
$R\!\operatorname{Hom}_R(F_p,R[d])$は次数$-d$に集中する．
従って$\mathcal Ext^{>0}(F,\OO_C)=0$である．局所大域スペクトル系列とSerre双対性から
\[
 H^0(F^*)\simeq H^1(F)^*,\qquad
 H^1(F^*)\simeq H^0(F)^*,\qquad
 \deg F^*=-\deg F
\]
を得る．また，非零切断$\OO_C\to F$は単射で，余核の長さは$\deg F$である．
従って負次数では$H^0(F)=0$，次数0で切断を持てば$F\simeq\OO_C$，
次数1なら$h^0(F)=1$である．

$M_L$に次数$\geq-1$のrank one torsion-free subsheaf $F$があると仮定する．
$\deg F>0$では$H^0(F^*)=0$なので$\Hom(F,V\otimes\OO_C)=0$となる．
$\deg F=0$では，$H^0(F^*)=0$なら同じ矛盾であり，
$F^*\simeq\OO_C$なら$\Hom(F,M_L)$は
$V\otimes H^0(F^*)\to H^0(F^*\otimes L)$の核なので零である．
$\deg F=-1$では$h^0(F^*)=1$であり，その非零切断による
$V\to H^0(F^*\otimes L)$は$C$のintegralityと$F^*$のtorsion-freenessにより単射である．
この場合も$\Hom(F,M_L)=0$となる．従って任意のrank one subsheafの次数は$-2$以下であり，
$M_L$はstableである．

$E=M_L\otimes L'$はstable，rank $2$，degree $3$である．
正次数のrank $2$のstable vector bundle $G$は$\Hom(G,\OO_C)=0$なので$H^1(G)=0$である．
実際，非零射の像$F\subset\OO_C$は次数$\leq0$のrank one torsion-free sheafであり，
その核はstabilityに反する次数のrank one subsheafとなる．
したがって$h^0(E)=\chi(E)=3$である．

任意の閉点$p$のideal sheafを$I_p$とする．$\deg I_p^*=1$である．
$E$から$I_p^*$への非零射があれば，その像の次数は$1$以下であるが，
$E$のrank one torsion-free quotientはstabilityにより次数$2$以上でなければならない．
従って$\Hom(E,I_p^*)=0$である．$E$は局所自由なので
\[
 H^1(E\otimes I_p)^*
 =\Hom(E\otimes I_p,\OO_C)
 =\Hom(E,I_p^*)=0.
\]
ゆえに$H^0(E)\to E\otimes k(p)$は全射である．これは特異点にも適用され，
$E$の全体生成が従う．
\end{proof}

\begin{lemma}[三次までの関係式と唯一のoverlap]\label{lem:cubic-relations-overlap}
$L$を次数3の可逆層，$\sigma\in\operatorname{Aut}C$とし，
$B=B(C,L,\sigma)$と置く．$V=B_1$，
$R=\ker(V\otimes V\to B_2)$，$S=T(V)/(R)$とする．このとき次が成立する．
\begin{enumerate}
\item $B$は次数1で生成され，その定義idealは次数2と3で生成される．
\item $L_i=\sigma^{i*}L$と置くと，標準の同一視の下で
\[
 W=(R\otimes V)\cap(V\otimes R)
 \simeq H^0\bigl(C,L_0\otimes L_1^{-2}\otimes L_2\bigr).
\]
従って$\dim W$は$0$または$1$である．
\item $W=0$なら$B=S$である．$\dim W=1$なら，任意の$0\ne w\in W$による
二つの外側のcontraction $V^*\to R$はともに同型である．
この場合$S$は三次元二次AS-regular代数で，
ある$h\in S_3$について$B=S/(h)$である．
\end{enumerate}
\end{lemma}
\begin{proof}
次数1生成は次数3以上の可逆層の混合積の全射性による．
関係式の次数を制御するため，$V_i=H^0(L_i)$とし，$n\geq2$について
\[
 K_n=\ker\left(V_0\otimes H^0(L_1\cdots L_{n-1})
 \longrightarrow H^0(L_0\cdots L_{n-1})\right)
 =H^0(M_{L_0}\otimes L_1\cdots L_{n-1})
\]
と置く．$n\geq3$なら
$K_n\otimes V_n\to K_{n+1}$は全射である．実際，
$L_n$のbase-point-free pencil $U\subset V_n$を選び，
$0\to L_n^{-1}\to U\otimes\OO_C\to L_n\to0$を
$E_n=M_{L_0}\otimes L_1\cdots L_{n-1}$とtensorする．
$E_n\otimes L_n^{-1}$はstable，rank $2$，degree $6n-15>0$なので$H^1=0$である．
従って$K_n\otimes U\to K_{n+1}$が全射となる．
具体的に，$I_{[a,b)}$を$V_a\otimes\cdots\otimes V_{b-1}$から
$H^0(L_a\cdots L_{b-1})$への積の核とする．尾部の積による商写像は
\[
 I_{[0,n)}\twoheadrightarrow K_n,
 \qquad
 \ker=V_0\otimes I_{[1,n)}
\]
を与える．$I_{[0,n+1)}$の元の$K_{n+1}$における像を，
$K_n\otimes V_n$の元で持ち上げ，さらに$I_{[0,n)}\otimes V_n$に持ち上げる．
差は尾部の関係に属するので，$n\geq3$について
\[
 I_{[0,n+1)}=
 I_{[0,n)}\otimes V_n+V_0\otimes I_{[1,n+1)}
\]
となる．$\sigma$による移送で任意の始点に同じ議論が適用できるから，
次数に関する帰納法により関係idealは次数$2,3$で生成される．

overlapを計算する．$R_{01}=\ker(V_0\otimes V_1\to H^0(L_0L_1))$とすると，
前補題による全体生成から
\[
 0\longrightarrow Q\longrightarrow R_{01}\otimes\OO_C
 \longrightarrow M_{L_0}\otimes L_1\longrightarrow0,
 \qquad Q=L_0\otimes L_1^{-2}
\]
を得る．ここでkernelは可逆層で，その表示は行列式から従う．
$L_2$とtensorして大域切断を取れば
\[
 \ker\left(R_{01}\otimes V_2\longrightarrow
 V_0\otimes H^0(L_1L_2)\right)
 =H^0(Q\otimes L_2).
\]
左辺が$(R_{01}\otimes V_2)\cap(V_0\otimes R_{12})$である．
$\sigma$による$V_i\simeq V$の同一視で主張の$W$を得る．
$Q\otimes L_2$は次数0の可逆層なので，その切断次元は$0$または$1$である．
また，$\dim B_3=9$，$\dim R=3$なので
\[
 \dim S_3=27-18+\dim W=9+\dim W.
\]
従って次数3の独立な関係の個数は$\dim W$である．関係式が三次までで生成されることから，
$W=0$なら$B=S$，$\dim W=1$なら$B=S/(h)$，$\deg h=3$を得る．

$\dim W=1$なら$Q\otimes L_2\simeq\OO_C$である．上の評価列を双対化すると，
$H^0((M_{L_0}\otimes L_1)^*)=0$から
\[
 R_{01}^*\longrightarrow H^0(Q^{-1})=H^0(L_2)
\]
が単射であり，両辺三次元なので同型である．これは$w$の右外側のcontractionが同型であることをいう．
$L_0,L_2$の役割を交換し，$R_{12}\otimes\OO_C\to L_1\otimes M_{L_2}$を用いると，
左外側も同様である．

最後に$S$のAS-regular性を示す．$V$の基底$x_1,x_2,x_3$を選び，
二つのcontractionが同型であることを使って
\[
 w=\sum_i x_i\otimes f_i=\sum_j f'_j\otimes x_j,
 \qquad f=Mx,\quad (f')^t=x^tM
\]
と書く．ここで$f_i,f'_j$はいずれも$R$の基底であり，$M$は一次式の$3\times3$行列である．
従って$S$はArtin--Tate--Van den Berghの意味でstandardである．
$f'=Q_0f$となる$Q_0\in GL_3(\C)$を取る．
$w=\sum_{ijk}w_{ijk}x_i\otimes x_j\otimes x_k$と書けば
$w_{ijk}=\sum_a(Q_0)_{ka}w_{aij}$である．
非零の$p=(p_j)\in V^*$に対し，
$r_p=\sum_jp_jf'_j\in R$の係数行列を$N_p$とすると
\[
 M(p)=N_p^tQ_0^t.
\]
$r_p\ne0$であり，$B$はdomainなので$R$は非零のrank one tensorを含まない．
従って$\operatorname{rank}N_p\geq2$，ゆえに$\operatorname{rank}M(p)\geq2$である．
これで$M$の$2\times2$ minorの共通零点は空となる．
Artin--Tate--Van den Berghのstandard algebraの判定
\cite[Theorem 1]{ATV1990}（\cite[Theorem 2.2]{IM}にも再掲）により$S$は三次元二次AS-regularである．
\end{proof}


\begin{theorem}[整な三次境界のnormal被覆の二分]
\label{thm:cubic-normal-cover-dichotomy}
$A$をlocally finite connected graded右noetherian $\C$-代数，
$g\in A_3$を両側regularなnormal元とする．次数付き代数として
\[
 B=A/(g)\simeq B(C,L,\sigma),\qquad p_a(C)=1,\quad\deg L=3
\]
と書けると仮定する．ここで$C$はintegral projective curve，$L$は可逆層，
$\sigma\in\operatorname{Aut}C$である．このとき次のいずれかが成立する．
\begin{enumerate}
\item $A$は三次元二次Koszul AS-regular代数で，Gorenstein parameterは$3$である．
\item あるgraded automorphism $\nu$について
$A\simeq B[z;\nu]$，$\deg z=3$であり，この同型は$z$を$g$に移す．
\end{enumerate}
特に$\qgr A$が有限ホモロジー次元を持てば第一の場合が成立する．
$A$のdomain性，ambientなSerre関手，六公理はこの定理の仮定ではない．
また，$g$の中心性，$C$の滑らかさ，$\sigma$の位数も仮定しない．
\end{theorem}
\begin{proof}
補題\ref{lem:integral-genus-one-sections}により$C$はGorensteinで$\omega_C\simeq\OO_C$であり，
$H_B(t)=(1-t^3)/(1-t)^3$である．$g$のregularityから
\[
 H_A(t)=\frac{H_B(t)}{1-t^3}=\frac1{(1-t)^3}.
\]
商写像$A\to B$は次数$1,2$で同型なので，
補題\ref{lem:cubic-relations-overlap}の二次包絡$S=T(B_1)/(R)$から$A$へのgraded homomorphismがある．
また$B$は一次生成なので，次数に関する帰納法から$A$は$A_1$と$g$で生成される．
実際，斉次元$a\in A_n$の商での像を次数1の元の積で表し，そのliftとの差を
$gA_{n-3}$に入れれば，次数が下がる．

同補題の$W$が零なら$B=S$であり，$B\to A$は商写像のsectionとなる．
$W\ne0$の場合には$S$は二次AS-regularで，$B=S/(h)$，$\deg h=3$である．
$h$の$A$における像は$\lambda g$，$\lambda\in\C$となる．
$\lambda\ne0$なら$S\to A$の像は$A_1$と$g$を含むので全射であり，
双方のHilbert級数が$(1-t)^{-3}$だから同型となる．
二次三次元AS-regular代数はKoszulでparameter $3$を持つので，第一の場合である．
$\lambda=0$なら再び$B\to A$がsectionとなる．

sectionがある場合を調べる．$ga=\nu(a)g$により定まるgraded automorphism $\nu$は
$A_1$を保つので，sectionの像$\langle A_1\rangle\simeq B$を保つ．
従って
\[
 B[z;\nu|_B]\twoheadrightarrow A,\qquad z\longmapsto g,\quad\deg z=3
\]
という全射があり，Hilbert級数の一致から同型となる．これが第二の場合である．

最後に$\qgr A$の有限ホモロジー次元を仮定する．第二の場合には系\ref{cor:split-pd}から
$\operatorname{pd}_{\Gr B}\C<\infty$となり，$1/H_B(t)$はLaurent多項式でなければならない．
しかし
\[
 \frac1{H_B(t)}=\frac{(1-t)^2}{1+t+t^2}
\]
は原始三乗根に極を持つ．よって第二の場合は排除される．
\end{proof}

\begin{theorem}[一次生成される整な三次normal境界からの認識]
\label{thm:cubic-generated-boundary-recognition}
六公理と固定したexactな三乗根の下で，$0\ne g\in A_3$がnormalで，
$B=A/(g)$がdomainであるとする．このとき次は同値である．
\begin{enumerate}
\item $B$は次数$1$で生成される．
\item $A$は三次元二次Koszul AS-regular代数で，Gorenstein parameterは$3$である．
\end{enumerate}
この条件が成立すれば全$n\in\Z$で$E_n=0$であり，
$(\OO,\OO(1),\OO(2))$はfull strong exceptional collectionとなる．
\end{theorem}
\begin{proof}
一次生成性を仮定すると，命題\ref{prop:cubic-boundary-full-identification}が
$B$をintegral plane cubicの次数$3$のTHCRと全次数で識別し，$E=0$を与える．
$A$はright noetherian domainであり，$g$はregularである．また
補題\ref{lem:dimension}により$\operatorname{hdim}\qgr A=2$なので，
定理\ref{thm:cubic-normal-cover-dichotomy}を適用できる．
逆に二次AS-regularな$A$は一次生成され，その商$B$も一次生成される．
full strong性は定理\ref{thm:finite}の同値条件による．
\end{proof}

\begin{corollary}[三次境界のregular triple条件]
\label{cor:cubic-regular-triple}
定理\ref{thm:cubic-normal-cover-dichotomy}の有限ホモロジー次元の場合，
特に定理\ref{thm:cubic-generated-boundary-recognition}の同値条件が成立する場合，
境界の偏極は
\[
 L\otimes\sigma^{2*}L\simeq(\sigma^*L)^{\otimes2}
\]
を満たす．
\end{corollary}
\begin{proof}
二分の証明で$W=0$の場合は分裂するため，有限ホモロジー次元の下では$W\ne0$である．
従って$L\otimes(\sigma^*L)^{-2}\otimes\sigma^{2*}L$は非零切断を持つ次数零の線束であり，自明である．
\end{proof}

\subsection{重み付き完全交叉境界の全normal三次被覆の排除}
\label{subsec:v6-weighted-cubic-cover}

以下の結果は，特定の可換な錐だけでなく，同じ境界を持つ全てのnormal三次被覆を扱う．
六公理やambientのSerre関手は仮定しない．証明は下に記す直接的な計算による．
この被覆分類と認識への応用の文献上の新規性は未確認である．

\begin{lemma}[中心的な局所化を持つOre拡大の頂点]
\label{lem:v6-ore-differential-vertex}
$R$をconnected locally finite graded右noetherian代数，
$\sigma$をgraded automorphism，$\delta$を次数$e>0$の
$\sigma$-derivationとし，
\[
 A=R[z;\sigma,\delta],\qquad zr=\sigma(r)z+\delta(r),
 \qquad \deg z=e
\]
が右noetherianであるとする．$\delta(R_{\geq1})\subset R_{\geq1}$を仮定し，
ある$m>0$で$y=z^m$が$A$の中心元であるとする．
$M=A/R_{\geq1}A\simeq\C[z]$，$N=\pi M$と置く．
$k_R$の最小次数付き自由分解の第$i$項を
$P_i=\bigoplus_j R(-a_{ij})$と書けば
\[
 \Ext^i_{\qgr A}(N,N(r))
 \simeq\bigoplus_j\C[z,z^{-1}]_{a_{ij}+r}
 \qquad(i\geq0,\ r\in\Z)
\]
である．特に$\qgr A$が有限homological dimensionを持てば
$\operatorname{pd}_{\operatorname{Gr}R}k_R<\infty$となる．
\end{lemma}
\begin{proof}
$R_{\geq1}A$は両側idealであり，その商は$\C[z]$である．
$D=A[y^{-1}]$，$\widetilde M=M[y^{-1}]=\C[z,z^{-1}]$と置く．
中心的局所化$L:\operatorname{Gr}A\to\operatorname{Gr}D$はexactであり，
exactな制限functorを右随伴に持つ．任意のtorsion $A$加群$T$と
$D$加群$V$について，導来随伴から
\[
 \Ext^i_{\operatorname{Gr}A}(T,\operatorname{res}V)
 =\Ext^i_{\operatorname{Gr}D}(T[y^{-1}],V)=0
 \quad(i\geq0)
\]
である．従って局所化から誘導されるexact functor
$\overline L:\operatorname{QGr}A\to\operatorname{Gr}D$の右随伴
$J=\pi\operatorname{res}$はexactかつ完全忠実であり，injectiveを保つ．
ゆえに
\[
 \Ext^i_{\operatorname{QGr}A}(JV,JW)
 \simeq\Ext^i_{\operatorname{Gr}D}(V,W).
\]
$\widetilde M/M$はtorsionである．実際，$R_{\geq1}$は零作用し，
各元はある正の$z$冪で消える．従って$N=J\widetilde M$となる．
局所化の随伴性を再び使うと
\[
 \Ext^i_{\operatorname{QGr}A}(N,N(r))
 \simeq\Ext^i_{\operatorname{Gr}A}(M,\widetilde M(r)).
\]
$A$は左$R$加群として自由なので，$P_\bullet\otimes_R A$は$M$の自由分解である．
その微分の行列成分は$R_{\geq1}$に入り，$\widetilde M$には零作用する．
従って$\Hom_A(-,\widetilde M(r))$を施した複体の微分は全て零であり，表示した式を得る．
右noetherian性により，noetherian対象間のYoneda Extは
$\operatorname{QGr}A$と$\qgr A$で一致する．
$P_i\ne0$なら$r\equiv-a_{ij}\pmod e$を選べば右辺は非零であり，最後の主張が従う．
\end{proof}

\begin{theorem}[重み$(1,1,1,2)$の三次完全交叉境界のnormal被覆]
\label{thm:v6-weighted-cubic-cover}
\[
 R=\C[u,v,w]/(uw-v^2),\qquad \deg u=\deg v=\deg w=1,
\]
とし，$0\ne\ell\in R_1$，$p\in R_3$について
\[
 B=R[z]/(\ell z-p),\qquad \deg z=2
\]
がdomainであるとする．$A$をconnected locally finite graded右noetherian代数，
$g\in A_3$を両側regularなnormal元とし，$A/(g)\simeq B$を仮定する．
このとき$\qgr A$のhomological dimensionは無限である．

より正確には，境界の生成元を次数を保って取り直した上で，$A$は次のいずれかである．
\begin{enumerate}
\item 分裂Ore拡大$B[t;\nu]$，$\deg t=3$．この表示では$t\mapsto g$である．
\item 可換な重み付き錐$R[z]$．この表示では$g$は$\ell z-p$の非零定数倍である．
\item $\ell=v$かつ$p=a u^3+b w^3$と書ける場合のOre拡大
\[
 R[z;\sigma,\delta],\qquad
 \sigma(u)=u,\quad\sigma(w)=w,\quad\sigma(v)=-v,
 \qquad\delta(u)=\delta(w)=0,\quad\delta(v)=2p.
\]
すなわち$u,w$は中心的で，$v^2=uw$，$zv+vz=2p$である．
この表示でも$g$は$vz-p$の非零定数倍であり，$z^2$は中心元である．
\end{enumerate}
\end{theorem}
\begin{proof}
$R[z]$はdomainであり，$\ell z-p$は非零なので
\[
 H_B(t)=\frac{1-t^3}{(1-t)^3},\qquad
 H_A(t)=\frac{H_B(t)}{1-t^3}=\frac1{(1-t)^3}.
\]
$u,v,w,z$を$A$へ持ち上げ，同じ記号で表す．
次数$1,2$では$A\to B$が同型なので，$u,v,w$は互いに可換で$uw=v^2$である．
また，$R\to B$は単射である．実際，$R[z]$では
$(\ell z-p)\cap R=0$だからである．従って$R\to A$も単射となる．
$B$の生成元を持ち上げて次数に関する帰納法を行えば，
$A$は$u,v,w,z,g$で生成される．

次数$3$における商写像の核は$\C g$なので，定数$c_u,c_v,c_w,\lambda\in\C$について
\begin{equation}\label{eq:v6-lift-relations}
 [z,x]=c_xg\quad(x=u,v,w),\qquad \ell z-p=\lambda g
\end{equation}
と書ける．$gx=\nu(x)g$でnormal automorphismを定める．

$\lambda=0$なら，$x\in\{u,v,w\}$と$\ell z-p=0$のcommutatorを取って
$c_x\ell g=0$を得る．$g$の右regularityと$\ell\ne0$から$c_x=0$である．
従って全ての境界関係が$A$で成り立ち，$B\to A$は商写像のsectionとなる．
$\nu$は$A_1,A_2$を保つので，これらで生成されるsectionの像も保つ．
よって$B[t;\nu|_B]\to A$，$t\mapsto g$は全射であり，Hilbert級数の一致により同型である．
これは第一の場合である．

以下$\lambda\ne0$とする．このとき$g$を消去できるので$A$は$u,v,w,z$で生成される．
式\eqref{eq:v6-lift-relations}のcommutatorとnormalityから
\begin{equation}\label{eq:v6-rank-one-normal-action}
 \nu(x)-x=\lambda^{-1}c_x\ell\quad(x=u,v,w)
\end{equation}
を得る．ここでは両辺に右から$g$を掛けた等式から，$g$のregularityによって約した．

$R=\C[s^2,st,t^2]$という表示を用いると，$\ell$は二変数二次形式である．
$\C$上の変数変換と非零定数倍により，$\ell=u$または$\ell=v$にできる．
前者は重根，後者は相異なる二根の場合に対応する．

$\ell=u$の場合には，式\eqref{eq:v6-rank-one-normal-action}を
\[
 \nu(u)=(1+\alpha)u,\quad\nu(v)=v+\beta u,
 \quad\nu(w)=w+\gamma u
\]
と書く．$\nu(uw-v^2)=0$を$R$内で計算すると
\[
 \alpha v^2-2\beta uv+
 \bigl((1+\alpha)\gamma-\beta^2\bigr)u^2=0.
\]
$u^2,uv,v^2$は一次独立なので$\alpha=\beta=\gamma=0$である．
従って$c_u=c_v=c_w=0$であり，$R[z]\to A$が全射となる．
Hilbert級数の一致から同型であり，第二の場合を得る．

$\ell=v$の場合には
\[
 \nu(u)=u+\alpha v,\quad\nu(v)=(1+\beta)v,
 \quad\nu(w)=w+\gamma v
\]
と書ける．同様に$\nu(uw-v^2)=0$から
\[
 \gamma uv+\alpha vw+
 \bigl(\alpha\gamma-(1+\beta)^2+1\bigr)v^2=0
\]
となる．従って$\alpha=\gamma=0$，$\beta=0$または$-2$である．
$\beta=0$なら再び第二の場合となる．

$\beta=-2$なら$c_u=c_w=0$，$c_v=-2\lambda$なので，
$u,w$は$z$とも可換であり，
\[
 zv-vz=-2\lambda g=-2(vz-p),\qquad zv+vz=2p.
\]
$R/(v)=\C[u,w]/(uw)$だから，ある$q\in R_2$と$a,b\in\C$について
$p=a u^3+b w^3+vq$と書ける．$z$を$z-q$で置き換えると，
$vz-p$と$[z,x]$は変化せず，$p=a u^3+b w^3$にできる．
ここで$\sigma$と$\delta$を定理の式で定めると，
$\delta(v^2-uw)=(-v)2p+2pv=0$であり，可換関係も保つので，
$\delta$は$R$上の$\sigma$-derivationを定める．
よって$R[z;\sigma,\delta]\to A$は全射である．
このOre拡大は左$R$自由で，Hilbert級数が$(1-t)^{-3}$なので同型である．
$z^2v=vz^2$は$zv+vz=2p$と$zp=pz$から従い，$u,w$も$z$と可換なので$z^2$は中心的である．
また$r=vz-p$は$u,w$と可換で，$rv=-vr$，$rz=-zr$を満たす．
従って第三の候補には実際にregular normalな$r$があり，その商は指定された$B$となる．

最後に有限homological dimensionを排除する．第一の場合には，
補題\ref{lem:split-vertex}と系\ref{cor:split-pd}から$\operatorname{pd}_{\operatorname{Gr}B}k<\infty$が必要となる．
しかし有限最小自由分解があれば$1/H_B(t)$はLaurent多項式となる一方，
\[
 \frac1{H_B(t)}=\frac{(1-t)^2}{1+t+t^2}
\]
は原始三乗根に極を持つので矛盾する．

第二，第三の場合には補題\ref{lem:v6-ore-differential-vertex}を使える．
第二の場合は$z$が中心的で，第三の場合は$z^2$が中心的である．
ところが
\[
 H_R(t)=\frac{1+t}{(1-t)^2},\qquad
 \frac1{H_R(t)}=\frac{(1-t)^2}{1+t}
\]
は$t=-1$に極を持つので，$k_R$の最小自由分解は有限にならない．
従って両場合とも$\qgr A$は無限homological dimensionを持つ．
\end{proof}

\begin{remark}[節点境界の場合]
研究ノートの具体例
\[
 B_{\mathrm{node}}=
 \C[u,v,w,z]/(uw-v^2,\ vz-u^3-w^3)
\]
は$\ell=v$，$p=u^3+w^3$の場合である．
この境界のnormal三次被覆には，可換な錐以外にも
\[
 u,w\text{ central},\qquad v^2=uw,
 \qquad zv+vz=2(u^3+w^3)
\]
という非可換domainが存在する．しかし上の定理はこの被覆も排除するので，
$B_{\mathrm{node}}$のambientを取り替えるだけでは有限homological dimensionを得られない．
これは境界だけの数値・飽和性に対する反例を，元の六公理の反例へ持ち上げる可能性を否定する結果である．
\end{remark}

\begin{lemma}[二次までで生成される可換境界の表示]
\label{lem:v6-commutative-boundary-presentation}
$B$をconnected graded可換$\C$-domainで，
$\dim B_n=3n$ for $n>0$，かつ次数$1,2$で生成されるものとする．
$B$が一次生成されないなら，定理\ref{thm:v6-weighted-cubic-cover}の表示を持つ．
\end{lemma}
\begin{proof}
$\dim B_1=3$，$\dim B_2=6$である．domain内の積部分空間のHopfの次元不等式から
$\dim B_1^2\geq2\dim B_1-1=5$である．
$\dim B_1^2=6$なら次数$2$の独立な生成元は不要なので，仮定により一次生成となる．
従って$\dim B_1^2=5$である．

$B_1$の基底と独立な$z\in B_2$を選ぶと，重み$(1,1,1,2)$の多項式環から$B$への全射を得る．
二次の核は一つの非零二次形式$q$で張られる．
$B$はdomainであり$B_1$の元は一次独立なので，$q$は一次式の積に分解しない．
$\C$上の三変数二次形式なのでrankは$3$であり，変数変換によって$q=uw-v^2$とできる．

$R[z]$から$B$への全射を考えると，その三次部分の次元は$10$から$9$へ下がるので，
ある非零$F\in R[z]_3$が核を張る．$F$は$\ell z-p$と書ける．
$R[z]$はdomainだから$F$はregularであり，$R[z]/(F)$と$B$のHilbert級数は一致する．
従って全次数で$B=R[z]/(F)$である．

$\ell=0$は不可能である．実際，$R=\C[s^2,st,t^2]$の下で
$0\ne p\in R_3$は二変数六次形式であり，$\C$上の因数分解を二つと四つに分ければ
$p=p_1p_2$，$0\ne p_1\in R_1$，$0\ne p_2\in R_2$と書ける．
$R[z]/(p)$では$p_1,p_2$はいずれも非零なので，domain性に反する．
よって$\ell\ne0$であり，所要の表示を得る．
\end{proof}

\begin{corollary}[二次までで生成される可換三次境界からの認識]
\label{cor:v6-commutative-degree-two-recognition}
$A$をconnected locally finite graded右noetherian代数，$g\in A_3$をregular normal元とする．
$B=A/(g)$が可換domainで，$\dim B_n=3n$ for $n>0$，かつ次数$1,2$で生成されると仮定する．
$\qgr A$が有限homological dimensionを持てば
\[
 A\simeq\C[x_0,x_1,x_2],\qquad\deg x_i=1.
\]
特に$B$は一次生成され，$g$は通常の斉次三次式に対応する．
\end{corollary}
\begin{proof}
一次生成されない場合は前補題と定理\ref{thm:v6-weighted-cubic-cover}に矛盾する．
従って$B$は一次生成される．$B$は三変数多項式環の一つの三次式$f$による商である．
実際，二次までの核は零，三次の核は一次元であり，$\C[x_0,x_1,x_2]/(f)$と
$B$のHilbert級数は一致する．

$B_1$の$A$へのliftは互いに可換なので，$\C[x_0,x_1,x_2]\to A$が定まる．
$f$の像は$\lambda g$である．$\lambda=0$なら$B\to A$がsectionとなり，
$A=B[g;\nu]$となるので，分裂Ore拡大の頂点Extにより有限homological dimensionに反する．
$\lambda\ne0$なら$A$は次数$1$で生成され，三変数多項式環からの全射は
$H_A=(1-t)^{-3}$により同型となる．
\end{proof}

\begin{remark}[二つの被覆認識の適用範囲]
上の系は標準Hilbert列，可換性，次数$2$までの生成を明示的に仮定するが，
六公理やambientのSerre関手は仮定しない．従って$B_{\mathrm{node}}$の全被覆の排除は，
より弱いambientの仮定の下でも成立する．後の定理\ref{thm:v7-commutative-cubic-recognition}は
六公理とexactな三乗根を使い，可換な三次domain境界の生成次数・飽和性に関する仮定を全て除く．
一般の非可換境界は，いずれの可換境界認識の結論にも含めない．
\end{remark}

\subsection{Balancedなstacky節点の排除}
\label{subsec:stacky-node-obstruction}

次の方法は重み付き完全交叉表示を必要としない．曲線stackへの表示を明示的に仮定する．
balancedなcyclic nodeの局所表示はLekili--Polishchuk\cite[Introduction, p.~3]{LPstack}の規約と一致する．
同文献のAuslander order上の単純対象と，以下で計算する構造層上の単純対象は区別する．
以下の局所Ext計算および有限homological dimensionとの非両立性には直接の証明を与える．

\begin{corollary}[Balancedなstacky節点の排除]
\label{cor:v6-balanced-node-obstruction}
補題\ref{lem:v6-normal-ext-periodicity}の設定と有限homological dimensionを仮定し，$(D,T^e)$はAZ生成条件を満たすとする．
さらに構造対象を保つ同値
\[
 (\qgr B,D)\simeq(\operatorname{coh}\mathcal C,\OO_{\mathcal C})
\]
があり，$\mathcal C$は$\C$上のproper separated tame Deligne--Mumford curveで，
その特異点は全てordinary node，各nodeのstabilizerはcyclicかつbalancedであるとする．
すなわち，形式的局所表示は
\[
 [\operatorname{Spec}\C[[x,y]]/(xy)\,/\,\mu_r],\qquad
 \zeta\cdot(x,y)=(\zeta x,\zeta^{-1}y)
\]
である．このとき各nodeで$r=1$である．
ここでは$T$または$\Phi$がnodeやそのcharacterを個別に固定することを仮定しない．
\end{corollary}
\begin{proof}
局所環$Q=\C[[x,y]]/(xy)$について，character $a\in\Z/r\Z$の単純対象を$S_a$とする．
その最小equivariant自由分解は，次数$0$の項$Q\otimes\chi^a$に続いて
\[
 P_{2j}=(Q\otimes\chi^a)^{\oplus2}\quad(j\geq1),\qquad
 P_{2j+1}=(Q\otimes\chi^{a+1})\oplus(Q\otimes\chi^{a-1})\quad(j\geq0)
\]
となる．最初の微分は$(x\ y)$で，その後の微分は
$\operatorname{diag}(y,x)$，$\operatorname{diag}(x,y)$を交互に使う．
$(x)\cap(y)=0$，$\operatorname{ann}(x)=(y)$，$\operatorname{ann}(y)=(x)$から完全性が直ちに分かる．
単純対象への$\Hom$を施すと微分は零であり，$\mu_r$の不変部分を取ることはexactなので，
\[
 \Ext^{2j}(S_a,S_b)=
 \begin{cases}\C^2&a=b,\\0&a\ne b\end{cases}\quad(j\geq1),
\]
\[
 \dim\Ext^{2j+1}(S_a,S_b)
 =\mathbf1_{b=a+1}+\mathbf1_{b=a-1}\quad(j\geq0)
\]
を得る．$r=1,2$では二つのcharacterが一致し得るが，偶数Extの対角性は変わらない．
異なるnode間のExtは零である．点支持の局所Ext層には高次cohomologyがないので，
この計算は$\operatorname{coh}\mathcal C$の大域Extを与える．

node以外の点では局所環はregularであり，tamenessにより単純対象のExtは次数$2$以上で零である．
従って全nodeのcharacter単純対象の集合$\Sigma$は，
任意に高い次数でExtを持つ単純対象の集合として圏から内在的に識別できる．
特に全てのexact autoequivalence，したがって$T$と$\Phi$で保たれる．
命題\ref{prop:v6-diagonal-simple-obstruction}から全ての$S_a\in\Sigma$に対して
$\Hom(\OO_{\mathcal C},S_a)\ne0$となる．
しかしこの空間は$S_a$のstabilizer不変部分であり，$a=0$の場合だけ非零である．
従って非自明なcharacterは存在せず，$r=1$である．
\end{proof}

\begin{remark}[三次境界への適用条件]
六公理の下で$e=3$なら，$(D,T^3)$のAZ生成条件はambientから降りる．
実際，$i_*F$への有限直和$\bigoplus_j\pi_AA(-3n_j)\twoheadrightarrow i_*F$の各射は，
normal元$g$を殺すので$\pi_BB(-3n_j)$を経由する．
この因子分解はtailsで，$ga=\nu(a)g$と標的への$g$の零作用を使って確認できる．
従って曲線stack表示を構造対象とともに得られれば，上の系は適用できる．
ただし，現時点でこの表示を任意の非可換domain境界から構成したわけではない．
また，表示が$\qgr B$の抽象的な同値だけで，$D$が$\OO_{\mathcal C}$に対応しない場合には，
非自明characterへの$\Hom$が消えるという最後の議論を使えない．
\end{remark}
\subsection{反標準Hilbert列による単一節点のbalance}

次の結果では，coarse curveと正規化の形，およびshiftの幾何学的表示を追加仮定する．
一般の三次domain境界がこの形であるという分類は得ていない．
特に，dualizing sheafが自明であることだけからbalanceを結論しない．
nodeの二つの接方向に任意のunit characterが作用しても，対数微分の不変性は成立するためである．

\begin{lemma}[単一orbifold nodeのRiemann--Roch公式]\label{lem:orbinode-riemann-roch}
$\mathcal X$をintegral proper tame DM curveとする．そのcoarse curveはintegral rational nodal curve，
stack構造はその節点だけにあり，formal nodeは$\mu_r$が接方向にunit character $a,b$で作用するものとする．
正規化$\nu:\widetilde{\mathcal X}\to\mathcal X$について，$\widetilde{\mathcal X}$はcoarse $\mathbb P^1$を持ち，節点上の二点だけにstabilizer $\mu_r$を持つとする．
$\mathcal N$を$\mathcal X$上の線束，$z\in\Z/r\Z$を節点fiberのcharacter，
$\delta=\deg_{\widetilde{\mathcal X}}\nu^*\mathcal N$をorbifold degreeとする．
$\alpha(z),\beta(z)\in\{0,\ldots,r-1\}$を
\[
 a\alpha(z)+z\equiv0,\qquad b\beta(z)+z\equiv0\pmod r
\]
で定めると
\begin{equation}\label{eq:v6-orbinode-RR}
 \chi(\mathcal X,\mathcal N)
 =\delta+1-\frac{\alpha(z)+\beta(z)}r-\mathbf1_{z=0}.
\end{equation}
\end{lemma}

\begin{proof}
$p:\widetilde{\mathcal X}\to\mathbb P^1$をcoarse mapとする．
一つのstack点のchartで，局所座標$u$のweightを$a$，線束fiberのweightを$z$とすれば，
不変sectionの最小単項式は$u^{\alpha(z)}$である．
従って$p_*\nu^*\mathcal N$は$\mathbb P^1$上の線束で，そのdegreeは
\[
 \delta-\frac{\alpha(z)+\beta(z)}r
\]
である．有限stabilizerはlinearly reductiveなので$p_*$は完全で，高次直像は零である．
従って正規化上のEuler標数は上のdegreeに$1$を加えたものとなる．

nodeの正規化完全列を$\mathcal N$でtensorすると
\[
 0\longrightarrow\mathcal N\longrightarrow\nu_*\nu^*\mathcal N
 \longrightarrow i_*\chi_z\longrightarrow0
\]
を得る．最後の項は節点の両branchの値の差である．
そのglobal sectionの次元は$z=0$のときだけ$1$で，高次cohomologyは零だから，
表示された公式が従う．
\end{proof}

\begin{proposition}[反標準Hilbert列とAZ3からのbalance]\label{prop:orbinode-balance}
前補題の設定とし，$S=\mathcal L\otimes\sigma^*(-)$とする．
$\sigma$には二つのbranchの交換も許す．
$(\OO_{\mathcal X},S^3)$がAZ1を満たし，全$m>0$で
\[
 h^0(\mathcal X,S^{3m}\OO)=9m,\qquad H^1(\mathcal X,S^{3m}\OO)=0
\]
が成立すると仮定する．このとき
\[
 \deg\nu^*\mathcal L=3,\qquad b\equiv-a\pmod r.
\]
従って系\ref{cor:v6-balanced-node-obstruction}のnormal被覆の仮定も満たす場合には$r=1$である．
\end{proposition}

\begin{proof}
$\mathcal N_n=S^n\OO$と置く．degreeは自己同型で不変なので
$\deg\nu^*\mathcal N_n=n\ell$，$\ell=\deg\nu^*\mathcal L$である．
\eqref{eq:v6-orbinode-RR}の補正項は有界だから，$n=3m$で比較して$\ell=3$を得る．
節点のfiberにおける$S$の作用は，characterの有限集合の置換$\psi$である．
$\mathcal N_n$のfiber characterは$\psi^n(0)$となる．
AZ1から$\{\psi^{-3m}(0):m>0\}$は全characterを尽くす．
実際，尽くさなければ，欠けたcharacterのsimpleへの全てのHomが零となる．
有限置換なので正の$3m$の軌道も同じ集合を尽くす．

従って\eqref{eq:v6-orbinode-RR}と仮定から，全$z\ne0$について
$\alpha(z)+\beta(z)=r$を得る．$r>1$なら$z=-a$を選ぶ．
$a$はunitなので$z\ne0$であり，$\alpha(-a)=1$，従って$\beta(-a)=r-1$である．
$b(r-1)-a\equiv0$から$b\equiv-a$が従う．$r=1$は自明である．
これで系\ref{cor:v6-balanced-node-obstruction}を適用できる．
\end{proof}

\begin{remark}[この排除で解決していない範囲]
系\ref{cor:cubic-boundary-duality}と系\ref{cor:cubic-saturated-hilbert}は，
構造対象とshiftを保つ上の表示がある場合，命題\ref{prop:orbinode-balance}の
反標準Hilbert値と正次数の$H^1$消滅を与える．AZ1もambientから降りる．
したがってこの正規化型の非自明なstacky nodeは排除される．
一方，coarse cusp，tacnode等のatlas，非可換orderのheart，有限section欠損を含む
一般の境界をこの正規化型へ分類したわけではない．
研究ノートにはcyclic stabilizerと無限位数自己同型を持つ明示的な境界族を置く．
この族の全体にも系\ref{cor:v6-balanced-node-obstruction}を適用できる．
\end{remark}

\subsection{高次Ext周期性が強制する局所平面性と安定化群の指標}\label{subsec:v7-local-geometry}

この節では，normal境界が構造対象を保って可換な曲線stackとして表示される場合を扱う．
shiftの幾何学的表示は仮定しない．有限ホモロジー次元は，局所特異点が
hypersurfaceであること自体にも制約を与える．

\begin{lemma}[構造対象の双対性と双対化層]\label{lem:v7-dualizing-trivial}
$\mathcal X$をprojective coarse curveを持つproper separated tameな純一次元Deligne--Mumford stackとする．
任意の$F\in\operatorname{coh}\mathcal X$について自然な同型
\[
 \Ext^1_{\mathcal X}(\OO_{\mathcal X},F)
 \simeq\Hom_{\mathcal X}(F,\OO_{\mathcal X})^*
\]
があるなら，双対化層は$\omega_{\mathcal X}\simeq\OO_{\mathcal X}$である．
\end{lemma}
\begin{proof}
Levchenkoの双対化層の存在定理\cite[Proposition 2]{Levchenko}は，全てのcoherent $F$について
\[
 \Hom(F,\omega_{\mathcal X})\simeq H^1(\mathcal X,F)^*
\]
を自然に与える．$\Ext^1(\OO,F)=H^1(F)$と仮定の双対性から
$\Hom(F,\omega_{\mathcal X})\simeq\Hom(F,\OO_{\mathcal X})$を得る．
Yonedaの補題によって結論する．この議論は，特異な曲線の有界導来圏全体に
Serre functorがあることを必要としない．
\end{proof}

\begin{proposition}[Normal周期性からの局所平面性]\label{prop:v7-local-planarity}
補題\ref{lem:v6-normal-ext-periodicity}の設定と
$\operatorname{hdim}\qgr A<\infty$を仮定する．さらに構造対象を保つ同値
\[
 (\qgr B,D)\simeq(\operatorname{coh}\mathcal X,\OO_{\mathcal X})
\]
があり，$\mathcal X$はprojective coarse curveを持つproper separated tameな
reduced純一次元Deligne--Mumford stackとする．このとき$\mathcal X$の全ての
完備局所atlasはregularまたはplane hypersurfaceである．
\end{proposition}
\begin{proof}
$\mathcal X$の特異点上の単純対象の同型類の集合を$\Sigma$とする．
特異点の数と各stabilizerの既約表現の数は有限なので，$\Sigma$は有限集合である．
これは圏から内在的に識別できる．実際，単純対象$S$が$\Sigma$に入ることと，
\[
 \{n\geq0:\text{ある単純対象 }S'\text{ に対して }\Ext^n(S,S')\ne0\}
\]
が非有界であることは同値である．regularな一次元atlasでは，tamenessにより
次数$2$以上の局所Extは消える．非regularな局所環の剰余体は無限projective dimensionを持ち，
各次数のTor表現を有限個の既約表現へ分解すれば，逆向きが従う．
異なる閉点の単純対象間のExtは零であり，点支持の局所Ext層の高次cohomologyは消える．
従って以上の局所計算は大域Extでも成立する．

$\Theta=\Phi T^{-e}$は$\Sigma$を置換する．高次Ext周期性を反復すると，
$S,S'\in\Sigma$について$\dim\Ext^n(S,S')$は$n$に関して有界である．
実際，$n\geq d$を$2$ずつ下げ，第一変数を$\Theta$で置き換えれば，
次数$d,d+1$と有限集合$\Sigma$上のExtの次元だけで全てを抑えられる．

一つの特異点の完備局所表示を$[\operatorname{Spec}R/G]$とする．
$G$の既約表現$\rho$に対応する単純対象を$S_\rho$と書く．
$G$がlinearly reductiveであることから
\[
 \dim_{\C}\Tor_n^R(k,k)
 =\sum_{\rho\in\operatorname{Irr}G}(\dim\rho)
       \dim\Ext^n_{[\operatorname{Spec}R/G]}(S_{\mathbf1},S_\rho).
\]
従って$k$の$R$上のBetti数は有界である．
Gulliksen\cite[Proposition 6, p.~300]{Gulliksen1967}によれば，これは
\[
 \operatorname{depth}R\geq\operatorname{edim}R-1
\]
と同値である．同論文の$\operatorname{codim}R$はこの$\operatorname{depth}R$を意味する．
$R$はreduced純一次元なので$\operatorname{depth}R=1$であり，
$\operatorname{edim}R\leq2$を得る．
完備等標数局所環の表示$\C[[x_1,\ldots,x_v]]\twoheadrightarrow R$において，
$v=1$ならregular，$v=2$なら核は高さ$1$のradical idealであり，
正則局所UFDの性質から一つの非零元$f$で生成される．
従って非regularな場合は$R=\C[[x,y]]/(f)$である．
\end{proof}

\begin{lemma}[平面特異点のequivariant剰余体分解]\label{lem:v7-plane-tate}
$Q=\C[[x,y]]$に有限群$G$が線形に作用し，
$0\ne f\in(x,y)^2$がcharacter $\chi$を持つsemi-invariantであるとする．
$R=Q/(f)$，$V=\langle x,y\rangle$と置く．$R$上の剰余体の最小$G$-equivariant自由分解は，
表現として
\[
 P_0=R,\qquad P_{2m+1}=R\otimes V\otimes\chi^m\quad(m\geq0),
\]
\[
 P_{2m}=R\otimes
 \bigl(\chi^m\oplus(\det V)\otimes\chi^{m-1}\bigr)\quad(m\geq1)
\]
という項を持つ．また双対化加群のfiber characterは$(\det V)\chi^{-1}$である．
\end{lemma}
\begin{proof}
$f=xu+yv$，$u,v\in(x,y)$と書く．有限群で平均を取ることにより，
この表示に対応する写像$\chi\to Q\otimes V$を$G$-equivariantに選べる．
最初の微分と次の二つの微分は，基底を選べば
\[
 d_1=(x\ y),\qquad
 M=\begin{pmatrix}-y&u\\x&v\end{pmatrix},\qquad
 N=\begin{pmatrix}v&-u\\-x&-y\end{pmatrix}
\]
である．$MN=NM=-fI_2$なので，$R$上では$M,N$を交互に並べた複体となる．
$d_1$の核が$M$の像であることは，$Q$で
$xa+yb=fc$を持ち上げ，$x(a-uc)+y(b-vc)=0$にKoszulの関係を使えば分かる．
$\ker M=\operatorname{im}N$も直接従う．実際，$M\widetilde w=fz$を$Q$で満たすliftを選び，
$N$を掛けると$-f\widetilde w=fNz$なので，$Q$上で$f$を消去できる．
$\ker N=\operatorname{im}M$も同様である．全ての係数は極大idealに入るから分解は最小である．

第一のsyzygyは$\det V$，第二のsyzygyは$\chi$を持ち，
二周期ごとに$\chi$をtensorする．これが表示された項の表現を与える．
同じ記述は，次数$1$に$V$，次数$2$に$\chi$を置くequivariant Tate分解でも得られる．
最後に$0\to Q\otimes\chi\xrightarrow{f}Q\to R\to0$へ
$\Hom_Q(-,\omega_Q)$を適用すれば，$\omega_R$のfiberは
$(\det V)\chi^{-1}$である．これは既知のlci双対化公式
\cite[Theorem 3.8]{NironiDuality}の余次元$1$の場合でもある．
\end{proof}

\begin{theorem}[特異点の方程式指標は忠実である]\label{thm:v7-faithful-equation-character}
命題\ref{prop:v7-local-planarity}の設定で，さらに$(D,T^e)$がAZ生成条件を満たし，
$\Ext^1(D,F)\simeq\Hom(F,D)^*$が全ての$F$に自然に成立するとする．
このとき各特異点のstabilizer $G$はcyclicである．
その完備局所表示を$[\operatorname{Spec}\C[[x,y]]/(f)/G]$とし，
$f$のsemi-invariant characterを$\chi$とすると，$\chi:G\to\C^*$は忠実である．
特に$G=\mu_r$，$x,y,f$のweightがそれぞれ$a,b,c$なら
\[
 c\equiv a+b\pmod r,\qquad \gcd(r,c)=1.
\]
この結論には，$T$または$\Phi$が幾何学的であることや，各点を固定することを仮定しない．
\end{theorem}
\begin{proof}
局所平面性は前命題による．有限群の作用を平均すれば，cotangent spaceの
$G$-stableなliftから形式座標を選べる．その座標で$G$の作用は線形となる．
またprincipalな関係idealの生成元を平均し直せば，$f$をsemi-invariantに選べる．
補題\ref{lem:v7-dualizing-trivial}より$\omega_{\mathcal X}\simeq\OO_{\mathcal X}$なので，
前補題の双対化公式は$\det V=\chi$を与える．
従って各既約表現$\rho,\sigma$について
\[
 \Ext^{2m}(S_\rho,S_\sigma)
 \simeq\Hom_G(\rho\otimes\chi^m,\sigma)^{\oplus2}\quad(m\geq1).
\]
異なる閉点のExtは零である．

特異点上の単純対象の有限集合$\Sigma$で
$K(S_\rho)=S_{\rho\otimes\chi}$と定める．
$K(S)$は$\Ext^2(S,-)$が非零となる唯一の$\Sigma$内の単純対象として特徴付けられる．
従って全てのexact autoequivalenceは$K$と可換であり，特に$T$と$\Phi$もそうである．
高次Ext周期性により，十分大きい$m$で
\[
 \Ext^{2m+2}(S,S')\simeq
 \Ext^{2m}(\Phi T^{-e}S,S')
\]
だから，非零になる第二変数を比較して
$K^{m+1}S=K^m\Phi T^{-e}S$を得る．よって$\Sigma$上で
\[
 \Phi T^{-e}=K,\qquad T^e=\Phi K^{-1}.
\]
$\Phi D\simeq D$と可換性から
\[
 \Hom(D,T^{en}S)\simeq\Hom(D,K^{-n}S)\quad(n\geq0).
\]
AZ生成条件によれば，任意の$S_\rho$について左辺が非零となる$n$が存在する．
一方，右辺が非零となるのは$\rho\otimes\chi^{-n}$が自明表現の場合だけである．
従って$G$の全ての既約表現は$\chi$の冪である．
既約表現の核の共通部分は自明なので，$\ker\chi=1$であり，
$G$は$\C^*$の有限部分群，すなわちcyclic群である．
\end{proof}

\begin{corollary}[$A_k$型の局所制約]\label{cor:v7-ak-character}
前定理の設定で，特異点の表示が
\[
 R=\C[[x,y]]/(y^2-x^N),\qquad N\geq2,
\]
かつ$\mu_r$の作用が$x,y$のweight $a,b$で与えられるとする．
このとき
\[
 r\mid(N-2),\qquad r\text{ は奇数},\qquad
 \gcd(a,r)=1,\qquad b\equiv a\pmod r.
\]
特にordinary cusp ($N=3$) とtacnode ($N=4$) ではstabilizerは自明である．
\end{corollary}
\begin{proof}
方程式のsemi-invarianceと双対化指標の自明性から
\[
 2b\equiv Na\equiv a+b\pmod r
\]
を得る．第一と第三の等式より$b\equiv a$であり，従って
$(N-2)a\equiv0$，$c\equiv2a$である．
$\gcd(r,c)=1$から$r$は奇数，$a$はunitであり，$r\mid(N-2)$となる．
\end{proof}

\begin{remark}[Unbalanced nodeと残る範囲]
ordinary node $xy=0$では$\chi$のweightは$a+b$なので，必要条件は
$\gcd(r,a+b)=1$である．balancedの場合$a+b=0$は$r=1$を強制し，
系\ref{cor:v6-balanced-node-obstruction}を含む．
一方，$a+b$がunitであるunbalanced nodeは，この局所条件だけでは排除しない．
また$N\geq5$の$A_{N-1}$型には奇数$r\mid(N-2)$が残り得る．
この定理は局所表示への必要条件であり，これらの型のglobalな境界またはnormal被覆の存在を主張しない．
\end{remark}

\begin{example}[Coarse cuspを持つtacnode quotientの排除]\label{ex:v7-tacnode-obstruction}
\[
 R=\C[[x,y]]/(y^2-x^4),\qquad
 \mu_2:(x,y)\longmapsto(-x,-y)
\]
を考える．不変環は
\[
 R^{\mu_2}=\C[[u,v]]/(v^2-u^3),\qquad u=x^2,\quad v=xy,
\]
であり，coarse germはordinary cuspである．双対化層のfiber characterは自明である．
この局所stackと双対化指標の計算はNironi\cite[Example 3.7]{NironiDuality}に既にある．
しかし方程式$f=y^2-x^4$のcharacterは自明なので，
定理\ref{thm:v7-faithful-equation-character}により，この局所型は上のnormal境界に現れない．
直接には二つのcharacter単純対象$S_0,S_1$について
\[
 \Ext^{2m}(S_i,S_j)=\C^2\delta_{ij}\quad(m\geq1)
\]
となり，AZ生成条件との矛盾が検出される．
従ってcoarse curveがcuspであることや双対化層の自明性だけではこのatlasを排除できず，
高次ExtとAZ生成条件を用いる必要がある．
\end{example}

\subsection{可換な三次境界の構造対象付き分類}
\label{subsec:v7-commutative-boundary}

可換な商環については，生成次数を制限せずに元の$qgr$，構造対象，shiftを
同時に幾何化できる．以下では高Veroneseの$qgr$を元の$qgr$と同一視しない．
非標準次数の可換環から直接作るweighted projective stackを用い，
normal被覆の有限homological dimensionによってそのstabilizerを除く．

\begin{lemma}[正規化とconductorの次数公式]\label{lem:v7-stack-conductor-degree}
$\mathcal X$を$\C$上のintegral proper tame Deligne--Mumford curveとし，
generic stabilizerは自明，$\omega_{\mathcal X}\simeq\OO_{\mathcal X}$とする．
正規化を$\eta:\widetilde{\mathcal X}\to\mathcal X$，そのcoarse curveの種数を$\gamma$とする．
このとき$\mathcal X$のsmooth locusには非自明なstabilizerはない．
特異点の正規化における逆像を$P_1,\ldots,P_N$，
$P_i$のstabilizerの位数を$h_i$，conductorの一様化元に関する指数を$c_i\geq1$とすると
\begin{equation}\label{eq:v7-conductor-degree}
 2-2\gamma
 =\sum_{i=1}^N\left(1+\frac{c_i-1}{h_i}\right).
\end{equation}
\end{lemma}
\begin{proof}
smoothなstack点の完備局所表示は$[\operatorname{Spec}\C[[t]]/\mu_h]$である．
generic stabilizerが自明なので，$t$のcharacterは忠実である．
双対化層のfiberは$dt$と同じcharacterを持つから，
$\omega_{\mathcal X}\simeq\OO_{\mathcal X}$は$h=1$を強制する．

局所的に$R$をGorenstein curve ring，$\overline R$をその正規化とすれば，有限射の双対性から
\[
 \omega_{\overline R}
 =\operatorname{Hom}_R(\overline R,\omega_R)
 =\mathfrak c\,\omega_R,
 \qquad
 \mathfrak c=\operatorname{Hom}_R(\overline R,R)
\]
となる．$\mathfrak c$はconductor idealであり，正規化上では
$\bigoplus_jt_j^{c_j}\C[[t_j]]$である．この同一視は有限群の作用と自然に両立する．
従ってstack上でも
\[
 \omega_{\widetilde{\mathcal X}}
 \simeq\eta^*\omega_{\mathcal X}\otimes\OO_{\widetilde{\mathcal X}}(-\Delta),
 \qquad \Delta=\sum_i c_iP_i
\]
を得る．有限射に対する双対性と，projective CM stackの双対化層については
\cite[Corollary 2.30, Proposition 2.31]{NironiDuality}も参照されたい．

滑らかなeffective stack curveのcanonical degreeは
\[
 \deg\omega_{\widetilde{\mathcal X}}
 =2\gamma-2+\sum_i(1-h_i^{-1})
\]
であり，$\deg P_i=h_i^{-1}$である．この公式は局所表示
$u=t^{h_i}$における$du=h_it^{h_i-1}dt$から従う．
上の双対化層の同型の次数を取れば，所要の等式になる．
\end{proof}

\begin{lemma}[整数次数の忠実な偏極を持つ境界stackの排除]
\label{lem:v8-integral-degree-stack}
$\mathcal X$をprojective coarse curveを持つintegral proper tame DM curveとし，generic stabilizerは自明，
$\omega_{\mathcal X}\simeq\OO_{\mathcal X}$とする．
各完備局所atlasはregularまたはplane hypersurfaceであり，各特異点の方程式指標は忠実であると仮定する．
さらに，すべてのstabilizer上で忠実なfiber characterを持つ線束$\mathcal M$が存在し，
$\deg\mathcal M\in\Z$とする．
このとき$\mathcal X$は通常のintegral projective Gorenstein curveであり，$p_a(\mathcal X)=1$である．
\end{lemma}
\begin{proof}
補題\ref{lem:v7-stack-conductor-degree}よりsmooth locusのstabilizerは自明であり，
正規化のcoarse genusを$\gamma$，特異点の正規化の逆像を$P_1,\dots,P_N$，
そのstabilizerの位数とconductor指数を$h_i,c_i$とすれば
\[
 2-2\gamma=\sum_{i=1}^N\left(1+\frac{c_i-1}{h_i}\right).
\]
右辺の各項は$1$以上である．特異点がなければ$\gamma=1$であり結論を得る．
特異点があれば$\gamma=0$，$N=1$または$2$である．

$N=2$なら$c_1=c_2=1$である．
局所atlasを$R\subset\overline R=\bigoplus_{j=1}^b\C[[t_j]]$と書くと，
conductor指数1と局所環であることから
\[
 R=\C(1,\dots,1)+\bigoplus_j t_j\C[[t_j]].
\]
embedding dimensionは$b$なので，局所平面性と特異性から$b=2$であり，ordinary nodeである．
双対化生成元$(dt_1/t_1,-dt_2/t_2)$への作用より，branchを交換する作用は許されない．
従って$N=2$の二点は一つのnodeの二本のbranchである．
そのstabilizerを$\mu_r$とし，$r>1$と仮定する．
二つの一様化元のweightを$a,b\in(\Z/r\Z)^\times$，$\mathcal M$のfiber characterを
$z\in(\Z/r\Z)^\times$とする．
\[
 a\alpha+z\equiv b\beta+z\equiv0\pmod r,
 \qquad 1\leq\alpha,\beta\leq r-1
\]
と置くと，正規化のcoarse map $p$について
\[
 \deg p_*\eta^*\mathcal M
 =\deg\mathcal M-\frac{\alpha+\beta}{r}\in\Z.
\]
$\deg\mathcal M\in\Z$だから$\alpha+\beta=r$であり，$b=-a\pmod r$となる．
nodeの方程式$xy$のcharacterは$a+b=0$なので，その忠実性から$r=1$でなければならない．

$N=1$なら，唯一の正規化点のstabilizerの位数を$h$として$c=h+1$となる．
$h>1$なら$\mathcal M$の忠実性から
\[
 \deg p_*\eta^*\mathcal M=\deg\mathcal M-\frac{\alpha}{h}\in\Z,
 \qquad 1\leq\alpha\leq h-1
\]
となり，整数次数に矛盾する．よって$h=1$，$c=2$である．
特異点のstabilizerは，方程式指標の忠実性から$\mu_r$である．
それはatlasのbranchに自由推移的に作用するので，branch数は$r$，各conductor指数は2である．
Gorenstein conductor公式
\[
 2\delta(R)=\dim_\C(\overline R/\mathfrak c)=\sum_jc_j
\]
から$\delta=r$となり，局所種数$\delta-r+1$は1である．
この公式は$0\to R\to\overline R\to Q\to0$を$\operatorname{Hom}_R(-,R)$で双対化し，
$\operatorname{Ext}^1_R(\overline R,R)=0$と
$\operatorname{length}\operatorname{Ext}^1_R(Q,R)=\operatorname{length}Q$を用いて
$\operatorname{length}(R/\mathfrak c)=\delta$を得ることによる．
Smythの分類\cite[Appendix A, Proposition A.3]{SmythGenusOne}より，atlasはelliptic $r$-fold pointである．
$r\geq4$ならembedding dimension $r-1\geq3$となるので，局所平面性に反する．

$r=2$ならordinary tacnodeである．equivariantな形式座標でその最低次項は一つの線形形式の平方であり，
位数2の作用ではこの平方のcharacterは自明である．従ってsemi-invariantな方程式のcharacter全体も自明であり，忠実性に反する．

$r=3$ならordinary triple pointであり，三本の異なる接線を$\mu_3$が巡回する．
$\gamma$を群の生成元とし，接線の一つを定義する線形形式を$\ell$とする．
有限群作用を形式的に線形化すれば，$\gamma^3=1$がcotangent space上で厳密に成立する．
三次接錐の方程式は
\[
 \ell\,\gamma(\ell)\,\gamma^2(\ell)
\]
の非零scalar倍であり，これは$\gamma$で不変である．従ってsemi-invariant方程式のcharacterは自明であり，やはり忠実性に反する．
以上より$r=1$だけが残り，ordinary cuspである．

全stabilizerが自明となったので$\mathcal X$はprojective coarse curveと一致する．
$\omega_{\mathcal X}\simeq\OO_{\mathcal X}$とintegral性から$H^0(\OO)=H^1(\OO)=\C$であり，$p_a=1$である．
\end{proof}

\begin{theorem}[可換三次境界の飽和環は平面三次曲線の環である]
\label{thm:v7-commutative-boundary-geometry}
六公理とexactな三乗根$s$を仮定し，$0\ne g\in A_3$をnormal元とする．
$B=A/(g)$が可換domainなら，integral projective curve $Y$と次数$3$の線束$\mathcal L$が存在し，
\begin{equation}\label{eq:v7-commutative-boundary-triple}
 (\qgr B,\pi_BB,(1))
 \simeq(\operatorname{coh}Y,\OO_Y,(-)\otimes\mathcal L),
 \qquad p_a(Y)=1
\end{equation}
となる．従って全次数のsection環は
\[
 \widehat B=\bigoplus_{n\geq0}H^0(Y,\mathcal L^n)
 \simeq\C[x_0,x_1,x_2]/(f),
 \qquad \deg x_i=1,\quad\deg f=3,
\]
であり，$f$は既約である．特に
\[
 \widehat b_n=3n\quad(n>0),\qquad
 B\hookrightarrow\widehat B\text{ の余核は有限次元},\qquad
 B^{(3)}=\widehat B^{(3)}
\]
が成立する．ここでは$B$の生成次数，飽和性，normal性を追加仮定しない．
\end{theorem}
\begin{proof}
\textbf{構造対象とshiftを保つstack表示．}
$B$はconnected locally finiteな可換noetherian graded環なので有限生成である．
三次境界のHilbert評価から$\operatorname{GKdim}B=\dim B=2$である．
斉次生成元の次数を$d_1,\ldots,d_q$とし，
\[
 U=\operatorname{Spec}B\setminus\{B_+\},\qquad
 \mathcal X=[U/\mathbb G_m],\qquad\mathcal L=\OO_{\mathcal X}(1)
\]
と置く．$\mathcal X$はweighted projective stack
$\mathbb P(d_1,\ldots,d_q)$の閉部分stackであり，integral proper separated tame DM curveである．
stabilizerは$\mathbb G_m$の有限部分群であり，$\mathcal L$のfiberへの作用は忠実である．
系\ref{cor:cubic-normal-dense-support}より正次数supportのgcdは$1$なので，generic stabilizerは自明である．

graded $B$加群は$\operatorname{Spec}B$上の$\mathbb G_m$-equivariant coherent sheafに対応する．
$U$への制限の核は頂点支持の加群，すなわち有限次元加群である．
$U$上のequivariant coherent sheafの開埋め込みによる直像はgraded準連接加群である．
$U$の有限個の斉次主開集合上の生成元を分母を払って持ち上げ，
有限個の斉次成分を選べば，$\operatorname{Spec}B$上のcoherentなequivariant延長を得る．
従って制限によるSerre商は
\[
 (\qgr B,\pi_BB,(1))
 \simeq(\operatorname{coh}\mathcal X,\OO_{\mathcal X},(-)\otimes\mathcal L)
\]
を与える．これは元のgradingに対する同値であり，Veroneseへの置き換えを含まない．

\textbf{双対化層，局所平面性，偏極の次数．}
境界双対性から，補題\ref{lem:v7-dualizing-trivial}により
$\omega_{\mathcal X}\simeq\OO_{\mathcal X}$である．特に$\mathcal X$はGorensteinである．
また$H^0(\mathcal X,\OO)=H^1(\mathcal X,\OO)=\C$である．
命題\ref{prop:v7-local-planarity}から各完備局所atlasはregularまたはplane hypersurfaceである．
ここで使う有限homological dimensionはambientの$\qgr A$のものである．

十分大きく可除な整数$M$に対し，$\mathcal L^M$はcoarse curve
$\operatorname{Proj}B$上のample線束$\OO_{\operatorname{Proj}B}(M)$へ降りる．
正規化と有限長のconductor商を使うRiemann--Rochから
\[
 \chi(\mathcal X,\mathcal L^n)=n\deg\mathcal L+O(1)
\]
である．系\ref{cor:cubic-boundary-duality}と
系\ref{cor:cubic-saturated-hilbert}により
$H^1(\mathcal X,\mathcal L^{3m})=0$，
$h^0(\mathcal X,\mathcal L^{3m})=9m$ for $m>0$だから，
\begin{equation}\label{eq:v7-cubic-polarization-degree}
 \deg\mathcal L=3.
\end{equation}

定理\ref{thm:v7-faithful-equation-character}により，各特異点の方程式指標は忠実である．
整数次数の忠実な偏極$\mathcal L$に補題\ref{lem:v8-integral-degree-stack}を適用して，
全てのstabilizerが自明であることを得る．

\textbf{通常の曲線とsection環の識別．}
以上で全てのstabilizerが自明となる．従って$\mathcal X$はそのprojective coarse curve
$Y$と同型である．構造対象のcohomologyから$p_a(Y)=1$であり，
\eqref{eq:v7-cubic-polarization-degree}から$\deg\mathcal L=3$である．
補題\ref{lem:integral-genus-one-sections}により$\mathcal L$はvery ampleで，
section環は一次生成されるintegral plane cubicの座標環となる．
有限余核と第三Veroneseの一致は既に証明した三次境界の欠損公式による．
\end{proof}

\begin{lemma}[一次normal切断の可換な反標準境界]
\label{lem:v8-degree-one-commutative-boundary}
$U$をconnected locally finite graded右noetherian domainとし，$0\ne g\in U_1$をnormal元，
$R=U/(g)$を可換domainとする．$\qgr U$のhomological dimensionは2以下で，Serre functorは$(-1)[2]$とする．
また$D=\pi_RR$について$(D,(1))$がampleであり，正整数$d\geq3$に対して
\[
 R_n=\Hom_{\qgr R}(D,D(n)),\qquad\dim R_n=dn\quad(n>0),
 \qquad\End(D)=\C
\]
と仮定する．このとき，あるintegral projective genus-one curve $C$と次数$d$の線束$\mathcal M$について
\[
 (\qgr R,D,(1))\simeq(\coh C,\OO_C,\mathcal M\otimes-),
 \qquad R=\bigoplus_{n\geq0}H^0(C,\mathcal M^n)
\]
となる．特に$R$は次数1で生成される．
\end{lemma}
\begin{proof}
可換graded環$R$に対するcanonical weighted Proj stack
\[
 \mathcal X=[(\operatorname{Spec}R\setminus\{R_+\})/\mathbb G_m]
\]
を取る．$\dim R=2$で，$R_1\ne0$だからgeneric stabilizerは自明である．
また$\mathcal M=\OO_{\mathcal X}(1)$のfiber characterは全stabilizerで忠実であり，
元のgradingを保つ同値$(\qgr R,D,(1))\simeq(\coh\mathcal X,\OO,\mathcal M\otimes-)$がある．
この同値は定理\ref{thm:v7-commutative-boundary-geometry}の冒頭のequivariant sheafとgraded加群のSerre商による構成と同じである．

normal切断の構造対象の双対性は，ここでは次数$e=1$とSerre parameter $\ell=1$なので
\[
 \Ext^1(D,F)\simeq\Hom(F,D)^*,\qquad\Ext^{\geq2}(D,F)=0
\]
を与える．補題\ref{lem:v7-dualizing-trivial}から$\omega_{\mathcal X}\simeq\OO$である．
命題\ref{prop:v7-local-planarity}と定理\ref{thm:v7-faithful-equation-character}は任意の正整数$e$に対する周期性と
$(D,(e))$のAZ生成条件を用いるので，$e=1$として適用できる．
従って全atlasが局所平面であり，全特異点の方程式指標が忠実である．

正規化とconductor商によるRiemann--Rochは
$\chi(\mathcal X,\mathcal M^n)=n\deg\mathcal M+O(1)$を与える．
十分可除な$n$で$\mathcal M^n$はcoarse上のample線束へ降り，Serre消滅が成立する．
仮定のsection次元から$\deg\mathcal M=d\in\Z$を得る．
前補題により$\mathcal X$は通常のgenus-one curve $C$である．
仮定のsection同型によって$R$はその全section環に等しい．
Very amplenessは補題\ref{lem:integral-genus-one-sections}の長さ2の部分スキームの議論で，
$3-\operatorname{length}Z$を$d-\operatorname{length}Z$へ置き換えれば従う．
平方の積写像の全射性は$d$に関して帰納する．$d=3$は同補題による．
$d>3$ならsmoothな閉点$p$を取り，$\mathcal N=\mathcal M(-p)$とする．
帰納法により$H^0(\mathcal N)^2=H^0(\mathcal M^2(-2p))$である．
$\mathcal M$のvery amplenessから，$p$で消えないsection $u$と，$p$でちょうど一次に消えるsection $v$を選べる．
$u^2,uv$は$\mathcal M^2|_{2p}$のsection空間を張るので，
$\operatorname{Sym}^2H^0(\mathcal M)\to H^0(\mathcal M^2)$は全射となる．
$n\geq2$では$\mathcal M$のbasepoint-free pencilの完全列へ$\mathcal M^n$をtensorし，
$H^1(\mathcal M^{n-1})=0$を使えば
$H^0(\mathcal M)\otimes H^0(\mathcal M^n)\to H^0(\mathcal M^{n+1})$の全射性を得る．
従って全section環は次数1で生成される．
\end{proof}

\subsection{Normal twistのscalar性と中心化}

\begin{lemma}[幾何学的shiftを持つ境界上のnormal自己同型]
\label{lem:v8-thcr-normal-scalar}\label{lem:v7-normal-scalar-boundary}
$U$をconnected locally finite graded右noetherian代数，
$g\in U_1$をregular normal元とし，$ga=\nu(a)g$，$R=U/(g)$と置く．
$\qgr U$のhomological dimensionが$2$以下で，Serre functorが$(-1)[2]$であるとする．
さらに，構造対象とshiftを保つ同値
\[
 (\qgr R,D,(1))\simeq
 (\coh C,\OO_C,\mathcal M\otimes\rho^*(-))
\]
があり，$C$はintegral projective Gorenstein curve，
$\omega_C\simeq\OO_C$，$\mathcal M$は正次数の線束，
$\rho\in\operatorname{Aut}C$であるとする．
$R$は対応する全次数のTHCRであり，次数$1$で生成されると仮定する．
このときある$c\in\C^*$に対し
\[
 \overline\nu|_{R_n}=c^n\operatorname{id}_{R_n}\qquad(n\geq0)
\]
となる．$R$の可換性，$C$の滑らかさ，$\rho$の位数は仮定しない．
\end{lemma}

\begin{proof}
$D=\pi_RR$，$T=(1)$，$\Phi=(-)_{\overline\nu}$と置き，
$i_*:\qgr R\to\qgr U$を忘却functorとする．
$D(n)=T^nD$は$C$上の線束であり，
$\Phi D(n)\simeq D(n)$である．後者は，shiftした正則加群に対する
\[
 R_{\overline\nu}\longrightarrow R,\qquad a\longmapsto\overline\nu^{-1}(a)
\]
から従う．従って$F=D(n)$，$G\in\qgr R$に対して
\[
 \Ext^j_{\qgr R}(F,G)=0\quad(j\geq2).
\]
Normal切断の自然なtriangle
\[
 \Phi T^{-1}F[1]\longrightarrow Li^*i_*F
 \longrightarrow F\longrightarrow\Phi T^{-1}F[2]
\]
に$\Hom(-,G[2])$を適用すると
\[
 \Ext^2_{\qgr U}(i_*F,i_*G)
 \simeq \Ext^1_{\qgr R}(\Phi T^{-1}F,G)
\]
を得る．この同型は$F$の間の全ての射，および$G$に関して自然である．
左辺にambientのSerre双対性を，右辺に線束についてのGorenstein曲線の双対性を適用すると
\[
 \Hom_{\qgr R}(G,T^{-1}F)^*
 \simeq\Hom_{\qgr R}(G,\Phi T^{-1}F)^*
\]
となる．ここでは$\omega_C\simeq\OO_C$の自明化を一つ固定した．
特異な$C$の有界導来圏全体にSerre functor $[1]$を仮定しているのではない．

Yonedaの補題により，全ての整数$k$について同型
\[
 \eta_k:\Phi D(k)\xrightarrow{\sim}D(k)
\]
を得る．これらは全ての射$D(k)\to D(l)$について自然である．
上記の正則加群の同型で$\Phi D(k)$を$D(k)$と同定すると，
$\Phi$は$\Hom(D,D(n))=R_n$に$\overline\nu^{-1}$として作用する．
一方$\End(D(k))=H^0(C,\OO_C)=\C$なので，各$\eta_k$は非零scalarである．
従って$\eta$の自然性から$\overline\nu$は各$R_n$にscalarとして作用する．
$R$の一次生成性と$\overline\nu$の乗法性から，そのscalarは$c^n$となる．
\end{proof}

\begin{corollary}[THCRを保つ中心化]
\label{cor:v8-thcr-centralization}\label{cor:v7-centralization}
前補題の下で，右Zhang twist $V=U^{\nu^{-1}}$を取ると，
$g$は$V$のcentral regular元であり，$V/(g)\simeq R$がgraded代数として成り立つ．
従って中心化によって境界の自己同型$\rho$は変わらない．
この中心化のために$\nu$の三乗根や有限位数性を仮定する必要はない．
\end{corollary}

\begin{proof}
$g^2=\nu(g)g$とregular性から$\nu(g)=g$である．
$\tau_0=\nu^{-1}$と書き，$a\in U_m$に対して
\[
 g*a=g\nu^{-1}(a)=ag,
 \qquad a*g=a\nu^{-m}(g)=ag
\]
を得る．よって$g$は中心的である．商の積は
\[
 r_m*s_n=c^{-mn}r_ms_n.
\]
次数$n$の元を$c^{n(n-1)/2}$倍する写像が，このtwistされた商から$R$へのgraded代数同型を与える．
この計算には元の積$r_ms_n$の可換性を使っていない．
\end{proof}

\subsection{可換な三次境界からの完全な認識}\label{subsec:v7-full-recognition}

以下では，境界環の可換性を仮定するが，飽和性，標準Hilbert列，生成次数は仮定しない．
有限section欠損を直接消す代わりに，欠損のない反標準Veroneseを先に認識する．

\begin{lemma}[種数1のsection環の負次数双微分]\label{lem:v7-negative-biderivations}
$C$を$\C$上のintegral projective curveで$p_a(C)=1$，
$\mathcal L$を次数$d\geq3$の線束とし，
$R=\bigoplus_{n\geq0}H^0(C,\mathcal L^n)$と置く．
任意の負整数$j$について，次数$j$のskew biderivation
$\beta:R\times R\to R$は零である．
同じ結論は，$R$の有限余次元graded部分環にも成立する．
\end{lemma}
\begin{proof}
$C$は滑らかな種数1曲線，rational nodal curve，またはrational cuspidal curveである．
線束$\mathcal L$はampleであり，その全section環のpunctured cone $U=\operatorname{Spec}R\setminus\{R_+\}$は
$\mathcal L^{-1}$の零切断を除いた$\mathbb G_m$-bundleである．
この同一視には一次生成を仮定する必要はない．十分大きい隣接する二次数の線束が
大域生成されることを使い，それぞれの非零sectionで局所化すると，
次数$1$の可逆な局所生成元が得られる．
有限余次元部分環も，十分高い次数の共通idealで局所化すると同じ$U$を持つ．
従ってbiderivationを$U$へ局所化すればよい．

$\mathcal L$を自明化した開集合では$U$の環は$\OO_C[t,t^{-1}]$である．
$K(C)$の超越次数は$1$なので，その上のskew biderivationは零である．
従って$\beta$の$\OO_C\times\OO_C$への制限は零であり，次数$j$の条件から
\[
 \beta(t,f)=t^{j+1}\partial(f)\qquad(f\in\OO_C)
\]
である．ここで$\partial$は$\OO_C$のderivationである．
$t'=ut$という自明化の変更では$\partial'=u^{-j}\partial$となるので，
$\partial$は$H^0(C,\mathcal T_C\otimes\mathcal L^j)$の元を定める．
逆にこの元が零なら，局所環を生成する$\OO_C$と$t,t^{-1}$上で$\beta=0$となる．

滑らかな場合は$\mathcal T_C\simeq\OO_C$なので，$j<0$からこのsection空間は零である．
特異な場合は正規化$\nu:\mathbb P^1\to C$を使う．
局所derivationは正規化上のregularなvector fieldへ延長される．
nodeでは，二つのbranchの定数項が一致する条件から，vector fieldは両方の原像で消える．
cusp $\C[[t^2,t^3]]$では，$\partial(t)=a(t)\in\C((t))$と置くと
$2ta(t),3t^2a(t)\in\C[[t^2,t^3]]$が必要であり，$a(t)\in t\C[[t]]$となる．
従ってcuspの原像でもregularで消える．この直接計算から，少なくとも単射
\[
 \mathcal T_C\hookrightarrow\nu_*\mathcal T_{\mathbb P^1}
\]
を得る．ゆえに
\[
 H^0(C,\mathcal T_C\otimes\mathcal L^j)
 \hookrightarrow H^0(\mathbb P^1,\OO(2+jd))=0
 \qquad(j<0, d\geq3).
\]
$U$への局所化で零となるbiderivationは，環がdomainなので元の環上でも零である．
\end{proof}

\begin{proposition}[中心的な被覆の可換性]\label{prop:v7-central-commutativity}
$R$を前補題のsection環またはその有限余次元graded部分環とする．
$V$をconnected nonnegatively graded代数，$z\in V_e$を正次数のcentral regular元とし，
$V/(z)\simeq R$と仮定する．このとき$V$は可換である．
\end{proposition}
\begin{proof}
商の可換性から$[V,V]\subset zV$である．帰納的に$[V,V]\subset z^mV$と仮定する．
$R$上で
\[
 \beta_m(\overline a,\overline b)
 =z^{-m}[a,b]\pmod{zV}
\]
と置く．$z$のregular性により右辺は一意であり，liftを$zc$だけ変えると差は
$z[c,b]\in z^{m+1}V$に入るから，well-definedである．
commutatorのLeibniz則から，これは次数$-me$のskew biderivationである．
前補題により$\beta_m=0$，従って$[V,V]\subset z^{m+1}V$となる．
全ての$m$に帰納して$[V,V]\subset\bigcap_{m\geq1}z^mV=0$を得る．
最後の等式は非負次数と$e>0$から従う．
\end{proof}

\begin{remark}[既知の楕円変形との関係]
滑らかな楕円曲線については，Rains\cite[Corollaries 2.5--2.6]{RainsFiltered}が
負次数multiderivationの消滅と，section環のfiltered deformationの可換性を証明している．
同論文Proposition 2.9は，さらにベクトル束の斉次自己準同型環の変形の中心を扱う．
上の補題は必要な負次数biderivationの消滅を直接示し，nodeとcuspの場合も含めている．
normal被覆を中心的なものへ移す議論，および内在的六公理への適用は別に証明する．
\end{remark}

\begin{lemma}[射影平面上のexactな三乗根の軌道]
\label{lem:v7-plane-root-orbit}
$s$を$\coh\mathbb P^2$の$\C$線形exact autoequivalenceとし，
$s^3\OO\simeq\OO(3)$と仮定する．このとき全$n\in\Z$について
\[
 s\OO(n)\simeq\OO(n+1),\qquad s^n\OO\simeq\OO(n).
\]
さらに軌道環$\bigoplus_{n\geq0}\Hom(\OO,s^n\OO)$は次数1で生成され，
そのHilbert関数は$P(n)=\binom{n+2}{2}$である．
この段階では軌道環の可換性は主張しない．
\end{lemma}
\begin{proof}
$\coh\mathbb P^2$の単純対象は閉点のskyscraper sheafであり，exact autoequivalenceはこれらを置換する．
線束$L$は全ての閉点$p$に対して
\[
 \Ext^i(L,k(p))=0\quad(i>0),\qquad\dim\Hom(L,k(p))=1
\]
を満たす．逆に，coherent sheaf $F$がこの条件を満たせば線束である．
実際，点支持のlocal-to-global Extによりこれは各局所環上の条件となる．
極小自由分解に$\Hom(-,k(p))$を施すと微分は零なので，
$\Ext^1(F_p,k(p))=0$から$F_p$は自由であり，$\dim\Hom(F_p,k(p))=1$からrankは1である．
従って$s$は線束を線束へ移す．

$\operatorname{Pic}(\mathbb P^2)=\Z[\OO(1)]$より，$s\OO(n)=\OO(d_n)$と書ける．
各$n$について
\[
 3=\dim\Hom(\OO(n),\OO(n+1))
 =h^0\bigl(\mathbb P^2,\OO(d_{n+1}-d_n)\bigr)
\]
なので，$d_{n+1}-d_n=1$である．従って，ある整数$c$について$d_n=n+c$となる．
これを3回適用すると$s^3\OO=\OO(3c)$であり，仮定から$c=1$となる．
これで二つの軌道の式とHilbert関数を得る．

各$L_i=s^i\OO$と$\OO(i)$の同型を選ぶ．shiftによる同型
$\Hom(L_0,L_1)\simeq\Hom(L_i,L_{i+1})$により，軌道環の$n$重積の像は
\[
 \Hom(L_0,L_1)\otimes\cdots\otimes\Hom(L_{n-1},L_n)
 \longrightarrow\Hom(L_0,L_n)
\]
という全ての連続した射の合成の像に等しい．選んだ線束の同型の下で，これは
\[
 H^0(\OO(1))^{\otimes n}\longrightarrow H^0(\OO(n))
\]
という通常のsectionの積であり，全射である．従って軌道環は次数1で生成される．
\end{proof}

\begin{lemma}[射影平面上の三乗根と線形twist]
\label{lem:v8-plane-root-twist}
$s$を$\coh\mathbb P^2$の$\C$線形exact autoequivalenceとし，関手の自然同型
$s^3\simeq(-)\otimes\OO(3)$を仮定する．このとき，$\alpha\in\operatorname{PGL}_3(\C)$，$\alpha^3=1$が存在して
\[
 s\simeq\alpha^*(-)\otimes\OO(1)
\]
となる．軌道環$A_s$は三変数多項式環$S$の右Zhang twist $S^\sigma$であり，$\sigma$は$S_1$上の線形自己同型から誘導され，$\sigma^3=1$となるように取れる．
\end{lemma}
\begin{proof}
補題\ref{lem:v7-plane-root-orbit}により$s\OO(n)\simeq\OO(n+1)$である．
これらの同型を選び，$W=H^0(\mathbb P^2,\OO(1))$と置く．
$s$が$\Hom(\OO(n),\OO(n+1))$に誘導する線形同型を$M_n:W\to W$と書く．
通常のsectionの積の可換性と$s$のfunctorialityから
\[
 M_{n+1}(v)M_n(u)=M_{n+1}(u)M_n(v)\qquad(u,v\in W)
\]
を得る．$L=M_{n+1}M_n^{-1}$とすると$L(v)u=L(u)v$である．
互いに独立な一次式$u,v$は互いに素なので，$u\mid L(u)$となる．
従って$L(u)$は$u$のscalar倍であり，線形性から$L$はscalar自己同型である．
よって全$M_n$のprojective classは一定となる．
同型$s\OO(n)\simeq\OO(n+1)$をscalar倍して調整すれば，全$n$について$M_n=M$とできる．
高次のHomは一次のHomの積で生成されるので，$s$の作用はすべての線束間のHomで$\operatorname{Sym}^jM$となる．

$\alpha^*$の$W$上の作用が$M$であるように$\alpha\in\operatorname{PGL}_3$を取る．
すると$s$と$\alpha^*(-)\otimes\OO(1)$は，すべての線束とそれらの間の射からなる部分圏上で自然に同型である．
任意のcoherent sheafは線束の有限直和によるpresentationを持つので，この同型はcokernelにより全coherent sheafへ延長される．
presentationの比較には，任意の二つの全射への共通refinementを取り，十分負の線束からの射の持上げにSerre消滅を用いる．
これにより延長の独立性と射に関する自然性が従う．

ここでは対象の同型$s^3\OO\simeq\OO(3)$だけでなく，関手の自然同型を使う．
自然同型$s^3\simeq(-)\otimes\OO(3)$を一次のHomへ適用すると$M^3$はscalarとなり，従って$\alpha^3=1$である．
$M$をscalar倍して$M^3=1$としてよい．選んだ線束の同型の下で軌道環の積は
\[
 a\cdot b=M^n(a)b\qquad(a\in S_m,\ b\in S_n)
\]
である．各次数$n$で$a\mapsto M^{-n}(a)$とすれば，これは右Zhang積
\[
 a*b=a\,\sigma^m(b),\qquad\sigma=M^{-1},\quad\sigma^3=1
\]
へ移る．従って$A_s\simeq S^\sigma$である．
\end{proof}

\begin{theorem}[可換な第三Veronese境界からの認識]
\label{thm:v8-commutative-veronese-boundary}
六公理とexactな三乗根$s$を仮定し，$0\ne g\in A_3$をnormal元，$B=A/(g)$をdomainとする．
次の条件は同値である．
\begin{enumerate}[label=(\roman*)]
\item $B^{(3)}$は可換である．
\item $A^{(3)}$は可換である．
\item ある$\alpha\in\operatorname{PGL}_3(\C)$，$\alpha^3=1$について
\[
 (\Cat,\OO,s)\simeq
 (\coh\mathbb P^2,\OO,\alpha^*(-)\otimes\OO(1))
\]
である．
\item $A$は$\C[x_0,x_1,x_2]$の線形自己同型$\sigma$，$\sigma^3=1$による右Zhang twistである．
\end{enumerate}
これらが成立すると$E=0$，$B=\widehat B$となり，$A$は三次元二次Koszul AS-regularである．
$B$は$\alpha$-不変なintegral plane cubicの次数3のTHCRである．
境界の一次生成性，飽和性，幾何学的表示，$g$の中心性を追加仮定しない．
\end{theorem}
\begin{proof}
(i)を仮定する．$U=A^{(3)}$は元の反標準ample pairの軌道環であるから，
$(\Cat,\OO,\tau)\simeq(\qgr U,\pi U,(1))$であり，$\qgr U$のSerre functorは$(-1)[2]$である．
$g\in U_1$はregular normalである．
三次境界の飽和に関する系\ref{cor:cubic-saturated-hilbert}とample pairの降下から
\[
 R=U/(g)=B^{(3)}=\widehat B^{(3)},\qquad
 R_n=\Hom_{\qgr B}(D,D(3n)),\qquad\dim R_n=9n\quad(n>0)
\]
を得る．境界のample pair $(D,(3))$にAZ再構成を適用すれば
$(\qgr B,D,(3))\simeq(\qgr R,\pi R,(1))$となる．
従って補題\ref{lem:v8-degree-one-commutative-boundary}を$d=9$として適用でき，
$R$は通常のintegral genus-one curveの次数9の全section環である．

$R$の一次生成から，次数に関する帰納法で$U$も一次生成となる．
補題\ref{lem:v7-normal-scalar-boundary}と系\ref{cor:v7-centralization}により，
$U$のZhang twist $V$において$g$を中心化でき，$V/(g)\simeq R$を保つ．
命題\ref{prop:v7-central-commutativity}により$V$は可換である．

$V$はconnected noetherian graded domainで，次数1で生成される．
従って$Y=\operatorname{Proj}V$はintegral projective surfaceであり，
$\qgr V\simeq\coh Y$となる．曲面であることは
$\dim V_m=\dim U_m=\binom{3m+2}{2}$の二次成長から従う．

ここでtwist同値が構造対象のshiftをどのように移すかを明示する．
$\tau_0=\nu^{-1}$と書き，graded right module $M$を同じgraded vector spaceに
\[
 m*a=m\tau_0^{\deg m}(a)
\]
という作用を入れた$M^{\tau_0}$へ送る．このexact同値は有限次元加群を保つので$qgr$に降りる．
整数$j$について
\[
 (U(j))^{\tau_0}\xrightarrow{\sim}V(j),\qquad x\longmapsto\tau_0^j(x)
\]
はgraded module同型である．実際，$x\in U(j)_l=U_{l+j}$なら，
写像の左側の積を移したものは
$\tau_0^j(x)\tau_0^{l+j}(a)$となり，右側の積と一致する．
従って合成同値$\Cat\simeq\coh Y$は$\OO$を$\OO_Y$へ，$\tau^j\OO$を$\OO_Y(j)$へ移す．
一般対象上で元のshiftと$V$の標準shiftが一致すると仮定する必要はない．

$Y$は滑らかである．任意の閉点$y$について，点支持のlocal-to-global Extと
$\Cat$のhomological dimensionから
\[
 \Ext^i_{\OO_{Y,y}}(k(y),k(y))
 =\Ext^i_Y(k(y),k(y))=0\qquad(i>2)
\]
となる．極小自由分解と正則局所環の判定
\cite[Tag 065U, Proposition 10.110.5]{Stacks}により，全ての閉点で局所環は正則である．
$\C$上の有限型なので$Y$は滑らかである．
Serre関手の一意性を構造対象へ適用すると
\[
 \omega_Y[2]=\SSer_Y\OO_Y\simeq\OO_Y(-1)[2],
 \qquad\omega_Y^{-1}\simeq\OO_Y(1).
\]
$\OO_Y(1)$はvery ampleだから$Y$はdel Pezzo曲面である．
Serreの定理とHilbert関数から，$m\gg0$で
\[
 h^0(Y,\OO_Y(m))=\dim V_m=\tfrac92m^2+\tfrac92m+1.
\]
Riemann--Rochの最高次係数を比較すると$(-K_Y)^2=9$である．
滑らかな次数9のdel Pezzo曲面の古典的分類\cite[\S2.1, \S3.1.1]{DolgachevMartin}により$Y\simeq\mathbb P^2$となり，
$\OO_Y(1)\simeq\OO_{\mathbb P^2}(3)$である．

Serre関手の自然な一意性は，移された$s$について関手の自然同型
$s^3\simeq(-)\otimes\OO(3)$を与える．
補題\ref{lem:v8-plane-root-twist}から(iii)と(iv)を得る．
(iv)から$A^{(3)}=S^{(3)}$なので(ii)となり，(ii)から商を取って(i)を得る．
(iii)から(iv)は同じ補題による．これで同値を示した．

$H^1(\mathbb P^2,\OO(n))=0$ for all $n$なので$E=0$であり，
命題\ref{prop:cubic-defect-mass}より$B=\widehat B$である．
また補題\ref{lem:v7-plane-root-orbit}から$A$は一次生成なので，商$B$も一次生成である．
定理\ref{thm:cubic-generated-boundary-recognition}により$A$は三次元二次Koszul AS-regularである．

境界の表示を直接確認する．(iv)の同型で$g$が対応する三次式を$f\in S_3$とする．
$\sigma^3=1$なので$f*x=fx$であり，normal性から$fS_1=S_1\sigma(f)$となる．
$S$の一意分解性により$\sigma(f)\mid fx$がすべての$x\in S_1$について成立するため，
$\sigma(f)$は$f$のscalar倍である．従って$(f)$は$\sigma$-不変であり，
\[
 B\simeq(S/(f))^\sigma.
\]
$f$が非自明に斉次因数分解すれば，twistした商にも非零斉次零因子が生じるので，
$B$のdomain性から$f$は既約である．従って$C=V(f)$はintegral plane cubicである．
この商のtwist表示は$C$上の次数3のTHCR表示である．
\end{proof}

\begin{theorem}[可換な整三次normal境界からの完全な認識]
\label{thm:v7-commutative-cubic-recognition}
六公理とexactな三乗根$s$を仮定する．$0\ne g\in A_3$がnormalであり，
$B=A/(g)$が可換domainならば
\[
 A\simeq\C[x_0,x_1,x_2],\qquad\deg x_i=1.
\]
従って構造対象と固定したshiftを含めて
\[
 (\Cat,\OO,s)\simeq
 (\coh\mathbb P^2,\OO_{\mathbb P^2},(-)\otimes\OO(1)).
\]
また$B$は通常のintegral plane cubicの座標環であり，
$B=\widehat B$，$E_n=0$ for all $n\in\Z$となる．
境界の生成次数，標準Hilbert列，飽和性，曲線の滑らかさ，$g$の中心性は追加仮定ではない．
\end{theorem}
\begin{proof}
$B^{(3)}$も可換なので，定理\ref{thm:v8-commutative-veronese-boundary}から
$A$は次数1で生成され，$\dim A_n=P(n)$となる．
$g$の次数は3なので$A_2\to B_2$は単射である．$B$の可換性から$[A_1,A_1]=0$であり，
$A$は三変数多項式環のgraded商となる．Hilbert関数の一致により，この全射は同型である．
AZ再構成が構造対象と固定したshiftも保つことから，tripleの主張を得る．
境界の表示，$E=0$，$B=\widehat B$も前定理による．
\end{proof}

\begin{example}[第三Veroneseだけが可換な整三次境界]
\label{ex:v8-order-three-fermat}
$\zeta$を原始三乗根とし，$S=\C[x,y,z]$，
$\sigma(x)=x$，$\sigma(y)=\zeta y$，$\sigma(z)=\zeta^2z$と置く．
$A=S^\sigma$，$g=x^3+y^3+z^3$とすれば，$g$はcentralであり，
\[
 A^{(3)}=S^{(3)},\qquad
 B=A/(g)\simeq(S/(g))^\sigma
\]
となる．$B$はdomainで$B^{(3)}$は可換だが，
$x*y=\zeta xy$，$y*x=xy$なので$B$自体は可換でない．
$\qgr A\simeq\coh\mathbb P^2$で，移されたshiftの三乗は$\OO(3)$によるtensorとなるので，六公理とexactな三乗根を満たす．
従って定理\ref{thm:v8-commutative-veronese-boundary}の仮定は，境界自体の可換性より真に弱い．
この例の環と曲線の記述は既知のskew Fermat cubicの範囲に属し，新しい例とは扱わない．
\end{example}

\begin{remark}[既知のtwistと本稿の認識の区別]
線形自己同型によるZhang twist，そのgraded module圏の同値，AS-regular性の保存は既知の一般理論である\cite{ZhangTwist}．
本稿の三乗根の識別と必要な同値は上で直接証明した．
Baumann--Belmans--van Garderen\cite[\S6.1, Proposition 50, Lemma 52]{BBvGCentral}は，
既にskew polynomialであるambientについて，三乗根係数の特定の条件の下で
第三Veroneseが可換であり，skew cubicの$qgr$が通常の三次曲線の$\coh$と同値になることを示している．
本稿では六公理とnormal境界の第三Veroneseの可換性から，ambientの一次生成・多項式環のtwist表示・境界飽和性を導く．
この認識の組合せの文献上の新規性は未確認である．
\end{remark}

\subsection{幾何学的な三次境界と中心的な反標準被覆}
\label{subsec:v8-geometric-central-cover}
以下の還元では，境界環そのものの可換性を必要としない．
構造対象とshiftを保つ通常の種数1曲線表示が得られた場合に，
有限section欠損を残したまま，反標準環を中心的なTHCR被覆へ移す．

\begin{proposition}[幾何学的な三次境界からの中心的次数9被覆]
\label{prop:v8-geometric-central-reduction}
六公理とexactな三乗根$s$を仮定し，$0\ne g\in A_3$をnormal元，
$B=A/(g)$をdomainとする．さらに
\[
 (\qgr B,D,(1))\simeq
 (\coh C,\OO_C,\mathcal L\otimes\sigma^*(-))
\]
があり，$C$はintegral projective curveで$p_a(C)=1$，
$\deg\mathcal L=3$，$\sigma\in\operatorname{Aut}C$であると仮定する．
$B$の飽和性および一次生成性は仮定しない．
$\mathcal L_3=\mathcal L\otimes\sigma^*\mathcal L\otimes\sigma^{2*}\mathcal L$と置く．
このとき次が成立する．
\begin{enumerate}[label=(\roman*)]
\item $U=A^{(3)}$は次数$1$で生成され，$g\in U_1$であり，
\[
 U/(g)=B^{(3)}=\widehat B^{(3)}
 \simeq B(C,\mathcal L_3,\sigma^3).
\]
\item $U$には右Zhang twist $V$が存在し，
\[
 g\in Z(V)_1,\qquad V/(g)\simeq B(C,\mathcal L_3,\sigma^3),
 \qquad \dim V_m=\binom{3m+2}{2}.
\]
$V$は次数$1$で生成される左右noetherian domainであり，
$\qgr V$は$\Cat$にexact同値である．
\item 元の$A$も左右noetherianであり，左$U$加群として有限生成である．
\end{enumerate}
\end{proposition}

\begin{proof}
仮定したtriple同値から
\[
 \widehat B\simeq B(C,\mathcal L,\sigma)
\]
となる．補題\ref{lem:integral-genus-one-sections}を，線束
$\mathcal L_n$と$\sigma^{n*}\mathcal L$に適用すると，この全section環は次数$1$で生成される．
従ってその第三Veroneseも次数$1$で生成される．
また$C$はGorensteinで$\omega_C\simeq\OO_C$である．

三次境界の欠損公式から$B_{3m}=\widehat B_{3m}$ for $m\geq0$である．
従って上の(i)の商の同型を得る．$R=U/(g)$と置けば，
$R$の一次生成性から$U$も一次生成される．
実際，任意の$u\in U_m$の像を$R_1$の元の積で表し，そのliftとの差を$gU_{m-1}$に入れて帰納すればよい．

元の反標準ample pair $(\OO,s^3)$にAZ再構成を直接適用すると
\[
 (\Cat,\OO,s^3)\simeq(\qgr U,\pi U,(1))
\]
である．従って$\qgr U$のhomological dimensionは$2$で，Serre functorは$(-1)[2]$となる．
通常のTHCRの同値は
\[
 (\qgr R,\pi R,(1))\simeq
 (\coh C,\OO_C,\mathcal L_3\otimes\sigma^{3*}(-))
\]
を保つ．補題\ref{lem:v8-thcr-normal-scalar}と系\ref{cor:v8-thcr-centralization}を適用して$V$を得る．
Twistは次数ごとのベクトル空間，domain性，右noetherian性，および$qgr$のexact同値を保つ．
Hilbert関数は元の反標準条件そのものである．

両側noetherian性と有限性は，幾何学的表示を仮定しない
系\ref{cor:v8-cubic-ambient-noetherian}と補題\ref{lem:v8-veronese-finite}から従う．
Zhang twistは左noetherian性も保つので，$V$も両側noetherianである．
\end{proof}

\begin{remark}[中心化で消えない問題]
上の命題は，幾何学的な境界表示が得られた場合の構造還元である．
$B=\widehat B$や$A$の一次生成性をまだ結論していない．
滑らかな$C$と無限位数の平行移動$\sigma^3$の場合，$V$は
Rogalski--Sierra--Staffordの意味の次数$9$のelliptic algebraとなる．
同論文\cite[Notation 3.2, Theorem 4.3]{RSS}のSklyanin第三Veronese認識には，
特定の$g$-divisible idealとovermodule，それらの自己準同型環のHilbert関数，
および境界上の二つのdivisor等式(4.1)--(4.2)がさらに必要である．
次数$9$と有限$\operatorname{hdim}\qgr V$だけからの認識は，
同論文Remark 4.13でも未解決として区別される．
この引用から有限section欠損の消滅を無条件に推論してはならない．

\end{remark}

\begin{corollary}[通常の曲線表示の下で残る一つの障害]
\label{cor:v8-geometric-one-obstruction}
命題\ref{prop:v8-geometric-central-reduction}の仮定の下で，Euler orbit rankは有限で，
次は同値である．
\begin{enumerate}[label=(\roman*)]
\item $E_{-1}=0$．
\item $B=\widehat B$．
\item $B$は次数$1$で生成される．
\item 固定した$A$は三次元二次Koszul AS-regularである．
\item Euler orbit rankは$3$である．
\item Euler orbit rankは$4$以下である．
\item $\dim A_1=3$．
\item $\dim A_2=6$．
\end{enumerate}
これらが成立しなければ，Euler orbit rankは$5$であり，$\dim A_1\leq2$，$\dim A_2\leq5$となる．
特に，$D^b(\Cat)$のEuler形式がrank $4$以下の群を経由すれば認識が成立する．
\end{corollary}
\begin{proof}
定理\ref{thm:v8-cubic-standard-sections}と系\ref{cor:v8-cubic-automatic-euler-rank}により，
幾何学的表示を用いずEuler orbit rankは有限で，$\epsilon=\dim E_{-1}$が零なら$3$，正なら$5$である．
また三次境界の欠損公式により(i)と(ii)は同値である．
$\widehat B=B(C,\mathcal L,\sigma)$は次数$1$で生成されるので(ii)から(iii)を得る．
逆に(iii)なら，既に証明した一次生成三次境界の全次数識別から飽和性(ii)を得る．
(iii)と(iv)は三次normal被覆の認識定理で同値である．
(iv)から(vii), (viii)は標準Hilbert関数による．
逆に(vii)なら$B_1=A_1=\widehat B_1$なので，$\widehat B$の一次生成性から$B=\widehat B$となる．
(viii)なら$B_2=\widehat B_2$であり，$B_3=\widehat B_3$は一般の三次境界で成立する．
一次生成環$\widehat B$の次数$2,3$成分は全ての次数$n\geq2$を生成するので，
$B_n=\widehat B_n$ for $n\geq2$である．
従ってsection欠損は次数$1$だけに支持されるが，剰余類$1,2$の欠損総量が等しいことから，この欠損も零となる．
これで全ての同値を得る．一般の上界$h_1\leq3,h_2\leq6$と合わせて，残りの主張も従う．
\end{proof}

\begin{corollary}[中心化後の認識から元の三乗根の認識へ]
\label{cor:v8-recognized-cover-root}
命題\ref{prop:v8-geometric-central-reduction}の設定で，
中心化後の$V$について$\qgr V$が三次元二次AS-regular代数$S$の$qgr$にexact同値なら，
元の固定した$A$も三次元二次Koszul AS-regularである．
特に，$C$が滑らかで$\sigma^3$が無限位数の平行移動の場合に，
$V$がRogalski--Sierra--Stafford\cite[Notation 3.2, (4.1)--(4.2)]{RSS}の条件を満たせば，この結論が成立する．
\end{corollary}
\begin{proof}
$\Cat\simeq\qgr V\simeq\qgr S$なので，命題\ref{prop:v9-marking-independence}から直ちに結論を得る．
ここでは同値が元の$s$を標準shiftへ送ることを仮定する必要はない．

この境界設定では，独立した数値的な証明もある．
$\qgr S$の標準full exceptional collectionは長さ$3$なので，
$K_0(\qgr S)\simeq\Z^3$である．従って$\Cat\simeq\qgr V\simeq\qgr S$のEuler形式は
rank $3$の群を経由する．元の$s$の軌道のEuler rankも$3$以下なので，
系\ref{cor:v8-geometric-one-obstruction}を適用できる．
この議論は，圏同値が元の$s$を$S$の標準shiftへ送ることを必要としない．

最後の場合にはRSSのTheorem 4.3から$V\simeq S^{(3)}$を得る．
ここで$S$はSklyanin代数で，次数$1$で生成されるので，
通常のVeronese同値により$\qgr V\simeq\qgr S$となり，前半を適用できる．
\end{proof}

\subsection{中心的な滑らかな楕円境界による認識}

\begin{proposition}[既知の被覆認識の適用]\label{prop:elliptic}
さらに$0\ne g\in A_3$が$A$の中心元であり，次数付き環同型
\[
 A/gA\simeq B(C,\mathcal L,\sigma)
\]
が存在すると仮定する．ここで$C$は滑らかな楕円曲線，$\deg\mathcal L=3$，
$\sigma$は非torsion平行移動，$B$はtwisted homogeneous coordinate ringである．
このとき$A$は対応する三次元二次Sklyanin代数である．
\end{proposition}
\begin{proof}
$\operatorname{hdim}\qgr A=2$なので，定理\ref{thm:cubic-normal-cover-dichotomy}から
$A$は境界の二次包絡と同型である．この境界の二次包絡は対応するSklyanin代数である．
後者の既知の識別と三次中心元による商表示については
\cite[Proposition 4.5, Remark 4.6]{RSS}を参照する．
\end{proof}

Sklyaninという型の識別は，ここでは滑らかな楕円曲線と非torsion平行移動という仮定の下で述べた．
三次元二次AS-regularという認識自体は，定理\ref{thm:cubic-normal-cover-dichotomy}により
特異なintegral曲線や任意の位数の自己同型も扱う．
公理からnormal境界の存在を証明したわけではない．
また，$A^{(3)}$の中心にある元が$A$でも中心的であるとは限らず，
抽象的な境界圏の同値だけでは，全次数の商環の同型の代わりにならない．

\section{一般認識問題の残余}\label{sec:remaining}
元の問いは，定義\ref{def:axioms}とexactな三乗根の存在から，非可換射影平面としての表示が従うか，というものである．
命題\ref{prop:v9-marking-independence}により，次の二点は固定した三乗根の座標環だけでなく，
非標識のabel圏の認識にも直結する未解決箇所である．
\begin{enumerate}
\item $E_{-1}=\Ext^1(\OO,s^{-1}\OO)$は自動的に零か．
\item $E_{-1}=0$の下で，$s^3\OO$は$\thick(\OO,s\OO,s^2\OO)$に属するか．
\end{enumerate}
前者の十分条件は節\ref{sec:euler}で与えたが，その数値仮定を公理から導く問題は残る．
有限性の側では，$E_{\geq0}$の左第三Veronese有限生成と一様右消滅までを証明した．
残る問題は左torsion性，または命題\ref{prop:basscriterion}の第二Bass空間の有限次元性である．
反標準方向だけのeventual消滅も，既に全対象Serre消滅と同値であることに注意を要する．
ただし$A$または$A^{(3)}$に正次数の非零normal元がある範囲では，
命題\ref{prop:normal-vanishing}がこの有限性を解決する．
従って$E$が無限次元となる真の例があれば，両環とも正次数normal元を持たない必要がある．

内在的な線束対象については，六公理だけで例外性とHom半順序を得る．
条件\eqref{eq:v10-SV}の下では数値類が対象を一意に決め，さらに全数値rank $3$なら
線束対象全体は$\{s^n\OO\}_{n\in\Z}$である．この追加条件の下でもfull性は未証明であり，
heartの正の対象の数値的検出を残余導来圏の消滅と取り違えない．

生成の側では，$E_{-1}=0$の下で$\PP$の回転軌道が残余圏全体を生成する．
全整数Euler二次性と$E$の有限次元性を確保した範囲なら，
命題\ref{prop:puritykillsdefect}により二重mutationの純粋性を調べればよい．
これは評価射の核と余核という実対象の問題であり，$K_0$上の零性では代用できない．
そのEuler二次性の下では，積写像は低次数を四候補まで制限するが，$(h_1,h_2,h_3)=(4,7,10)$の排除，
全環の一次生成性，二次関係性は一般には未証明である．

次数1・2の整なnormal切断，および一次生成された整な三次normal切断では，
節\ref{sec:normal-boundary}，\ref{sec:degree2-hilbert}，\ref{sec:elliptic}により認識は閉じている．
さらに三次domain境界の第三Veroneseが可換なら，生成次数や飽和性を仮定せず，
定理\ref{thm:v8-commutative-veronese-boundary}によって三乗が恒等な自己同型による多項式環のtwistと認識できる．
境界自体が可換なら$A=\C[x_0,x_1,x_2]$である．
従ってnormal境界による未解決の核心は，第三Veroneseが非可換で，かつ一次生成を仮定しない三次domain境界である．

この一般の境界でも$E_{3m}=0$は全$m\in\Z$で成立し，$\widehat B_n=\Hom_{\qgr B}(D,D(n))$に対して
\[
 0\longrightarrow B_n\longrightarrow\widehat B_n
 \longrightarrow E_{n-3}\longrightarrow E_n\longrightarrow0\qquad(n>0),
 \qquad b_n\leq\dim\widehat B_n\leq3n
\]
である．命題\ref{prop:cubic-defect-mass}は，$C=\widehat B/B$の両非零剰余類の欠損総量が
ともに$\epsilon=\dim E_{-1}$であることを与え，$E_{-1}=0$と$B=\widehat B$を同値にする．
定理\ref{thm:v8-cubic-standard-sections}は$\widehat b_n=3n$を一般に与え，Euler rankは自動的に有限である．
また$A,B$の両側noetherian性と高Veroneseの全次数THCR表示も確保される．
残る問題はrank $5$，$\epsilon>0$を排除すること，および元のgradingに対する一次生成を導くことである．

可換境界ではcanonical weighted Projが元の$qgr$，構造対象，shiftを同時に表す．
高次Ext周期性，方程式指標の忠実性，conductor次数式と整数次数の偏極によって全stabilizerを除き，
欠損のない反標準Veroneseを中心化して認識へ至る．
高Veroneseの曲線が滑らかなら，定理\ref{thm:v10-smooth-veronese-heart}により
元のheart，構造対象，次数shiftは通常の次数3の曲線tripleとして復元される．
特異曲線でも，関数体自己同型の曲線への延長と，一つの原始剰余層の可逆性があれば同じ結論となる．
従ってこの復元に残る一般の課題は，この二つの局所条件を導くか，またはそれらを不要にすることである．
cohomologyの保存だけでheart同値を推論しているわけではなく，剰余層の積と次数一致が追加の実質である．

通常の種数1曲線表示が得られれば，系\ref{cor:v8-geometric-one-obstruction}により
$\epsilon=0$一つで認識が閉じ，$h_1=3$または$h_2=6$も同値な判定となる．
命題\ref{prop:v10-unimodular-cubic-defect}は，rank $5$以下の全Euler格子の判別式が
周期成分の平方類$[3]$と両立しない場合に$\epsilon=0$を与える．
特に滑らかな高Veronese境界と長さ$5$以下のfull exceptional collectionがあれば認識が閉じる．
全数値rankやunimodularity自体の公理からの導出は未解決である．
中心化した反標準環がAS-regular曲面の圏として認識できれば，その$K_0$のrankから$\epsilon=0$を得る．
ただしRSSの認識に必要な$g$-divisible ideal，overmodule，divisor等式は一般にはまだ構成していない．
次数9と有限homological dimensionだけで既知の認識定理を適用してはならない．

有限余次元の境界部分環については，abstractな$qgr$の同値だけから元のgradingの飽和性や一次生成性を結論してはならない．
研究ノートの可換な有限欠損ambient例はこの論理的障害を示すが，
六公理とexactな三乗根を満たすnormal被覆ではない．それらの全仮定を満たす可換境界の持上げは，
上の完全認識定理で排除される．

別の範囲として，一次生成で両側noetherianかつ双有理的に幾何的な場合は節\ref{sec:birational}で認識した．
normal元の存在を使わない一般への拡張では，障害加群の有限性または評価射の純粋性を直接示す課題が残る．

研究ノートには，AZ生成条件を除くとEuler式と全Ext消滅を満たしても$\PP\ne0$となる
重み付き射影stack，左右有限性の形式的な誤推論を防ぐYoneda型代数，
全反標準Hilbert列を持つ非normalなGorenstein domain，
任意の有限個の反標準Hilbert検査を通るnormalな環の族を記載した．
全ての例について欠ける仮定を明記しており，全基礎公理を満たす反例ではない．
二つの認識障害が実際に消えれば，定理\ref{thm:finite}により
strong性，geometricity，一次生成性，二次関係性を別々に仮定する必要はない．

\begingroup\small
\begin{thebibliography}{99}
\bibitem{CKSTHCR} A.~Chirvasitu, R.~Kanda and S.~P.~Smith,
\emph{Maps from Feigin and Odesskii's elliptic algebras to twisted homogeneous coordinate rings},
arXiv:1908.06525v3 (2021), \url{https://arxiv.org/abs/1908.06525}.
\S2.3，Proposition 2.2とTheorem 2.4で積・左右規約・THCRのSerre同値の再掲を確認．
同定理が引用するArtin--Van den Bergh原論文の全証明を直接確認したという意味ではない．
\bibitem{AOUCompact} T.~Abdelgadir, S.~Okawa and K.~Ueda,
\emph{Compact moduli of noncommutative projective planes},
arXiv:1411.7770v1 (2014),
\url{https://arxiv.org/abs/1411.7770}.
Theorems 7.1, 7.3，Appendix A，特にDefinition A.3，Corollary A.7，Theorem A.8を本文で確認．
本稿はこのpreprintの定理番号を用いる．出版情報は未確認であり推測しない．
\bibitem{KachiKollar} Y.~Kachi and J.~Koll\'ar,
\emph{Characterizations of $\mathbb P^n$ in arbitrary characteristic},
Asian J. Math. \textbf{4} (2000), 115--122,
\url{https://doi.org/10.4310/AJM.2000.v4.n1.a8},
\url{https://arxiv.org/abs/math/0003230}.
Corollary 2(1)の数値同値による最大指数の特徴付けを確認．

\bibitem{AZ} M. Artin and J. J. Zhang,
\emph{Noncommutative projective schemes}, Adv. Math. \textbf{109} (1994), 228--287.
\href{https://doi.org/10.1006/aima.1994.1087}{doi:10.1006/aima.1994.1087}.
\bibitem{BS} T. Bridgeland and D. Stern,
\emph{Helices on del Pezzo surfaces and tilting Calabi--Yau algebras}, Adv. Math. \textbf{224} (2010), 1672--1716.
\href{https://doi.org/10.1016/j.aim.2010.01.018}{doi:10.1016/j.aim.2010.01.018};
\href{https://arxiv.org/abs/0909.1732}{arXiv:0909.1732v1}.
\bibitem{MU} I. Mori and K. Ueyama,
\emph{A categorical characterization of quantum projective spaces}, J. Noncommut. Geom. \textbf{15} (2021), 489--529.
\href{https://doi.org/10.4171/JNCG/403}{doi:10.4171/JNCG/403};
\href{https://arxiv.org/abs/1708.00167}{arXiv:1708.00167v4}.
\bibitem{ATV1990} M.~Artin, J.~Tate and M.~Van den Bergh,
\emph{Some algebras associated to automorphisms of elliptic curves},
The Grothendieck Festschrift, Vol.~I, Progress in Mathematics \textbf{86},
Birkh\"auser Boston, 1990, 33--85.
\href{https://doi.org/10.1007/978-0-8176-4574-8_3}{doi:10.1007/978-0-8176-4574-8\_3}（2007年再版）．
使用箇所：Theorem 1，Theorem 6.6，\S7（特にProposition 7.22）．
\bibitem{IM} A. Itaba and M. Matsuno,
\emph{AS-regularity of geometric algebras of plane cubic curves}, J. Aust. Math. Soc. \textbf{112} (2022), 193--217.
\href{https://doi.org/10.1017/S1446788721000070}{doi:10.1017/S1446788721000070};
\href{https://arxiv.org/abs/1905.02502}{arXiv:1905.02502v3}.
\bibitem{RSS} D. Rogalski, S. J. Sierra and J. T. Stafford,
\emph{Ring-theoretic blowing down II: Birational transformations}, J. Noncommut. Geom. \textbf{17} (2023), 1045--1097.
\href{https://doi.org/10.4171/JNCG/510}{doi:10.4171/JNCG/510};
\href{https://ems.press/content/serial-article-files/31044}{publisher PDF}.
\bibitem{Keeler} D. S. Keeler,
\emph{The rings of noncommutative projective geometry},
\href{https://arxiv.org/abs/math/0205005}{arXiv:math/0205005v2}, 2002.
\bibitem{VdB} M. Van den Bergh,
\emph{Non-commutative quadrics},
\href{https://arxiv.org/abs/0807.3753}{arXiv:0807.3753v3}, 2010.
\bibitem{KS} A. Kuznetsov and M. Smirnov,
\emph{On residual categories for Grassmannians}, Proc. London Math. Soc. (3)
\textbf{120} (2020), 617--641.
\href{https://doi.org/10.1112/plms.12294}{doi:10.1112/plms.12294};
\href{https://arxiv.org/abs/1802.08097}{arXiv:1802.08097v3}.
\bibitem{YZ} A. Yekutieli and J. J. Zhang,
\emph{Serre duality for noncommutative projective schemes},
Proc. Amer. Math. Soc. \textbf{125} (1997), 697--707.
\href{https://doi.org/10.1090/S0002-9939-97-03782-9}{doi:10.1090/S0002-9939-97-03782-9};
\href{https://www.math.bgu.ac.il/~amyekut/publications/nc-serre/paper.pdf}{author PDF}.
\bibitem{Mizuno} Y. Mizuno,
\emph{Dualizing complexes over $\Z$-algebras},
\href{https://arxiv.org/abs/2509.13073}{arXiv:2509.13073v3}, revised 4 September 2026.
\bibitem{IS} N. Iyudu and S. Shkarin,
\emph{Classification of quadratic and cubic PBW algebras on three generators},
\href{https://arxiv.org/abs/1806.06844v1}{arXiv:1806.06844v1}, 18 June 2018.
使用箇所：クラス$\Omega$の定義，Theorem 1.11，\S12.1，p.~78．
\bibitem{RSsurface} D. Rogalski and J. T. Stafford,
\emph{A class of noncommutative projective surfaces},
Proc. London Math. Soc. (3) \textbf{99} (2009), no.~1, 100--144.
\href{https://doi.org/10.1112/plms/pdn054}{doi:10.1112/plms/pdn054};
\href{https://arxiv.org/abs/math/0612657}{arXiv:math/0612657v1}.
使用箇所：Theorem 1.1，Corollary 1.2(I)(b), (II)(f)．
\bibitem{Stacks} The Stacks Project Authors, \emph{The Stacks Project}.
\href{https://stacks.math.columbia.edu/tag/065U}{Tag 065U}（正則局所環と剰余体の有限射影次元），
\href{https://stacks.math.columbia.edu/tag/01PR}{Tag 01PR}（ample可逆層の同値条件），
\href{https://stacks.math.columbia.edu/tag/0B5A}{Tag 0B5A}（Cohen--Macaulay双対性），
\href{https://stacks.math.columbia.edu/tag/0BS2}{Tag 0BS2}・
\href{https://stacks.math.columbia.edu/tag/0BS3}{Tag 0BS3}（曲線の双対性），
\href{https://stacks.math.columbia.edu/tag/0E3A}{Tag 0E3A}・
\href{https://stacks.math.columbia.edu/tag/0H2V}{Tag 0H2V}（切断の消滅とvery ampleness）．
確認日：2026年11月16日．
\bibitem{Hartshorne} R. Hartshorne, \emph{Algebraic Geometry},
Graduate Texts in Mathematics \textbf{52}, Springer, 1977.
ISBN 0-387-90244-9．使用箇所：Chapter V，Theorems 1.9, 1.10（pp.~364--365）．
\bibitem{StanleyReciprocity} R. P. Stanley,
\emph{Combinatorial reciprocity theorems}, Adv. Math. \textbf{14} (1974), 194--253.
\href{https://doi.org/10.1016/0001-8708(74)90030-9}{doi:10.1016/0001-8708(74)90030-9};
\href{https://math.mit.edu/~rstan/pubs/pubfiles/22.pdf}{author PDF}.
使用箇所：\S6，Proposition 6.1，p.~208．
\bibitem{BeckDevelin} M. Beck and M. Develin,
\emph{On Stanley's reciprocity theorem for rational cones},
\href{https://arxiv.org/abs/math/0409562v3}{arXiv:math/0409562v3}, 4 August 2005.
使用箇所：\S4，Theorem 5，p.~8．

\bibitem{ASquadratic} M.~Artin and J.~T.~Stafford,
Noncommutative graded domains with quadratic growth,
Invent. Math. 122 (1995), 231--276.
\url{https://doi.org/10.1007/BF01231444}.
Theorems 0.1--0.2; statements also in \cite[Theorems 5.1--5.2]{SVdBsurvey}.
\bibitem{SVdBsurvey} J.~T.~Stafford and M.~Van den Bergh,
Noncommutative curves and noncommutative surfaces,
Bull. Amer. Math. Soc. (N.S.) 38 (2001), 171--216.
\url{https://doi.org/10.1090/S0273-0979-01-00894-1}.
\href{https://arxiv.org/abs/math/9910082}{arXiv:math/9910082v2}.
使用箇所：Theorem 3.2，Theorems 5.1--5.2，Corollary 5.3；
一次生成を外す際の適用範囲はTheorems 5.5--5.6と区別する．
\bibitem{CassidyShelton} T.~Cassidy and B.~Shelton,
Generalizing the notion of Koszul algebra,
Math. Z. 260 (2008), no.~1, 93--114.
\url{https://doi.org/10.1007/s00209-007-0263-8}.
arXiv:0704.3752v1.
\S6, equations (6.1)--(6.2), p.~20; \S9.


\bibitem{SheltonTingey} B.~Shelton and C.~Tingey,
On Koszul algebras and a new construction of Artin--Schelter regular algebras,
J. Algebra 241 (2001), no.~2, 789--798.
\url{https://doi.org/10.1006/jabr.2001.8781}.
\bibitem{HuMoriWu} H.~Hu, I.~Mori and W.~Wu,
Classification of noncommutative central conics,
arXiv:2602.03236v1 (3 February 2026).
\url{https://arxiv.org/abs/2602.03236}.
Definitions 1.1--1.2; \S4.

\bibitem{CCKM} T.~Cassidy, A.~Conner, E.~Kirkman and W.~F.~Moore,
\emph{Periodic free resolutions from twisted matrix factorizations},
J. Algebra \textbf{455} (2016), 137--163.
\href{https://doi.org/10.1016/j.jalgebra.2016.01.037}{doi:10.1016/j.jalgebra.2016.01.037};
\href{https://users.wfu.edu/moorewf/Papers/twistedMFs.pdf}{著者公開の出版PDF}.
確認箇所：Theorem 4.2，p.~150．同定理はambient代数の有限global dimensionを仮定する．
\bibitem{LPstack} Y.~Lekili and A.~Polishchuk,
\emph{Auslander orders over nodal stacky curves and partially wrapped Fukaya categories},
J. Topology \textbf{11} (2018), 615--644.
\href{https://doi.org/10.1112/topo.12064}{doi:10.1112/topo.12064};
\href{https://lekili.freemath.xyz/research/pwrap.pdf}{著者公開版}.
確認箇所：Introduction，著者公開版p.~3のcyclic nodeとbalanceの定義．

\bibitem{Gulliksen1967} T.~H.~Gulliksen,
\emph{A note on the homology of local rings}, Math. Scand. \textbf{21} (1967), 296--300.
\url{https://doi.org/10.7146/math.scand.a-10864}.
使用箇所：Proposition 6，p.~300．同論文のcodimensionはdepthの意味である．
\bibitem{Levchenko} D.~Levchenko,
\emph{Serre duality for tame Deligne--Mumford stacks}, Res. Math. Sci. \textbf{9} (2022), article 67.
\url{https://doi.org/10.1007/s40687-022-00367-7}.
使用箇所：Proposition 2とその証明．
\bibitem{NironiDuality} F.~Nironi,
\emph{Grothendieck duality for Deligne--Mumford stacks},
\href{https://arxiv.org/abs/0811.1955v2}{arXiv:0811.1955v2}, 9 September 2009.
使用箇所：Corollary 2.30，Proposition 2.31，Example 3.7，Theorem 3.8，pp.~26, 30--32．
\bibitem{SmythGenusOne} D.~I.~Smyth,
\emph{Modular compactifications of the space of pointed elliptic curves I},
Compos. Math. \textbf{147} (2011), 877--913.
\url{https://doi.org/10.1112/S0010437X10005014}.
使用箇所：Appendix A，Proposition A.3とその証明，pp.~910--913；conductorの判定はp.~910．
\bibitem{RainsFiltered} E.~M.~Rains,
\emph{Filtered deformations of elliptic algebras},
\href{https://arxiv.org/abs/2107.13540v2}{arXiv:2107.13540v2}, 17 February 2022.
使用箇所：Lemma 2.4，Corollaries 2.5--2.6，Proposition 2.9，pp.~7--9．
一般のelliptic algebra認識との境界はConjecture 1も参照．
\bibitem{DolgachevMartin} I.~Dolgachev and G.~Martin,
\emph{Automorphisms of del Pezzo surfaces in odd characteristic},
J. London Math. Soc. \textbf{109} (2024), no.~5, e12905.
\url{https://doi.org/10.1112/jlms.12905}.
使用箇所：\S2.1，\S3.1.1に再掲された次数9の古典的分類（代数閉体，標数$\ne2$）．

\bibitem{BvdBGenerators} A.~Bondal and M.~Van den Bergh,
\emph{Generators and representability of functors in commutative and noncommutative geometry},
Moscow Math. J. \textbf{3} (2003), no.~1, 1--36.
\href{https://arxiv.org/abs/math/0204218v2}{arXiv:math/0204218v2}.
使用箇所：\S4.1，式(4.1)，Lemmas 4.1.4--4.1.6，著者版pp.~16--17．
このcohomology公式に後続の有限cohomological dimension仮定は不要．
\bibitem{ZhangTwist} J.~J.~Zhang,
\emph{Twisted graded algebras and equivalences of graded categories},
Proc. London Math. Soc. (3) \textbf{72} (1996), 281--311.
\url{https://doi.org/10.1112/plms/s3-72.2.281}.
一般理論との比較のための引用．必要なtwist同値と三乗根の識別は本文で直接証明する．
\bibitem{BBvGCentral} T.~Baumann, P.~Belmans and O.~van Garderen,
\emph{Central curves on noncommutative surfaces},
Trans. Amer. Math. Soc. \textbf{379} (2026), 3537--3589.
\url{https://doi.org/10.1090/tran/9571};
\href{https://arxiv.org/abs/2410.07620v1}{arXiv:2410.07620v1}.
参照番号・頁は確認済みのarXiv版による：\S6.1，Proposition 50，Lemma 52，式(135)，pp.~34--35．
出版情報は\href{https://pbelmans.ncag.info/writing/}{著者の公開文献一覧}で確認．出版社最終版本文は未取得．
\end{thebibliography}
\endgroup
\end{document}
